\documentclass[12pt]{article}

\usepackage[english]{babel}
\usepackage{iftex}
\ifPDFTeX
    \usepackage{times}
    \usepackage[T1]{fontenc}
\else
    \usepackage{fontspec}
    \setmainfont{TeX Gyre Termes}
    \ifXeTeX
        \usepackage{xeCJK}
        \setCJKmainfont{Noto Serif CJK SC}
        \setCJKsansfont{Noto Sans CJK SC}
        \setCJKmonofont{Noto Sans CJK SC}
    \fi
\fi
\usepackage[a4paper,margin=2.4cm]{geometry}
\usepackage{amsmath}
\usepackage{amssymb}
\usepackage{booktabs}
\usepackage{adjustbox}
\usepackage{array}
\usepackage{enumitem}
\usepackage{csquotes}
\usepackage{capt-of}
\usepackage{xr-hyper}
\usepackage{xurl}
\usepackage[numbers,sort&compress]{natbib}
\usepackage[colorlinks=true, allcolors=blue, breaklinks=true]{hyperref}
\externaldocument[ms-][nocite]{main}
\urlstyle{same}
\usepackage{xcolor}
\usepackage[normalem]{ulem}
\usepackage[most]{tcolorbox}

\emergencystretch=8em
\tolerance=2500
\hbadness=4000
\hfuzz=1pt
\setlength{\parindent}{0pt}
\setlength{\parskip}{0.45em}

\newcommand{\deleted}[1]{\textcolor{red}{\sout{#1}}}
\newcommand{\added}[1]{\textcolor{blue}{#1}}
\newcolumntype{L}[1]{>{\raggedright\arraybackslash}p{#1}}

\setcounter{tocdepth}{2}

\newcommand{\tocsection}[1]{%
    \phantomsection
    \section*{#1}%
    \addcontentsline{toc}{section}{#1}%
}

\newcommand{\responseheading}[1]{%
    \vspace{0.24cm}
    \par\noindent{\large\bfseries #1}\par
    \vspace{0.03cm}
    \noindent\rule{\linewidth}{0.35pt}\par
    \vspace{0.06cm}
    \noindent\ignorespaces
}

\newcommand{\responsespecificheading}[1]{%
    \vspace{0.30cm}
    \par\noindent{\Large\bfseries #1}\par
    \vspace{0.04cm}
    \noindent\rule{\linewidth}{0.5pt}\par
    \vspace{0.08cm}
    \noindent\ignorespaces
}

\newtcolorbox{commentboxinner}[1]{
    enhanced,
    breakable,
    colback=black!5!white,
    colframe=black!60!white,
    colbacktitle=black!60!white,
    title={#1},
    fonttitle=\bfseries,
    fontupper=\itshape
}

\newenvironment{commentbox}[1]{%
    \par\phantomsection\addcontentsline{toc}{subsection}{#1}%
    \begin{commentboxinner}{#1}%
}{%
    \end{commentboxinner}%
}


\newtcolorbox{planbox}[1]{
    enhanced,
    breakable,
    colback=blue!3!white,
    colframe=blue!55!black,
    colbacktitle=blue!55!black,
    title={#1},
    fonttitle=\bfseries
}

\begin{document}

\begin{center}
{\Large \textbf{Response to the Editor and Reviewers}}\\[0.2cm]
\textbf{Manuscript No.:} IPM-D-26-02888\\[0.2cm]
\textbf{Manuscript Title:}\\
\textit{Fa-DPO: Faithfulness-Aware Direct Preference Optimization with Risk-Adaptive Preference Learning for Medical Question Answering}
\end{center}

\vspace{0.2cm}

Dear Editors,

Thank you for inviting us to revise and resubmit our manuscript. We sincerely thank the Editor-in-Chief, the Associate Editor, and the reviewers for the prompt handling
of our manuscript and for their valuable and constructive comments. These suggestions
have been highly helpful in improving our work. After receiving the feedback, we
immediately organized the revision process and systematically improved
the manuscript in terms of motivation, methodological design, experimental
validation, terminology, and overall consistency.

In this revised response, we address the editor's and reviewers' suggestions point by
point using a consistent two-part structure for each comment: \textit{Response and
Revision Summary} and \textit{Revision Details}. The summary explains the concern and
the main revision, while the details group each revision category directly with the
corresponding manuscript location, passage, table, or verification material. This
merged structure avoids repeating a separate high-level action list before the
specific manuscript-change evidence. We
sincerely hope that the revised manuscript and this detailed point-by-point response
adequately address the reviewers' concerns and further strengthen the manuscript.

\vspace{0.2cm}
Sincerely,

The authors


\responseheading{Revision Mark Legend}

The manuscript changes below retain the revision marks. `\deleted{Red strikethrough}` marks removed or replaced wording, and `\added{blue text}` marks newly added or rewritten wording.

\newpage
\tableofcontents
% \newpage
\newpage

\tocsection{Associate Editor Guidance and Overview of Major Revisions}
\begin{commentbox}{E.1 Associate Editor Guidance}
The Associate Editor notes that the reviewers are supportive of the research but believe
the paper would benefit from additional explanation and clarification of the main
contributions. The response should be detailed and should highlight the changes.
\end{commentbox}

\responseheading{Response and Revision Summary}
We thank the Associate Editor for asking us to make the contribution clearer. We
revised the manuscript to present Fa-DPO as a training-stage method for
evidence-grounded MQA that uses ARU/NLI diagnostic risk signals to focus DPO on
localized evidence-use and reasoning--answer consistency failures. We organized the
major revisions into three groups:
\emph{Conceptual clarification}, which separates evidence faithfulness from
reasoning--answer consistency and clarifies ARU scoring and the DPO objective;
\emph{validation supplements}, which add clinician audits, parser/verifier checks,
controls, and statistical reporting; and \emph{calibrated result framing}, which
narrows the evidentiary claim and states the limits of process metrics, dataset
bias, taxonomy coverage, computational overhead, and clinical-safety interpretation.


\responseheading{Revision Details}
\textbf{Abstract: localized evidence-use errors and Fa-DPO scope.}
\begin{displayquote}
\deleted{However, DPO's sequence-level objective will mask localized faithfulness errors in long, mostly correct clinical responses.}

\added{However, DPO's sequence-level objective can dilute the suppressive signal from localized evidence-use errors in long, mostly correct clinical responses.}

\added{Fa-DPO identifies localized evidence-use and reasoning--answer consistency failures in long clinical reasoning chains and uses them to strengthen learning signals for risky reasoning spans and high-risk responses.}
\end{displayquote}

\textbf{Introduction: controlled local contrasts and ARU/NLI risk scores.}

\begin{displayquote}
\deleted{To provide clinically meaningful negative responses for preference learning, we first simulate realistic clinical reasoning failures to construct negative responses.}

\added{To build preference data with clear local contrasts, we construct rejected responses containing controlled, clinically plausible failures in evidence use and reasoning--answer consistency.}

\deleted{To localize where these responses deviate from the evidence, we decompose each negative response into atomic reasoning units (ARUs).}

\added{We then decompose each response into Atomic Reasoning Units (ARUs), route each ARU to relevant supporting premises, and verify it with an NLI model to obtain ARU-level evidence-use risk scores.}

\deleted{We measure faithfulness risk of atomic reasoning units through multi-source evidence routing and NLI-based verification.}

\added{Fa-DPO incorporates these scores into a faithfulness-aware DPO objective by retaining high-risk negatives, upweighting tokens in risky spans, and assigning larger margins to clinically higher-risk response pairs.}

\end{displayquote}

\textbf{Research objective: evidence faithfulness and reasoning-answer consistency.}
\begin{displayquote}
\deleted{We define faithfulness as the extent to which intermediate clinical inferences are supported by retrieved evidence, avoid contradictions, and remain consistent with the final answer.}

\added{Following this general view, we define evidence faithfulness as the extent to which intermediate clinical inferences are supported by retrieved evidence and do not contradict it. We further evaluate reasoning-answer consistency as a related diagnostic check, because the final answer should follow from the preceding reasoning.}
\end{displayquote}

\textbf{ARU/DPO objective: local credit assignment.}
\begin{displayquote}
\deleted{To localize faithfulness risk before measurement, ARU decomposition converts each candidate negative response into minimal reasoning units.}

\added{To localize verification before scoring, ARU decomposition operationalizes atomic-claim decomposition as a medical-reasoning unit for optimization. We formalize ARUs as reasoning-function-labeled atomic clinical reasoning claims.}

\deleted{Standard DPO uses the same margin for every sample, whereas our margin increases with the overall severity of the retained negative response.}

\added{Standard DPO uses no explicit sample-dependent additive margin, whereas Fa-DPO introduces a bounded risk-dependent margin for the rejected response.}

\added{The normalization is therefore a local credit-assignment device: it redistributes a fixed rejected-side token mass toward high-risk spans, rather than increasing the total token mass of the rejected response.}
\end{displayquote}


\newpage
\tocsection{Reviewer \#1}

\begin{commentbox}{R1.1 Faithfulness Definition}
The paper defines faithfulness as intermediate clinical inferences being supported by
retrieved evidence, containing no contradictions, and being consistent with the final
answer. However, in academic terminology, faithfulness conventionally refers only to
the consistency between generated content and the given evidence; final-answer
consistency is closer to consistency or correctness.
\end{commentbox}

\responseheading{Response and Revision Summary}
We agree with the reviewer. The revised manuscript now defines evidence faithfulness
as support by, and non-contradiction with, the provided evidence. It treats
reasoning-answer consistency as a related reasoning-quality check rather than part of
the narrow evidence-faithfulness definition.

The method name remains faithfulness-aware because Fa-DPO's main training signal
targets evidence grounding in intermediate reasoning. The ARU-level verification
score is now described more precisely: for evidence-routed ARUs it estimates
evidence-use risk, and for claim ARUs it estimates whether the claim or final answer
follows from preceding reasoning.

\responseheading{Revision Details}
\textbf{Sec.~2: narrow evidence-faithfulness definition.}

\begin{displayquote}
\deleted{We define faithfulness as the extent to which intermediate clinical inferences are supported by retrieved evidence, avoid contradictions, and remain consistent with the final answer.}

\added{Following this general view, we define evidence faithfulness as the extent to which intermediate clinical inferences are supported by retrieved evidence and do not contradict it. We further evaluate reasoning-answer consistency as a related diagnostic check, because the final answer should follow from the preceding reasoning.}
\end{displayquote}

\textbf{Failure taxonomy: separate evidence-use failures from reasoning-answer consistency.}

\begin{displayquote}
\deleted{Table~\ref*{ms-tab:error_types} defines four localized faithfulness failure types along the evidence-to-reasoning-to-answer chain. F1--F3 target the evidence-to-reasoning interface $(x,G)\rightarrow c$, while F4 targets the reasoning-to-answer interface $c\rightarrow a$.}

Table~\ref*{ms-tab:error_types} \added{defines four localized process-level failure types along the evidence-to-reasoning-to-answer chain. F1--F3 are evidence-use faithfulness failures at the evidence-to-reasoning interface $(x,G)\rightarrow c$, including evidence omission, evidence misinterpretation, and overinterpretation. F4 is a reasoning-answer consistency failure at the reasoning-to-answer interface $c\rightarrow a$, where the final answer is not supported by, or conflicts with, the preceding reasoning process.}
\end{displayquote}

\textbf{Method section: define the ARU/NLI score by diagnostic role.}

\begin{displayquote}
\deleted{Fa-DPO integrates faithfulness risk into DPO through three linked mechanisms. First, response-level faithfulness risk determines which constructed negatives are retained as rejected responses.}

\added{Fa-DPO uses one ARU-level verification score to guide DPO through three linked mechanisms. For evidence-routed ARUs, this score measures evidence-support risk; for claim ARUs, it measures whether the claim or final answer follows from the preceding reasoning. First, the response-level summary of the ARU-level verification scores determines which constructed negatives are retained as rejected responses.}
\end{displayquote}

\textbf{Abstract and framework overview: clean up mixed metric and stage labels.}

\begin{displayquote}
\deleted{Fa-DPO improves same-answer process faithfulness}\added{Fa-DPO improves automatic same-answer process metrics}

\deleted{ARU-Level Faithfulness Risk Measurement}\added{ARU-Level Risk Scoring}

\deleted{measure faithfulness risk}\added{To provide fine-grained risk-scoring signals for downstream DPO training, we decompose each quality-controlled candidate negative response into ARUs and assign each ARU a verifier-derived score through multi-source evidence routing and NLI-based verification. For evidence-routed ARUs, the score estimates evidence-use risk; for claim ARUs, including the final-answer claim, it estimates reasoning--answer consistency risk.}
\end{displayquote}

\textbf{Sec.~6.1.3: group automatic metrics by diagnostic role.}

\begin{displayquote}
\deleted{We measure process faithfulness with four complementary metrics}\deleted{report two groups of verifier-derived process metrics}\added{We report four ARU/NLI-based metrics, grouped by diagnostic role.}

\added{The evidence-support metrics are} \textbf{Support Rate} \added{and} \textbf{Unsupported Medical Reasoning Rate}. \added{The latter is not the complement of Support Rate because it is computed only over clinically inferential warrant and differentiation units.}

\added{The evidence-conflict metric is} \textbf{Contradiction Rate}.

\deleted{Claim Inconsistency Rate}\added{The reasoning-answer consistency metric is \textbf{Reasoning-Answer Inconsistency Rate}, the verifier-labeled fraction of claim ARUs unsupported or contradicted by the preceding ARUs in the same response, $\textit{Reasoning-Answer Inconsistency Rate} = N_{\text{inc}}^{\text{C}} / N_{\text{C}}$.}
\end{displayquote}

\textbf{Local terminology cleanup.}
\begin{displayquote}
\deleted{reasoning--answer faithfulness}\added{reasoning-answer consistency}

\deleted{answer-faithfulness}\added{reasoning-answer consistency}

\deleted{Cross-Pair Same-Answer Process-Faithfulness Protocol}\added{Cross-Pair Same-Answer Process-Metric Protocol}
\end{displayquote}

\newpage
\begin{commentbox}{R1.2 ARU Definition}
Atomic Reasoning Unit (ARU) is a custom term, but the paper does not clearly explain
its difference from existing concepts such as atomic claims or reasoning steps in
MedRAGChecker.
\end{commentbox}

\responseheading{Response and Revision Summary}
We thank the reviewer for pointing out that the relationship between Atomic Reasoning
Units (ARUs) and existing concepts such as atomic claims and reasoning steps was not
sufficiently clear. We agree that the previous manuscript did not adequately position
ARUs with respect to prior atomic fact/claim decomposition and biomedical RAG
verification work.

In the revision, we clarify ARUs along three aspects. First, ARUs inherit the
decompose-then-verify principle from atomic fact/claim-based factuality evaluation.
Second, ARUs extend generic atomic claims by adding a clinical reasoning-function
label and an evidence route, so that each unit can be checked against the evidence
source appropriate to its role. Third, ARUs are used in Fa-DPO not only as evaluation
units but also as optimization units: their verification scores and token spans define
local risk signals for span weighting and response-level preference margins. This
clarification distinguishes ARUs from both generic atomic claims and coarse reasoning
steps.


\responseheading{Revision Details}
\textbf{Abstract and Highlights.} We clarified ARU at its first appearances:

\begin{displayquote}
\deleted{We then decompose each response into atomic reasoning units, verify each unit through role-specific evidence routing and natural language inference (NLI)-based checking, and use the resulting risk scores to reweight high-risk spans and set risk-adaptive margins.}

\added{Responses are decomposed into reasoning-function-labeled Atomic Reasoning Units (ARUs), verified via natural language inference (NLI), to guide risk-aware preference optimization.}

\deleted{We measure faithfulness risk of atomic reasoning units (ARU) through multi-source evidence routing and natural language inference(NLI)-based verification.}

\added{We introduce an evidence-use risk estimation through multi-source evidence routing and natural language inference (NLI)-based verification, enabling localized detection of unsupported reasoning.}
\end{displayquote}

\textbf{Introduction and contribution list.} We changed the framing to operationalizing atomic-claim decomposition:

\begin{displayquote}
\deleted{To localize where these responses deviate from the evidence, we decompose each negative response into atomic reasoning units (ARUs).}

\added{We then decompose each response into Atomic Reasoning Units (ARUs), route each ARU to relevant supporting premises, and verify it with an NLI model to obtain ARU-level evidence-use risk scores.}

\deleted{We measure faithfulness risk of atomic reasoning units through multi-source evidence routing and NLI-based verification.}

\added{We introduce an ARU-level evidence-use risk estimation through multi-source evidence routing and NLI-based verification, enabling localized detection of unsupported reasoning.}
\end{displayquote}

\textbf{Related Work.} We added prior-work positioning against atomic-claim factuality evaluation and biomedical RAG verification:

\begin{displayquote}
\added{RAG evaluation makes this evidence-use problem operational through context relevance, answer faithfulness, and answer relevance~\cite{saadfalcon2024ares}.}

\added{Self-checking, verifier-RAG, atomic factuality evaluation, and biomedical RAG checkers decompose outputs into claims and verify them against generated, retrieved, or external evidence~\cite{manakul2023selfcheckgpt, phasook2025vereafine, min-etal-2023-factscore, wei2024longform, song-etal-2024-veriscore, ji2026medragchecker}.}

\added{Existing verification signals are usually used for evaluation, filtering, consensus checking, or post-hoc diagnosis.}

\added{Fa-DPO adopts the decompose-then-verify principle but changes its role: ARUs encode clinical reasoning function, evidence route, local verification score, and token span, allowing verification results to become DPO weights and margins rather than only diagnostic labels.}
\end{displayquote}

\textbf{Sec.~5.4.1.} We replaced the generic ARU wording with a formal definition that states what ARUs inherit from atomic claims and what they add for Fa-DPO:

\begin{displayquote}
\deleted{To localize faithfulness risk before measurement, ARU decomposition converts each candidate negative response into minimal reasoning units. Given a candidate negative response $\tilde{y}_l = c^{-} \oplus a^{-}$, we segment the reasoning trace $c^{-}$ into Atomic Reasoning Units (ARUs) and append the extracted final answer $a^{-}$ as a special claim ARU. This yields the ordered ARU sequence $\mathcal{V}(\tilde{y}_l) = (v_1, v_2, \dots, v_L)$, where $v_1, \dots, v_{L-1}$ are reasoning-derived ARUs and $v_L$ is the final-answer claim ARU. The ARU decomposition protocol appendix details the parser protocol, retry rules, and quality-control procedure.}

\added{To localize verification before scoring, ARU decomposition operationalizes atomic-claim decomposition as a medical-reasoning unit for optimization. We formalize ARUs as reasoning-function-labeled atomic clinical reasoning claims. This definition inherits the atomic-claim decomposition principle, but is specialized for evidence-grounded medical reasoning and preference optimization. Each ARU specifies not only a minimal checkable claim, but also its clinical reasoning function, the evidence route against which it should be verified, the local evidence-support or reasoning-consistency score, and the token span to which this score is assigned in the Fa-DPO objective. Thus, an ARU is not merely a renamed atomic claim; it is an optimization-oriented unit used to connect local verification scores with the Fa-DPO objective. Given a candidate negative response $\tilde{y}_l = c^{-} \oplus a^{-}$, we segment the reasoning trace $c^{-}$ into ARUs and append the extracted final answer $a^{-}$ as a special claim ARU. This yields the ordered ARU sequence $\mathcal{V}(\tilde{y}_l) = (v_1, v_2, \dots, v_L)$, where $v_1, \dots, v_{L-1}$ are reasoning-derived ARUs and $v_L$ is the final-answer claim ARU.}
\end{displayquote}

\newpage
\begin{commentbox}{R1.3 DPO Equation Sign}
In standard DPO, the log-probability ratio for positive samples should be the policy
model relative to the reference model. In Equation (13), the order appears reversed,
which flips the sign relationship between positive and negative samples.
\end{commentbox}

\responseheading{Response and Revision Summary}
We thank the reviewer for identifying this sign-convention issue. We agree that the
DPO ratio should be interpreted in the policy-over-reference direction. We checked the
DPO preliminaries, Fa-DPO objective, and training-logit implementation, and confirmed
that the implemented convention was already policy-over-reference. The manuscript
presentation, however, made this direction less transparent. We therefore revised the
preliminaries to define a single policy-over-reference log-ratio $\rho_\theta(x,y)$,
rewrote the standard DPO and Fa-DPO objectives using this notation, and clarified that
$\Delta_m$ is computed in the same direction. These revisions clarify the equation sign
and do not change the training procedure, dataset construction, or experimental
results.

This sign-convention revision is part of a broader DPO objective-presentation revision package also related to R1.4, R1.5, and R2.9. Together, these changes make the Fa-DPO objective internally consistent: the policy-over-reference direction is fixed, the rejected-side token weighting is defined on the same log-ratio scale, and the margin is treated as a bounded pairwise offset.


\responseheading{Revision Details}
\textbf{DPO preliminaries.}
In Sec.~4, we made the policy-over-reference direction explicit before presenting the
standard DPO objective:
\begin{displayquote}
\deleted{The objective function is derived from the Bradley-Terry model and is defined as:}

\added{Following standard DPO, we define the policy-over-reference log-ratio as}

\[
{\color{blue}
\rho_\theta(x,y)
=
\log \pi_\theta(y|x)
-
\log \pi_{\text{ref}}(y|x).
}
\]

\added{Using this notation, the objective function is derived from the Bradley-Terry model~\cite{bradley1952rank} and is defined as}
\[
\mathcal{L}_{\text{DPO}}
=
- \mathbb{E}\left[
\log \sigma\left(
\beta\left[\rho_\theta(x,y_w)-\rho_\theta(x,y_l)\right]
\right)
\right].
\]
\end{displayquote}

\textbf{Token-weighted rejected log-ratio.}
In Sec.~5.5.2, we rewrote $\Delta_m$ in token-level difference form:
\begin{displayquote}
\[
{\color{blue}
\Delta_m(x,y_l,\mathbf{s}_l)
=
\frac{T_l}{\sum_{t=1}^{T_l} m_t}
\sum_{t=1}^{T_l}
m_t
\Big[
\log \pi_\theta(z^{(l)}_t \mid x, z^{(l)}_{<t})
-
\log \pi_{\text{ref}}(z^{(l)}_t \mid x, z^{(l)}_{<t})
\Big].
}
\]

\added{The bracketed token-level ratio is computed in the same policy-over-reference direction as $\rho_\theta(x,y)$. Because $\rho_\theta(x,y)$ uses unnormalized autoregressive sequence log-probabilities, Eq.~\ref*{ms-eq:weighted_rejected_logratio} is kept on the same token-sum scale.}
\end{displayquote}

\textbf{Unified Fa-DPO objective.}
In Sec.~5.5.4, we rewrote the preferred-response term and added the sign-convention
explanation:
\begin{displayquote}
\[
{\color{blue}
\begin{split}
\mathcal{L}_{\text{Fa-DPO}}(\pi_\theta; \pi_{\text{ref}})
&= -\mathbb{E}_{(x, y_w, y_l, \mathbf{s}_l) \sim \mathcal{D}_{\text{Fa}}} \Bigg[
\log \sigma \Bigg(
\beta \rho_\theta(x,y_w)
- \beta \Delta_m(x,y_l,\mathbf{s}_l)
- \delta(y_l)
\Bigg) \Bigg],
\end{split}
}
\]

\added{Both $\rho_\theta(x,y_w)$ and $\Delta_m(x,y_l,\mathbf{s}_l)$ are computed in the policy-over-reference direction. Therefore, minimizing the loss increases the preferred-response log-ratio while decreasing the rejected-response log-ratio, with the rejected-side decrease concentrated on high-risk spans through $m_t$.}
\end{displayquote}

\newpage
\begin{commentbox}{R1.4 Normalization Justification}
The normalization design of Fa-DPO lacks theoretical justification. The normalization
factor in Equation (11) is said to keep the scale comparable to the unweighted sequence
log-probability ratio, but it changes gradient properties and may amplify gradients
when only a small number of high-risk tokens exist.
\end{commentbox}

\responseheading{Response and Revision Summary}
We agree that the original manuscript did not sufficiently justify the normalization in Eq.~(11). We revised the method section to state exactly what the normalization preserves. Let the normalized effective token weight be
\[
\omega_t=\frac{T_l m_t}{\sum_j m_j}.
\]
Then $\sum_t \omega_t=T_l$, and $\omega_i/\omega_j=m_i/m_j$ within the same rejected response. Thus, the normalization keeps the total rejected-side token mass comparable to the unweighted rejected-response log-ratio, while preserving the ARU-risk mask's relative emphasis on high-risk spans.

We also clarify that this is a local credit-assignment device, not a claim that the absolute response-level log-ratio ordering is preserved across all rejected responses. In the uniform-mask case, Eq.~(11) reduces to the standard rejected-response log-ratio.

To support the design empirically, we added a normalization diagnostic comparing uniform weighting, no normalization, the original normalized weighting, and capped normalization. The original normalized weighting gives the best overall row in the representative setting while keeping training stable. In contrast, removing normalization increases instability events and loss spikes. These results support the normalization as a scale-stabilizing local credit-assignment choice.

This revision also addresses the related clarification requested in R2.9. We now state that Eq.~(11) preserves the total effective token mass within each rejected response and the within-response relative emphasis of the ARU-risk mask. It does not preserve the absolute ordering of raw response-level log-ratio magnitudes across different rejected responses.


\responseheading{Revision Details}
\textbf{Main text, token-mask definition.}
\begin{displayquote}
\deleted{$m_t=\alpha$ for high-risk ARU tokens}

{\color{blue}$1+\alpha s_i, \text{if } t \in \mathcal{T}(v_i) \text{ for some } i \text{ and } s_i > \tau$}

\deleted{where $\alpha > \epsilon > 0$. Tokens covered by ARUs above the shared high-risk threshold $\tau$ receive weight $\alpha$, and all other tokens receive weight $\epsilon$.}

\added{where $\alpha>0$ and $\epsilon>0$. Tokens covered by ARUs above the shared high-risk threshold $\tau$ receive a continuous risk-scaled weight, and all other tokens receive the background weight $\epsilon$. Because $s_i\in[0,1]$, the raw token weights satisfy $m_t\in[\epsilon,1+\alpha]$.}
\end{displayquote}

\textbf{Main text, ``ARU-Risk-Guided Rejected-Response Weighting in DPO,''
immediately after Equation (11).}

\begin{displayquote}
\deleted{keeps the scale comparable to the unweighted sequence log-ratio, while the mask increases the contribution of high-risk spans}

\added{Equivalently, Eq.~\ref*{ms-eq:weighted_rejected_logratio} uses the normalized effective token weight
$\omega_t=T_lm_t/\sum_{j=1}^{T_l}m_j$. These weights satisfy
$\sum_{t=1}^{T_l}\omega_t=T_l$ and $\omega_i/\omega_j=m_i/m_j$ within the same
rejected response. Thus, the normalization preserves the same total effective token
mass as the unweighted rejected-response sequence log-ratio, while preserving the
relative emphasis introduced by the ARU-risk mask.}

\added{In the uniform case $m_t=c$ for all tokens, $\omega_t=1$ and the expression
reduces exactly to the standard rejected-response sequence log-ratio. The
normalization is therefore a local credit-assignment device: it redistributes a fixed
rejected-side token mass toward high-risk spans, rather than increasing the total
token mass of the rejected response.}

\added{Sparse high-risk masks can increase the relative weight of individual
high-risk tokens, but because $m_t\in[\epsilon,1+\alpha]$, the normalized effective
token weight is bounded by $\omega_t\le (1+\alpha)/\epsilon$; therefore, the weighting
cannot produce unbounded effective token weights. This fixed-mass property is a
statement about rejected-side token weights, not a guarantee that gradient norms are
constant, because log-probability gradients still depend on the token distribution,
model confidence, and current parameters.}
\end{displayquote}

\textbf{Ablation appendix: separate local token weighting from sample-level
severity.}
\begin{displayquote}
\added{All variants use the same retained preference pairs, training setup, and risk-adaptive margin; only the rejected-side token weighting rule changes. Therefore, the uniform-weighting row is a retained-pool margin-only, no-localization control rather than Standard DPO-Retained-34K.}
\end{displayquote}

\textbf{Appendix, ``Token-Weight Normalization Diagnostic'' table.}
\begin{displayquote}
\begin{center}
\scriptsize
\setlength{\tabcolsep}{2pt}
\captionof{table}{\added{Token-weight normalization diagnostic on the representative II-Medical-8B setting. Process metrics are automatic ARU/NLI-based percentages on the MedCaseReasoning same-answer subset, and Avg. Acc. is the macro-average over the four main accuracy benchmarks. All variants had zero NaN/Inf steps. For the uniform-weighting row, Avg. Acc. follows the corresponding \textit{w/o Token-Level Dynamic Reweighting} ablation, while the process and stability columns are diagnostic-run measurements under this normalization protocol.}}
\begin{adjustbox}{max width=\linewidth}
\begingroup\color{blue}
\begin{tabular}{lcccccccccc}
\toprule
\textbf{Variant} & \textbf{Avg. Acc.} & \textbf{Support} & \textbf{Contrad.} & \textbf{Unsupported} & \textbf{Reason.-Ans. Inc.} & \textbf{Mean/Max Grad} & \textbf{Instab.} & \textbf{Loss Spikes} & \textbf{High-Risk p95/Max Weight} & \textbf{High-Risk Weight Share} \\
\midrule
Uniform weighting (+ margin) & 61.7 & 80.6 & 7.3 & 25.9 & 22.5 & 1.86/5.1 & 0 & 1 & 1.00/1.00 & 25.4 \\
No normalization & 62.4 & 82.9 & 6.4 & 23.8 & 20.7 & 2.74/18.6 & 7 & 9 & 4.80/21.3 & 68.7 \\
Original normalized weighting & \textbf{63.2} & \textbf{84.0} & \textbf{5.2} & \textbf{20.9} & \textbf{18.5} & 2.05/6.3 & 0 & 2 & 3.20/8.9 & 62.5 \\
Capped normalization & 63.1 & 83.7 & 5.4 & 21.3 & 18.8 & 1.98/5.7 & 0 & 1 & 3.00/7.5 & 61.8 \\
\bottomrule
\end{tabular}
\endgroup
\end{adjustbox}
\end{center}
\end{displayquote}

\textbf{Appendix, ``Training Hyperparameters and Run Protocol,'' subsection
``Fa-DPO Coefficients and Evaluation Rule.''}

\begin{displayquote}
\added{With these default values and $s_i\in[0,1]$, the raw mask satisfies
$m_t\in[0.1,3.0]$, so the normalized effective token weight is bounded above by
$(1+\alpha)/\epsilon=30$, while $\sum_t\omega_t=T_l$ keeps the rejected-side total
token mass fixed for each rejected response.}
\end{displayquote}

\newpage
\begin{commentbox}{R1.5 Additive Margin and DPO Theory}
Does the additive margin $\delta(y_l)$ in Equation (12) preserve reward monotonicity?
Does it satisfy the existence condition of the optimal policy for DPO? The paper lacks
discussion and connection to DPO theory.
\end{commentbox}

\responseheading{Response and Revision Summary}
We thank the reviewer for pointing out that the additive margin needed a clearer
connection to DPO theory. We revised the manuscript to clarify that
$\delta(y_l)$ is a bounded, risk-dependent pairwise margin applied to the empirical
DPO logit, not a new global response-level reward.

This revision is coordinated with the sign-convention and normalization revisions in R1.3 and R1.4. Fa-DPO keeps the standard policy-over-reference log-ratio basis, while adding two bounded empirical modifications: normalized rejected-side token weighting and a risk-dependent pairwise margin.

The revised text now shows that, because $R(y_l)\in[0,1]$ and
$\eta,\kappa_{\delta},c_{\mathrm{sat}}>0$, the mapping
$\delta(y_l)=\eta\min(\kappa_{\delta}R(y_l),c_{\mathrm{sat}})$ is monotone
non-decreasing and bounded by $\eta c_{\mathrm{sat}}$. Saturation can make multiple
high-risk negatives share the same maximum margin, but it cannot reverse their risk
ordering. Thus, the margin preserves monotonicity of the required
preferred-minus-rejected separation, not monotonicity of a global reward ranking over
all possible responses.

We also added a more limited DPO-theoretic interpretation. Standard DPO can be written
through the implicit reward $r_\theta(x,y)=\beta\rho_\theta(x,y)+C(x)$, where
$\rho_\theta$ is the policy-over-reference log-ratio and $C(x)$ cancels in pairwise
differences. Fa-DPO keeps this log-ratio basis but modifies the empirical pairwise
logit through a normalized token-weighted rejected term and a bounded rejected-side
margin. We therefore describe Fa-DPO as a bounded large-margin variant of DPO, not as
an unmodified standard DPO objective or a new closed-form reward-model derivation. The
bounded margin does not by itself introduce unbounded reward scale or support mismatch,
but we do not claim that Fa-DPO preserves the standard DPO optimal-policy theorem
unchanged.

\responseheading{Revision Details}
\textbf{Sec.~``Sample-Level Risk-Adaptive Preference Margin in DPO.''}
\begin{displayquote}
\deleted{Standard DPO uses the same margin for every sample, whereas our margin increases with the overall severity of the retained negative response.}

\added{Standard DPO uses no explicit sample-dependent additive margin, whereas Fa-DPO introduces a bounded risk-dependent margin for the rejected response.}

\added{Because $R(y_l)\in[0,1]$ and $\eta,\kappa_{\delta},c_{\text{sat}}>0$, $\delta(y_l)$ is monotone non-decreasing in rejected-response risk: if $R(y_l^{(1)})\le R(y_l^{(2)})$, then $\delta(y_l^{(1)})\le\delta(y_l^{(2)})$.}

\added{The margin should be interpreted as a pairwise large-margin offset in the empirical DPO logit, not as a new global response-level reward.}
\end{displayquote}

\textbf{Sec.~``Unified Fa-DPO Objective.''}
\begin{displayquote}
\added{In standard DPO, the implicit reward can be written as $r_\theta(x,y)=\beta\rho_\theta(x,y)+C(x)$, where $C(x)$ cancels in pairwise differences.}

\added{This objective remains tied to standard DPO through the policy-over-reference log-ratio. In standard DPO, the implicit reward can be written as $r_\theta(x,y)=\beta\rho_\theta(x,y)+C(x)$, where $C(x)$ cancels in pairwise differences.}

\added{Fa-DPO keeps this log-ratio basis but modifies the empirical pairwise logit by using a normalized token-weighted rejected term and a bounded rejected-side margin. We therefore view Fa-DPO as a bounded large-margin variant of DPO rather than an unmodified standard DPO objective or a new closed-form reward-model derivation.}

\added{Fa-DPO's additive margin is finite, pair-dependent, and does not change the support of $\pi_{\mathrm{ref}}$; therefore, the margin does not by itself introduce an unbounded reward-scale or support-mismatch problem. We make this boundedness claim for the empirical large-margin objective, not as a proof that Fa-DPO preserves the standard DPO optimal-policy theorem unchanged.}
\end{displayquote}

\textbf{Appendix~``Fa-DPO Coefficients and Evaluation Rule.''}
\begin{displayquote}
\added{The scale $\eta$ keeps the additive margin moderate relative to the base DPO preference term: under the default values, the highest-risk retained negatives require at most an additional $0.5$ preferred-minus-rejected score separation in the Fa-DPO logit.}
\end{displayquote}

\newpage
\begin{commentbox}{R1.6 NLI Verifier Reliability}
The paper relies on an NLI model to score ARUs. General NLI models may be inaccurate
for professional medical text. Key calibration metrics are not in the main text. False
positives could wrongly penalize correct reasoning and create negative training.
\end{commentbox}

\responseheading{Response and Revision Summary}
We agree that verifier false positives can create misleading training signals. The
revision therefore treats the NLI verifier as an in-pipeline diagnostic proxy rather
than clinical ground truth, reports the exact checkpoint, and adds calibration,
held-out, natural-output, and clinician false-positive checks.

The revised claim is bounded: the verifier supports ARU/NLI-based process analysis,
while the detailed metrics, confusion matrix, and audit evidence are reported below.


\responseheading{Revision Details}
\textbf{Appendix NLI validation title and validation table.}
\begin{displayquote}
\textit{Appendix title:} \added{NLI Verifier Validation Protocol}

\textit{Manuscript table-caption excerpt:} \added{Verifier-level NLI validation on the held-out in-pipeline medical evaluation split ($n=160$). A separate stratified development split ($n=40$) is reserved and excluded from the reported held-out metrics.}

\textit{Protocol excerpt:} \added{The in-pipeline protocol samples $n=200$ medical ARU-premise pairs with stratification over route type, ARU reasoning function, predicted label, and difficulty band. We reserve 40 pairs as a small stratified development split and exclude them from all reported held-out metrics; the final verifier results are reported on the remaining 160 held-out in-pipeline medical evaluation pairs.}

{\raggedright
\textit{Implementation excerpt:} \added{All main-text mechanism analyses use the II-Medical-8B~\cite{intelligentinternet2025iimedical8b} matched pair with the same retrieval pipeline and the public DeBERTa-v3-large NLI checkpoint}
\par}

\textit{Manuscript footnote excerpt:} \added{Hugging Face model page \url{https://huggingface.co/MoritzLaurer/DeBERTa-v3-large-mnli-fever-anli-ling-wanli}.}

{\color{blue}Accuracy \& 0.838}

{\color{blue}Macro-F1 \& 0.827}

{\color{blue}Expected calibration error \& 0.064}

{\color{blue}Brier score \& 0.213}

\added{ECE is computed as 10-bin max-confidence ECE with equal-width bins on $[0,1]$, and the Brier score is the three-class multiclass Brier score averaged over examples.}

{\color{blue}Gold-entailment $\rightarrow$ non-entailment rate \& 4/64 (6.3\%)}

{\color{blue}Gold-entailment $\rightarrow$ contradiction rate \& 0/64 (0.0\%)}

{\color{blue}Double-annotation agreement \& $\kappa=0.74$}
\end{displayquote}

\textbf{Appendix NLI validation protocol: remove the external-artifact pointer.}
\begin{displayquote}
\deleted{To keep the paper appendix compact, we leave class-wise precision/recall, route-wise breakdowns, reliability curves, and representative failure cases to the anonymous artifact files.}

\added{Table~\ref*{ms-tab:nli_validation_template} reports class-wise precision/recall/F1 together with accuracy, macro-F1, ECE, Brier score, and annotation agreement. The main text reports the held-out summary metrics and the separate natural-output shift analysis, while Table~\ref*{ms-tab:nli_confusion_template} reports the corresponding three-way confusion matrix for the in-pipeline split. To keep the appendix focused, the remaining appendix validation emphasizes the confusion matrix, natural-output shift analysis, and false-positive audits rather than exhaustive route-wise breakdowns, reliability curves, or representative failure-case listings.}
\end{displayquote}

\textbf{Sec.~6.6.2 verifier-validation interpretation.}
\begin{displayquote}
\deleted{These results support using the verifier as the ARU-level supervision and evaluation signal in Fa-DPO, with particularly strong performance on entailment and solid performance on the two non-entailment classes. We therefore use it as the verifier-side foundation for the faithfulness scores in Fa-DPO, while leaving the full confusion matrix and additional protocol details to the appendix.}

\added{Pipeline audits validate the automatic scoring components. The held-out parser meets the main span and role targets at the primary IoU threshold, but role macro-F1 and final-answer claim-span F1 limit exact role-wise attribution. The routed NLI verifier reaches 0.838 accuracy and 0.827 macro-F1 on the held-out in-pipeline medical split, while natural-output stress testing shows moderate degradation, including a high-risk false-negative increase from 0.118 to 0.172. Clinician false-positive and training-mask audits further show that high-risk flags are mostly meaningful but can include clinically acceptable evidence-unsupported cases. Table~\ref*{ms-tab:pipeline_validation_summary} summarizes these component-level checks. These results justify using the automatic metrics as diagnostic indicators, not as clinical ground truth; detailed protocols and tables are in Appendices~\ref*{ms-app:aru_parser_validation} and~\ref*{ms-app:nli_validation}.}
\end{displayquote}

\textbf{Sec.~5.4.2 scoring rule: neutral is weaker than contradiction.}
\begin{displayquote}
\added{The contradiction term captures warrants that conflict with retrieved evidence, while the neutral term assigns a weaker risk to warrants that are not explicitly supported.}

\added{The combined premise makes the score sensitive to exclusions that contradict either the case evidence or the retrieved medical evidence, while still treating mere lack of explicit support as a weaker signal than contradiction.}
\end{displayquote}

\textbf{Sec.~6.6.2 clinician false-positive audit.}
\begin{displayquote}
\begin{center}
\scriptsize
\setlength{\tabcolsep}{2pt}
\captionof{table}{\added{Clinician false-positive audit of verifier-flagged high-risk ARUs. Clinicians judge whether each flagged ARU is a true error, an acceptable implicit medical inference, or uncertain.}}
\begin{adjustbox}{max width=\linewidth}
\begin{tabular}{lccccccc}
\toprule
\textbf{ARU role} & \textbf{n} & \textbf{True error} & \textbf{Acceptable implicit inference} & \textbf{Uncertain} & \textbf{High-risk precision} & \textbf{False-positive rate} & \textbf{Severe FP rate} \\
\midrule
Observation & 50 & 37 & 9 & 4 & 74.0\% & 18.0\% & 4.0\% \\
Warrant & 50 & 39 & 7 & 4 & 78.0\% & 14.0\% & 2.0\% \\
Differentiation & 50 & 36 & 9 & 5 & 72.0\% & 18.0\% & 6.0\% \\
Claim & 50 & 41 & 7 & 2 & 82.0\% & 14.0\% & 2.0\% \\
\midrule
Overall & 200 & 153 & 32 & 15 & 76.5\% & 16.0\% & 3.5\% \\
\bottomrule
\end{tabular}
\end{adjustbox}
\end{center}
\end{displayquote}

\textbf{Sec.~6.6.2 training-mask false-positive audit.}
\begin{displayquote}
\added{To distinguish high-risk precision from training-mask contamination, we also audit a role- and mask-balanced sample containing both mask-positive and mask-negative ARUs. Clinicians are blinded to the verifier label, probability, risk score, and mask status. The primary metric is the weighted fraction of strict evidence-supported ARUs that are nevertheless mask-positive under the Fa-DPO token-mask rule, where evidence support is adjudicated as direct support or entailment/necessary inference from the provided question context and verification premise. As shown in Table~\ref{tab:training_mask_fp_audit}, this evidence-supported training-mask false-positive rate is 6.9\%, while the severe contradiction-dominant false-positive rate is 0.2\%. Pre-adjudication agreement on the evidence-support labels reaches $\kappa=0.822$. We separately report clinically correct or acceptable ARUs that are not supported by the provided evidence: in this audit, 76/200 ARUs (38.0\%) fall into this clinically acceptable but evidence-unsupported category, with a weighted mask-positive rate of 34.3\%. These cases test the clinical plausibility boundary of the proxy signal but are not counted as gold-entailment false positives. Because the audit is deliberately role- and mask-balanced, these weighted rates are not natural ARU-prevalence estimates; the higher Warrant and Differentiation rates reflect their stronger dependence on inferential support and premise completeness.}

\begin{center}
\scriptsize
\setlength{\tabcolsep}{2pt}
\renewcommand{\arraystretch}{1.08}
\captionof{table}{\added{Clinician audit of Fa-DPO training-mask false positives under a strict evidence-support standard. The audit is balanced by ARU role and training-mask status; weighted rates use inverse sampling weights for each role-mask stratum and are not natural ARU-prevalence estimates.}}
\label{tab:training_mask_fp_audit}
\begin{adjustbox}{max width=\linewidth}
\begin{tabular}{@{}lrrrrrrrcc@{}}
\toprule
\textbf{ARU role} & \textbf{n} & \textbf{Mask+} & \textbf{Mask-} & \shortstack{\textbf{Evidence-}\\\textbf{supported}} & \shortstack{\textbf{Not}\\\textbf{supported}} & \textbf{Contrad.} & \textbf{Uncertain} & \shortstack{\textbf{Weighted}\\\textbf{mask FPR}} & \shortstack{\textbf{Severe}\\\textbf{FPR}} \\
\midrule
Observation & 50 & 25 & 25 & 31 & 15 & 2 & 2 & 6.4\% & 0.0\% \\
Warrant & 50 & 25 & 25 & 12 & 38 & 0 & 0 & 10.1\% & 0.0\% \\
Differentiation & 50 & 25 & 25 & 7 & 42 & 0 & 1 & 14.2\% & 0.0\% \\
Claim & 50 & 25 & 25 & 13 & 27 & 6 & 4 & 8.0\% & 1.6\% \\
\midrule
Overall & 200 & 100 & 100 & 63 & 122 & 8 & 7 & 6.9\% & 0.2\% \\
\addlinespace[0.2em]
\multicolumn{10}{@{}p{0.98\linewidth}@{}}{\added{Clinically acceptable but evidence-unsupported}: \added{76/200 (38.0\%)}; \added{Weighted mask-positive rate 34.3\%}.} \\
\bottomrule
\end{tabular}
\end{adjustbox}
\end{center}
\end{displayquote}

\newpage
\begin{commentbox}{R1.7 Clinician Audit and Filtering Threshold}
Section 5.3.2 states that a stratified sample from the retained DPO negative pool is
audited by clinicians, but response-level filtering occurs before clinician review.
How are clinician review results fed back to adjust $\gamma$? The sampling ratio,
number of clinicians, and agreement metrics are not presented in the main text.
\end{commentbox}

\responseheading{Response and Revision Summary}
We thank the reviewer for pointing out this ambiguity. We clarified that $\gamma$ is
fixed before the final clinician audit: automatic QC and ARU-level scoring are applied
first, response-level retention follows, and the clinician audit validates a
post-retention sample rather than tuning the threshold.

We also aligned the stage name and audit reporting with this order. The manuscript now
summarizes the audit size, reviewer setup, and agreement reporting, with the detailed
criteria and rates shown below.


\responseheading{Revision Details}
\textbf{Pipeline order and stage name.}

\begin{displayquote}
\deleted{\textbf{Clinician-Audited Negative Response Construction:} we construct candidate negative responses through local rewriting under four controlled process-level failure types, apply automatic quality control, and validate a sampled subset with clinician audit.}

\added{\textbf{Controlled Negative Construction with Clinician Audit:} we construct candidate negative responses through local rewriting under four controlled process-level failure types, apply automatic quality control before risk-based retention, and validate a post-retention sample through clinician audit.}
\end{displayquote}

\textbf{Fixed threshold and no audit feedback into $\gamma$.}

\begin{displayquote}
\added{This threshold is fixed before the final clinician audit; the audit validates the retained pool rather than tuning $\gamma$.}
\end{displayquote}

\textbf{Appendix audit protocol and agreement reporting.}

\begin{displayquote}
\added{To validate the retained negatives used for preference training, we audit whether sampled negatives realize the intended controlled failure while preserving contextual coherence, clinical plausibility, model-likeliness/naturalness of the injected error, and the required answer policy. The audit uses 400 sampled responses, corresponding to approximately 1.2\% of the retained pool, with 100 items per error type. Five board-certified physicians in internal medicine and emergency medicine serve as reviewers, and 100 items receive double annotation.}

\deleted{The clinician audit includes three binary judgments---target-type match, single-error dominance, and answer-policy compliance---and two 1--5 ordinal judgments---clinical plausibility and faithfulness degradation.}\added{The expanded local rewrite naturalness audit includes binary judgments for target-type match, single-error dominance, answer-policy compliance, and strict accept-for-training, as well as 1--5 ordinal judgments for contextual coherence, model-likeliness/naturalness, clinical plausibility, and faithfulness degradation.} Table~\ref*{ms-tab:negative_quality} reports the \added{adjudicated} audit outcomes, while this appendix reports \deleted{agreement on the two ordinal judgments using weighted Cohen's kappa, together with exact and within-1 agreement on}\added{inter-rater agreement on} the 100 doubly annotated items. \added{The naturalness and coherence judgments are subjective, so we report Cohen's $\kappa$ for binary judgments and quadratic-weighted Cohen's $\kappa$, exact agreement, and within-1 agreement for ordinal judgments.}
\end{displayquote}

\newpage
\begin{commentbox}{R1.8 Similarity to Existing Work}
The overall modeling pipeline is highly similar to existing work. Fa-DPO claims its
core innovation is converting fine-grained faithfulness risk into preference signals,
but related ideas have been explored by HA-DPO (2024) and MedRAGChecker (2026).
\end{commentbox}

\responseheading{Response and Revision Summary}

We thank the reviewer for pointing out the relationship between Fa-DPO and two important related lines of work: HA-DPO and MedRAGChecker. We agree that our previous wording may have made the novelty sound broader than intended. In the revised manuscript, we therefore clarify the contribution as a more specific technical claim: Fa-DPO converts localized ARU-level evidence-use verification into preference signals for offline DPO in long-form medical reasoning.

\begin{enumerate}
\item \textbf{Difference from HA-DPO.} HA-DPO provides an important reference for hallucination-aware preference optimization. Fa-DPO differs from it in three main aspects:
\begin{enumerate}
\item \textbf{Target problem.} HA-DPO focuses on hallucination-aware preference learning, whereas Fa-DPO focuses on localized evidence-use errors in long-form medical reasoning. In our setting, the key issue is not only whether a response is globally hallucinated, but whether specific clinical reasoning units are supported by the provided evidence.

    \item \textbf{Supervision granularity.} HA-DPO mainly operates with response-level hallucination labels or preference signals. Fa-DPO instead decomposes long medical rationales into reasoning-function-labeled ARUs and estimates ARU-level evidence-use risk, so that localized unsupported reasoning spans can be identified.

    \item \textbf{Integration into DPO.} Fa-DPO does not simply train DPO on faithful versus unfaithful responses. ARU-level risk affects retained-negative selection, rejected-span token weighting, and risk-adaptive preference margins, giving stronger learning signals to local high-risk errors that would otherwise be diluted in sequence-level DPO.
\end{enumerate}

\item \textbf{Difference from MedRAGChecker.} MedRAGChecker provides an important reference for decomposing biomedical RAG outputs and verifying their evidence support. Fa-DPO differs from it in three main aspects:
\begin{enumerate}
    \item \textbf{Role in the pipeline.} MedRAGChecker mainly serves as a post-generation verification and diagnostic framework. Fa-DPO uses ARU-level verification before and during preference optimization, converting verification results into training signals for offline DPO.

    \item \textbf{Verification unit and purpose.} MedRAGChecker verifies generated medical claims to assess whether they are supported. Fa-DPO verifies reasoning-function-labeled ARUs, which are designed not only for support checking but also for localizing evidence-use risk within long medical rationales.

    \item \textbf{Use of verification scores.} In MedRAGChecker, verification scores are mainly used for evaluation and diagnosis. In Fa-DPO, these scores directly determine which negative samples are retained, which rejected spans receive stronger token-level penalties, and which preference pairs receive larger risk-adaptive margins.
\end{enumerate}

\end{enumerate}

Overall, the revised manuscript narrows the contribution from a broad faithfulness claim to a specific technical claim: converting localized ARU-level evidence-use risk into offline DPO signals for long-form medical reasoning. We revised the Abstract, Highlights, Introduction, Research Objectives, Method, and Related Work sections to state this distinction explicitly.



\responseheading{Revision Details}
\textbf{1. Abstract: narrowed the novelty claim from generic faithfulness-aware DPO to
localized evidence-use credit assignment.}
The original abstract described Fa-DPO broadly as converting fine-grained faithfulness
risk into preference signals. We revised this claim to specify the actual training
mechanism: localized evidence-use errors guide rejected-span weighting and
risk-adaptive preference margins.

\begin{displayquote}
\deleted{However, DPO's sequence-level objective will mask localized faithfulness errors in long, mostly correct clinical responses.}\added{However, DPO's sequence-level objective can dilute the suppressive signal from localized evidence-use errors in long, mostly correct clinical responses.}

\deleted{In this paper, we propose \textbf{Fa-DPO}, a \textbf{F}aithfulness-\textbf{A}ware \textbf{D}irect \textbf{P}reference \textbf{O}ptimization framework that converts fine-grained faithfulness risk into preference signals for evidence-grounded MQA.}\added{In this paper, we propose \textbf{Fa-DPO}, a \textbf{F}aithfulness-\textbf{A}ware \textbf{D}irect \textbf{P}reference \textbf{O}ptimization framework for evidence-grounded MQA.}

\added{Fa-DPO identifies localized evidence-use and reasoning--answer consistency failures in long clinical reasoning chains and uses them to strengthen learning signals for risky reasoning spans and high-risk responses.}

\deleted{We then decompose each response into atomic reasoning units, verify each unit through role-specific evidence routing and natural language inference (NLI)-based checking, and use the resulting risk scores to reweight high-risk spans and set risk-adaptive margins.}\added{Responses are decomposed into reasoning-function-labeled Atomic Reasoning Units (ARUs), verified via natural language inference (NLI), to guide risk-aware preference optimization.}
\end{displayquote}

\textbf{2. Highlights, Introduction, and Research Objectives: removed remaining broad novelty wording.}
We also revised the highlights, the main introduction claim, the contribution list, and
the research-objective statement
so that the manuscript consistently presents Fa-DPO as localized evidence-use credit
assignment and faithfulness-aware optimization using ARU-level verification scores
rather than generic conversion wording.

\begin{displayquote}
\deleted{Our Fa-DPO converts fine-grained faithfulness risk into preference signals for evidence-grounded medical question answering to suppress clinical hallucinations.}\added{We propose Fa-DPO, a faithfulness-aware direct preference optimization framework that uses local evidence-use risk to strengthen preference learning for long-context medical reasoning.}

\deleted{Our faithfulness-risk-guided strategy dynamically penalizes tokens in high-risk spans and enforces a risk-dependent margin to mitigate signal dilution in sequence-level DPO.}\added{We design a faithfulness-risk-guided optimization strategy that dynamically penalizes tokens in high-risk spans and enforces a risk-dependent margin to mitigate signal dilution in sequence-level DPO.}

\deleted{We measure faithfulness risk of atomic reasoning units through multi-source evidence routing and NLI-based verification.}\added{We introduce an ARU-level evidence-use risk estimation through multi-source evidence routing and NLI-based verification, enabling localized detection of unsupported reasoning.}

\deleted{In this paper, we propose \textbf{Fa-DPO}, a \textbf{F}aithfulness-\textbf{a}ware \textbf{DPO} framework for MQA that converts fine-grained faithfulness risk into preference signals.}
\added{In this paper, we propose \textbf{Fa-DPO}, a \textbf{F}aithfulness-\textbf{a}ware \textbf{DPO} framework that strengthens preference learning for evidence-grounded MQA by modeling local evidence-use risk.}

\deleted{We propose Fa-DPO, a faithfulness-aware DPO framework for evidence-grounded MQA, which converts fine-grained faithfulness risk into preference signals for evidence-grounded medical question answering to suppress clinical hallucinations.}\added{We propose Fa-DPO, a faithfulness-aware DPO framework for evidence-grounded MQA that improves faithful long-context clinical reasoning.}

\deleted{Convert ARU-level faithfulness risk into preference signals by retaining high-risk negatives, reweighting tokens in high-risk spans, and setting risk-adaptive preference margins.}\added{Use ARU-level verification scores for localized DPO credit assignment by retaining high-risk negatives, reweighting tokens in high-risk rejected spans, and setting risk-adaptive preference margins.}
\end{displayquote}

\textbf{3. Related Work on preference optimization: narrowed the gap against fine-grained and hallucination-aware DPO variants.}
We now state that prior variants move preference signals toward local units or
hallucination labels, whereas Fa-DPO uses reasoning-function-labeled medical units to
decide which rejected spans and responses receive stronger DPO penalties.

\begin{displayquote}
\added{Fine-grained, context-faithful, and hallucination-aware variants move preference signals toward tokens, steps, context conflicts, or response-level hallucination labels~\cite{gu2025maskdpo, zhou2025treg, yang2025sepo, xu2025fullstepdpo, bi2025contextdpo, zhao2024hadpo}.}

\deleted{The remaining gap is risk calibration: existing fine-grained methods improve process supervision or self-correction, but they do not calibrate local faithfulness risk inside offline DPO.}\added{These methods show that preference learning can target factual or contextual errors, but they do not use reasoning-function-labeled medical units to decide which rejected spans and responses should receive stronger DPO penalties.}

\deleted{Our method follows the offline preference-learning route, but differs from standard and fine-grained DPO by applying ARU-level evidence checking to retain rejected responses, weight spans, and calibrate the penalty assigned to each rejected response.}\added{Fa-DPO fills this gap by using ARU-level verification scores to retain high-risk negatives, weight rejected spans, and set risk-adaptive margins in offline DPO.}
\end{displayquote}

\textbf{4. Related Work on evidence verification: added a direct HA-DPO and MedRAGChecker comparison.}
We now acknowledge the shared decompose-then-verify principle, and then distinguish
Fa-DPO by the function of the unit: HA-DPO mainly uses response-level hallucination
preference labels, MedRAGChecker verifies generated claims after generation, whereas
Fa-DPO maps ARU-level verification scores to retained-negative selection, token-level
weights, and sample-level margins during DPO training.

\begin{displayquote}
\added{Self-checking, verifier-RAG, atomic factuality evaluation, and biomedical RAG checkers decompose outputs into claims and verify them against generated, retrieved, or external evidence~\cite{manakul2023selfcheckgpt, phasook2025vereafine, min-etal-2023-factscore, wei2024longform, song-etal-2024-veriscore, ji2026medragchecker}.}

\added{Existing verification signals are usually used for evaluation, filtering, consensus checking, or post-hoc diagnosis.}

\deleted{Our method differs from these evaluation and verification studies by converting fine-grained faithfulness measurements into training-time preference supervision, so that evidence-use failures become optimization signals rather than only diagnostic outputs.}\added{Fa-DPO adopts the decompose-then-verify principle but changes its role: ARUs encode clinical reasoning function, evidence route, local verification score, and token span, allowing verification results to become DPO weights and margins rather than only diagnostic labels.}

\added{Compared with HA-DPO, which primarily relies on response-level hallucination preference labels, and MedRAGChecker, which performs post-hoc medical RAG verification, Fa-DPO converts ARU-level verification outcomes into retained-negative selection, rejected-span token weights, and risk-adaptive DPO margins during training.}

\deleted{Taken together, these gaps motivate our method, which targets the unresolved bottleneck of faithful clinical reasoning in evidence-grounded MQA by converting fine-grained faithfulness risk into stable preference learning signals.}\added{Fa-DPO addresses this bottleneck by moving from evidence access to evidence use, from post-hoc faithfulness checking to training-time optimization, and from sequence-level preference learning to risk-adaptive span-aware preference learning.}
\end{displayquote}

\newpage
\begin{commentbox}{R1.9 Local Rewriting and Preference Data Generalization}
The local rewriting strategy may break clinical plausibility and create unnatural
negative samples. If negatives are manually injected rather than sampled from model
generations, DPO preferences may not generalize to actual model error patterns.
\end{commentbox}

\responseheading{Response and Revision Summary}
We agree. We revised the manuscript to describe locally rewritten negatives as
controlled counterfactual supervision, not samples from the full natural error
distribution of medical LLMs.

We also added a retained-negative naturalness audit, removed the construct-mismatched
Med-HALT pilot, and replaced it with a task-aligned natural-generation analysis. The
generalization claim is now narrower: the evidence supports transfer within
evidence-grounded MQA, while F1--F4 remains an incomplete taxonomy of broader medical
errors.


\responseheading{Revision Details}
\textbf{1. Training-data framing in the Introduction and Methodology.}
\begin{displayquote}
\deleted{To provide clinically meaningful negative responses for preference learning, we first simulate realistic clinical reasoning failures to construct negative responses.}

\added{To build preference data with clear local contrasts, we construct rejected responses containing controlled, clinically plausible failures in evidence use and reasoning--answer consistency.}
\end{displayquote}

\textbf{2. Controlled-counterfactual interpretation boundary in the negative-construction section.}
\begin{displayquote}
\added{These constructed negatives should be interpreted as controlled counterfactuals rather than as a complete model-error distribution.}

\added{Local rewriting gives us control over the dominant failure type and keeps the clinical case, answer style, and most surrounding reasoning fixed, which helps isolate localized evidence-use and reasoning-answer consistency failures for preference learning.}

\added{This preservation constraint improves controllability, but it does not by itself guarantee that the rejected responses follow the full distribution of naturally generated medical-LLM errors.}

\added{We therefore describe F1--F4 as controlled evidence-use and reasoning-consistency failures in evidence-grounded MQA rather than as an exhaustive taxonomy of broader medical hallucination phenomena.}

\deleted{Appendix~\ref*{ms-app:negative_protocol} gives the full prompt and filtering details.}

Appendix~\ref*{ms-app:negative_protocol} \added{gives the prompt structure and filtering details.}
\end{displayquote}

\textbf{3. Expanded clinician audit for local-rewrite naturalness and plausibility.}
\begin{displayquote}
\deleted{We conduct a clinician audit on 400 sampled responses with 100 items per error type.}

\added{The audit uses 400 sampled responses, corresponding to approximately 1.2\% of the retained pool, with 100 items per error type. Five board-certified physicians in internal medicine and emergency medicine serve as reviewers, and 100 items receive double annotation.}

\added{whether one local error remains dominant, whether the full rewritten chain remains contextually coherent, whether the injected error is plausible as an error that a medical LLM could naturally produce,} whether the response remains clinically plausible, whether it follows the answer policy\added{, and whether it is acceptable as a controlled counterfactual negative}.

\added{In the retained-pool audit, most sampled constructed negatives remain coherent, clinically plausible, and model-like after risk-based retention, with a strict accept-for-training rate of 65.0\% (95\% Wilson CI: 60.2--69.5\%).}

\added{We therefore treat the retained dataset as clinically enriched controlled supervision, not fully clinician-ranked RLHF data.}
\end{displayquote}

\textbf{4. Scope of external evaluation and removal of the Med-HALT pilot.}
\begin{displayquote}
\deleted{We further use a sampled Med-HALT reasoning pilot set as a supplementary external check of natural-error coverage.}

\added{To test transfer beyond controlled counterfactual rewrites, we evaluate the representative II-Medical-8B matched pair on 3,338 model-sampled cases from MedXpertQA-test, MedBullets-op4, MedBullets-op5, and MedQA-test. This analysis compares Standard DPO-Retained-34K and Fa-DPO-Retained-34K responses on the same held-out prompts without local rewriting. The labels are automatic rather than double-clinician-adjudicated and may co-occur, so this analysis supports a fixed-protocol consistency check rather than definitive clinical ground truth.}
\end{displayquote}

\begin{displayquote}
\deleted{Pilot Natural-Error Coverage Protocol on Med-HALT.}

\deleted{This appendix records the pilot protocol used for the Med-HALT study. The goal is not to compare Fa-DPO against Standard DPO under matched conditions, but to ask a narrower external-validity question: whether sampled errors from an external medical reasoning benchmark can be described using the controlled F1--F4 taxonomy introduced in Sec. 3.}

\deleted{Sampling Scope. The pilot uses a sample of 120 Med-HALT reasoning responses. The labeling unit is the complete response, consisting of the medical prompt, the generated reasoning chain, and the final answer. This pilot is analyzed separately from the paired MedCaseReasoning same-answer evaluations and is intended only as supplementary external evidence.}
\end{displayquote}

\textbf{5. Held-out natural-error annotation protocol.}
\begin{displayquote}
\added{The held-out natural-error analysis is conducted as a frozen post hoc evaluation over already generated model responses. Official answers are used only after generation to determine answer correctness and assign response-level error labels for analysis; they are not included in prompts, training data, candidate retention, hyperparameter tuning, or model selection. Each row contains one generated response with the prompt, answer options, gold option, predicted option, and visible reasoning.}

\added{The label definitions follow the evidence-grounded F1--F4 schema but are applied to natural outputs. F1 marks omission or substantive weakening of clinically important question evidence, F2 marks incorrect medical interpretation of evidence that the response uses, F3 marks overstatement beyond what the available facts support, and F4 marks a selected answer that is unsupported by or inconsistent with the preceding reasoning.}

\added{An unmapped error is a substantive medical reasoning or knowledge error not cleanly captured by F1--F4; unsupported medical claim and harmful recommendation are recorded as separate safety-relevant flags. The labels are automatic rather than double-clinician-adjudicated, and they are multi-label: F1--F4 may co-occur with one another and with unmapped, unsupported-claim, or harmful-recommendation flags.}
\end{displayquote}

\textbf{6. Discussion and Future Work: task-aligned generalization and bounded scope.}
\begin{displayquote}
\added{The evidence also shows bounded task-family generalization. The matched benchmark results cover multiple evidence-grounded MQA datasets and general-instruction, medical-adapted, and reasoning-enhanced backbones. The held-out natural-error analysis on 3,338 model-sampled MQA cases further shows fewer mapped F1--F4 errors and fewer unsupported medical claims for the representative II-Medical-8B pair. These findings indicate transfer beyond controlled rewrites within the same task family, while keeping the scope limited: F1--F4 is an evidence-use taxonomy for this setting, not a complete account of all medical hallucination or clinical reasoning errors.}

\deleted{Several limitations remain. The current evidence is restricted to 8B models and multiple-choice MQA, and the pipeline depends on offline negative construction, ARU parsing, retrieval, and NLI verification.}

\added{The evaluation remains bounded: experiments use 8B models, evidence-grounded multiple-choice MQA, and controlled negative construction with ARU parsing, retrieval, and NLI verification. Some process metrics are computed using the same ARU/NLI verification framework that provides training signals. Therefore, these metrics should not be interpreted as fully independent proof of faithfulness. Reverse entailment is likewise treated only as a possible coverage or specificity diagnostic, not as a reported result or training signal. We address these limitations by reporting answer accuracy and clinician audits as complementary evidence.}

\added{The added audits are mostly retrospective and partly automatic; they do not establish broad hallucination detection or deployment safety, and the retained negatives do not cover the full distribution of naturally generated medical-LLM errors.}

\deleted{Overall, the present results establish a stronger process-faithfulness case for preference optimization in MQA, while future work should extend the evaluation to broader model scales and more open-ended clinical settings, and further strengthen verifier-side and clinician-side validation.}

\added{Future work should test Fa-DPO in larger, open-ended clinical settings, combine controlled negatives with model-sampled, clinician-ranked preference data, and strengthen parser, verifier, and clinician validation. Fa-DPO improves answer accuracy and automatic reasoning-process metrics, with clinician audits providing additional support that unsupported and safety-relevant reasoning patterns are reduced.}
\end{displayquote}

\newpage
\begin{commentbox}{R1.10 Standard DPO Baseline Uses Risk-Filtered Pairs}
Standard DPO and Fa-DPO use the same retained preference pairs, but the negatives are
the risk-filtered subset of Fa-DPO rather than all 47,393 quality-controlled candidates.
This may artificially inflate Fa-DPO gains.
\end{commentbox}

\responseheading{Response and Revision Summary}
We agree that the original baseline wording did not clearly separate retained-pool
selection from the Fa-DPO objective. We therefore renamed the matched systems
\textbf{Standard DPO-Retained-34K} and \textbf{Fa-DPO-Retained-34K}, making clear
that both use the same 34,204 retained pairs and differ in the objective.

We also added representative full-QC, random-size-matched, retained-pair, SFT, and
reranking controls on II-Medical-8B. Among these added controls, Fa-DPO gives the
strongest overall result and keeps single-generation inference, while reranking
requires online ARU/NLI verification.


\responseheading{Revision Details}
\textbf{Training preference dataset: separate the full QC candidate pool from the
retained training pool.}
\begin{displayquote}
\added{Automatic construction and quality control yield 47,393 candidate pairs; ARU-score-based retention yields 34,204 final preference pairs from 12,411 source positives.}

\deleted{Standard DPO and Fa-DPO}\added{Standard DPO-Retained-34K and
Fa-DPO-Retained-34K} \added{use the same retained pair content, while Fa-DPO additionally uses $\mathbf{s}_l$ for token-level weighting and the sample-level margin.}
\end{displayquote}

\textbf{Baselines subsection: make clear what the matched comparison isolates.}
\begin{displayquote}
\deleted{We organize the comparison systems into three categories and use them in two
ways: controlled matched Fa-DPO vs.\ Standard DPO comparisons and a broader
released-baseline comparison.}\added{We use three baseline layers. Released-model baselines contextualize competitiveness across general instruction, medical-adapted, and reasoning-oriented medical models; the complete released-model list is reported in Table~\ref*{ms-tab:competitiveness_reference}. Matched retained-pair DPO controls provide the primary causal comparison, and representative faithfulness-oriented controls test simpler explanations for the observed gain.}

For each starting point, we train one \deleted{Standard DPO variant and one
Fa-DPO variant}\added{Standard DPO-Retained-34K variant and one
Fa-DPO-Retained-34K variant}.

\deleted{Standard DPO uses these same retained preference pairs as Fa-DPO}\added{Within each pair, the two variants share the same 34,204 retained preference pairs, preprocessing, train/dev split, training budget, verifier pipeline, prompt template, decoding setup, and answer extraction rule; the intended difference is only the preference objective.}
\end{displayquote}

\textbf{Baselines subsection: add full-candidate and size-matched DPO controls.}
\begin{displayquote}
\added{On the representative II-Medical-8B backbone, we additionally compare against SFT-faithful, Standard DPO-QC-All, Standard DPO-Random-34K, Standard DPO-Retained-34K, and verifier reranking with $K=4$ or $K=8$. These controls separate positive-only imitation, full-candidate DPO, sample-count effects, retained-pool effects, and inference-time verification.}
\end{displayquote}

\textbf{Main results: relabel the paired comparison throughout tables, captions, and
interpretation.}
\begin{displayquote}
We next test whether the automatic process-metric gains coincide with better final-answer
accuracy. For this evaluation, we compare \deleted{Fa-DPO with Standard
DPO}\added{Fa-DPO-Retained-34K with Standard DPO-Retained-34K} under otherwise
identical settings across six paired 8B base models.

\deleted{Fa-DPO improves macro-average accuracy over Standard DPO}\added{Fa-DPO-Retained-34K improves macro-average accuracy for all six base models, with observed gains of 1.2--1.6 points; the benchmark-stratified paired bootstrap intervals reported in the last column exclude zero for every backbone.}
\end{displayquote}

\textbf{Faithfulness-oriented controls: report the strongest added-control result.}
\begin{displayquote}
\added{Among these added controls on the representative II-Medical-8B backbone, Fa-DPO gives the strongest overall result. Positive-only SFT improves over the base model, but remains below Fa-DPO by 1.0 macro-average point and leaves a 7.8-point gap in Unsupported Medical Reasoning Rate. The full-candidate, random-pool, and retained-pair DPO controls also remain below Fa-DPO; under identical retained-pair content, Fa-DPO adds +1.2 average accuracy and improves all four process metrics over Standard DPO-Retained-34K. Verifier reranking narrows the gap, especially at $K=8$, but requires multiple generations and online ARU/NLI verification; Fa-DPO keeps single-generation inference.}
\end{displayquote}

\textbf{Appendix training details: define the representative control suite
operationally.}
\begin{displayquote}
\added{The representative baseline controls in Table~\ref*{ms-tab:faithfulness_baseline_results} use the II-Medical-8B~\cite{intelligentinternet2025iimedical8b} backbone and the
same benchmark splits, prompts, answer-extraction rule, and process-metric pipeline
as the matched II-Medical-8B comparison. SFT-faithful trains on the exact
preferred-response side ($y_w$) used in the retained preference pairs, deduplicated
by prompt and target response; it uses the same tokenizer, optimizer family, LoRA
setup, sequence limits, and one-epoch supervised schedule, but it is not an
update-matched DPO run because it is a positive-only supervised control. Standard
DPO-QC-All uses all 47,393 quality-controlled candidate pairs, Standard
DPO-Random-34K uses a random 34,204-pair subset of that QC pool, and Standard
DPO-Retained-34K uses the 34,204 ARU-risk-retained pairs. Fa-DPO-Retained-34K uses
the same retained pairs as Standard DPO-Retained-34K and adds only the token-level
risk weighting and response-level margin. Together, these controls define the
representative II-Medical-8B attribution check reported in
Table~\ref*{ms-tab:faithfulness_baseline_results}.}
\end{displayquote}

\newpage
\begin{commentbox}{R1.11 Same-Answer Subset Selection Bias}
Process-faithfulness results are reported only on the subset where both models predict
the correct final answer. This may exclude cases where Fa-DPO improves or worsens
answer correctness and introduces conditional selection bias.
\end{commentbox}

\responseheading{Response and Revision Summary}
We agree that the same-answer subset is conditional and should not be presented as an
unconditional estimate of overall response quality. We revised the manuscript to
describe it as a controlled mechanism check that fixes final-answer correctness while
testing reasoning-process indicators.

To address the selection-bias concern, we added full-test and correctness-stratified
MedCaseReasoning checks using the same answer extraction and ARU/NLI pipeline. These
diagnostics move in the same direction but are interpreted conservatively, alongside
the same-answer clinician validation, rather than as a pooled significance claim.


\responseheading{Revision Details}
\textbf{Process-results subsection: clarify same-answer as a conditional
mechanism analysis.}
\begin{displayquote}
This same-answer control isolates reasoning quality rather than answer
accuracy\added{, but it is a conditional mechanism analysis rather than an
unconditional estimate of full-test response quality}.

\deleted{The main text reports paired deltas only, while detailed subset-construction rules and the representative detailed table are deferred to the appendix.}

\added{The same-answer table reports paired deltas with paired uncertainty; full-test and correctness-stratified checks are reported below and in Appendix~\ref*{ms-app:process_metrics}.}
\end{displayquote}

\textbf{Main results: add the full-test and correctness-stratified check.}
\begin{displayquote}
\added{The full MedCaseReasoning test split shows the same direction with smaller effects: average accuracy increases by 3.8 points, Support Rate increases by 3.9 points, and the three error-oriented process metrics decrease by 2.2--6.3 points. This check is only diagnostic because the same 160 questions are reused across backbones and the per-backbone McNemar tests are not individually significant ($p=0.296$--$0.522$). The representative Standard-DPO-only-correct stratum also shows worse process metrics for Fa-DPO, so we treat the same-answer results as controlled mechanism evidence, not as an unconditional estimate of full-test response quality.}
\end{displayquote}

\noindent The newly added full-test summary table is:

\begin{table}[htbp]
\centering
\small
\caption{R1.11 full-test selection-bias summary. Metrics are computed on the complete MedCaseReasoning test split without conditioning on both systems being correct. Accuracy significance checks use per-backbone McNemar tests; the pooled 960-entry total is diagnostic only because the same questions are reused across backbones. This table is separate from the four-benchmark downstream macro-average accuracy table discussed in R1.13.}
\label{tab:response_r111_full_test_summary}
\begin{adjustbox}{max width=\textwidth}
\begingroup\color{blue}
\begin{tabular}{@{}lccc@{}}
\toprule
\textbf{Full-test metric} & \textbf{Standard DPO-Retained-34K} & \textbf{Fa-DPO-Retained-34K} & \textbf{$\Delta$ / test} \\
\midrule
Accuracy & 61.3 & 65.0 & +3.8; per-backbone McNemar $p=0.296$--$0.522$ \\
Support Rate & 73.7 & 77.6 & +3.9 [ +1.3, +6.9 ] \\
Unsupported Medical Reasoning Rate & 36.9 & 30.7 & -6.3 [ -10.4, -2.1 ] \\
Contradiction Rate & 11.2 & 9.0 & -2.2 [ -4.3, -0.5 ] \\
Reasoning-Answer Inconsistency Rate & 30.1 & 25.5 & -4.6 [ -8.4, -0.9 ] \\
Correct-and-Process-Pass Rate & 22.3 & 29.6 & +7.3 [ +1.9, +12.8 ] \\
\bottomrule
\end{tabular}
\endgroup
\end{adjustbox}
\end{table}

\textbf{Appendix protocol: define the new proxy rate and stratified checks.}
\begin{displayquote}
\added{For the full-test selection-bias check, we additionally report \textit{Correct-and-Process-Pass Rate}: the percentage of questions for which the final answer is correct and the automatic ARU/NLI process trace has no contradiction ARU, no unsupported warrant/differentiation ARU, and no inconsistent claim ARU. This is a strict automatic diagnostic pass rate rather than clinician-adjudicated ground truth and is used only as a selection-bias diagnostic.}

\added{To test whether same-answer conditioning changes the interpretation, we also compute the same automatic ARU/NLI-based process metrics on the full MedCaseReasoning test split. Table~\ref*{ms-tab:full_set_process_metrics} reports full-test metrics without conditioning on both systems being correct. Table~\ref*{ms-tab:correctness_strata_counts} reports the paired correctness strata used to diagnose which cases the same-answer subset excludes. Table~\ref*{ms-tab:iimedical_correctness_stratified_process} reports the representative II-Medical-8B process metrics inside each correctness stratum.}
\end{displayquote}

\newpage
\begin{commentbox}{R1.12 Clinician Audit Scope}
Section 6.6.1 audits only 400 negative samples out of 34,204, and does not audit
positive samples or model-generated responses. Positive samples may be imperfect,
creating systematic bias.
\end{commentbox}

\responseheading{Response and Revision Summary}
We thank the reviewer for raising this point. To separate preferred-side reliability from pairwise preference strength, we added both preferred-response and pair-direction audits.

These audits show that the retained preferred responses are reliable training targets, while the constructed pair directions contain measured directional uncertainty under a strict non-tie criterion. We therefore revised the manuscript to describe the retained data as controlled, risk-enriched supervision rather than noise-free clinician-ranked RLHF labels.

We also added a w/o-F1 sensitivity analysis, because F1 had the lowest pair-direction confidence among the four failure types. The main comparison remains consistent after excluding this type, suggesting that the observed gains are not driven only by the lowest-confidence subset.


\responseheading{Revision Details}
\textbf{Main-text validation summary: add the DPO-label validation step.}
\begin{displayquote}
\deleted{We next assess whether the evidence used in the preceding analyses is reliable.}
\added{This section validates the process evidence used in the preceding analyses. It asks three linked questions: whether automatic same-answer improvements align with clinician judgments and data-audit evidence; whether the ARU parser, routed NLI verifier, and token masks are reliable enough for proxy process measurement; and whether the same signal remains directionally useful on natural outputs, answer-changing diagnostics, and retrospective safety checks. These experiments support bounded use of the process metrics, not prospective clinical validation.}

\added{The preference-label audit further shows that preferred responses are reliable, while pair directions remain noisy. Table~\ref*{ms-tab:clinician_preference_validation_summary} summarizes these checks. We therefore treat the retained dataset as clinically enriched controlled supervision, not fully clinician-ranked RLHF data.}
\end{displayquote}

\textbf{New subsection: validate preferred responses and preference-pair direction.}
\begin{displayquote}
\added{To validate the preference labels, we separately check the preferred side and the pair direction. The negative-response audit above does not by itself validate DPO labels, because DPO also assumes that the preferred response $y_w$ is a reliable alignment target and that the pairwise direction $y_w \succ y_l$ is clinically meaningful. The preferred responses audited here are the same MedCaseReasoning source positive responses paired with retained rejected responses, not outputs newly generated by the controlled-rewrite model. We therefore conduct two blinded clinician audits on retained training items: a 200-item preferred-response audit and a 200-item pair-level audit, each balanced across the paired F1--F4 rejected-response types with 50 items per type.}
\end{displayquote}

\textbf{Preference-label audit table: report positive-label reliability and pairwise label noise.}
\begin{table}[htbp]
\centering
\small
\caption{\added{Clinician validation of preferred-response reliability and pairwise preference-label validity. The two retained-pair samples are balanced across F1--F4 rejected-response types and are targeted label-reliability audits rather than exhaustive quality control over all 34,204 retained pairs. For the preferred-response row, major issues are rare and the label distribution is highly imbalanced, so exact agreement and the adjudicated reliable-target rate are the interpretable quantities.}}
\label{tab:response_r112_preference_label_audit}
\begin{adjustbox}{max width=\textwidth}
\begin{tabular}{lccccc}
\toprule
\textbf{Audit target} & \textbf{$n$} & \textbf{Primary validity rate} & \textbf{95\% CI} & \textbf{Secondary findings} & \textbf{Agreement} \\
\midrule
\added{Preferred responses} & \added{200} & \added{Reliable target: 98.5\%} & \added{[95.7, 99.5]\%} & \added{Major issue: 1.5\%; final answer correct: 98.5\%; contradiction: 0.0\%} & \added{Exact 82.0\%; $\kappa$ not informative under rare-error imbalance ($-0.018$)} \\
\added{Preference pairs} & \added{200} & \added{Direction valid: 76.0\%} & \added{[69.6, 81.4]\%} & \added{DPO-suitable: 86.0\%; direct reversal: 10.0\%; tie: 14.0\%} & \added{Exact 79.0\%; $\kappa=0.662$} \\
\bottomrule
\end{tabular}
\end{adjustbox}
\end{table}

\begin{displayquote}
\added{The preferred-response side is reliable in this stratified audit: 197/200 preferred
responses are judged acceptable or minor-issue-but-usable, with a 1.5\%
major-issue rate and no adjudicated contradictions. For inter-rater agreement on
this rare-error judgment, exact agreement is 82.0\%; Cohen's $\kappa$ is not
informative because major preferred-response issues are rare and the label
distribution is highly imbalanced. We therefore treat the adjudicated
reliable-target rate of 98.5\% as the primary evidence for preferred-response
quality, rather than treating $\kappa$ as the primary quality measure for this
imbalanced binary audit. Pair-level labels are noisier.
Clinicians validate the constructed direction in 152/200 pairs and judge 172/200
pairs suitable for DPO training, but 20/200 pairs are direct reversals and 28/200
are ties with no clear preference. The two rates answer different questions:
direction validity asks whether the constructed label $y_w \succ y_l$ is
clinically correct, whereas DPO suitability asks whether the two responses
contain a usable preference contrast after clinician review. Thus, a direct
reversal may still reveal a clear contrast but is label noise for the submitted
constructed direction; ties are the least suitable for pairwise DPO because no
clear preference exists. Among non-tie cases, the constructed direction agrees
with the adjudicated clinician direction in 88.4\% of cases. Taken together,
these results support the preferred side as a usable target while showing that
the retained pairs remain noisy controlled supervision, not clinician-ranked
gold-standard RLHF preferences. The audit also exposes a substantive limitation
for DPO training: direct reversals can push the pairwise objective in the wrong
direction, and ties provide weak or no pairwise preference signal. Because
reversal and tie cases are not confined to F1, the F1-exclusion analysis below
should be read as a lowest-confidence-type sensitivity check, not as a complete
label-denoising experiment.}
\end{displayquote}

\textbf{Sensitivity analysis: lowest-confidence type check, not full label denoising.}

\begin{table}[htbp]
\centering
\small
\caption{\added{F1-exclusion sensitivity analysis for lower-confidence preference labels. Process metrics are automatic ARU/NLI-based percentages on the MedCaseReasoning same-answer hard subset, and Avg. Acc. is the macro-average over MedBullets-op4, MedBullets-op5, MedXpertQA, and MedQA.}}
\label{tab:response_r112_wo_f1_marked}
\begin{adjustbox}{max width=\textwidth}
\begin{tabular}{lccccc}
\toprule
\textbf{Method} & \textbf{Support $\uparrow$} & \textbf{Contradiction $\downarrow$} & \textbf{Unsupported $\downarrow$} & \textbf{Reason.-Ans. Inconsistency $\downarrow$} & \textbf{Avg. Acc. $\uparrow$} \\
\midrule
Standard DPO-Retained-34K & 77.3 & 9.4 & 31.6 & 26.1 & 62.0 \\
Fa-DPO-Retained-34K & 84.0 & 5.2 & 20.9 & 18.5 & 63.2 \\
\added{Fa-DPO-w/o-F1} & \added{83.2} & \added{5.5} & \added{21.8} & \added{19.2} & \added{63.0} \\
\bottomrule
\end{tabular}
\end{adjustbox}
\end{table}

\begin{displayquote}
\added{Removing F1 slightly weakens the full Fa-DPO result, but the main conclusion remains
intact. Compared with Standard DPO-Retained-34K, Fa-DPO-w/o-F1 still improves
Support Rate by 5.9 points, reduces Contradiction Rate by 3.9 points, reduces
Unsupported Medical Reasoning Rate by 9.8 points, reduces Reasoning-Answer Inconsistency
Rate by 6.9 points, and improves average accuracy by 1.0 point.}

\added{We selected F1 as a conservative lowest-quality sensitivity check: in the retained-pool audit, F1 has the lowest strict accept rate, and internal type-wise checks also identified omission/down-weighting edits as the least stable controlled rewrite category. Because the reported preference-label table does not provide type-wise pair-direction validity, this experiment should not be read as a full type-wise label-validity comparison. F4 is kept because it directly defines the reasoning-answer consistency target; excluding it would answer a target-coverage ablation question rather than a lower-confidence-data sensitivity question. This sensitivity should therefore not be interpreted as removing all pair-level reversal or tie noise from the retained training set. A stronger label-noise test would require a clinician-denoised retained pool in which confirmed reversals are corrected and ties are removed before retraining. This remains a residual limitation: future work should add clinician-ranked natural reasoning-answer inconsistency pairs and denoised preference-pair training.}
\end{displayquote}

\textbf{Evidence-validation section and conclusion: state the remaining limitation.}
\begin{displayquote}
\added{We therefore treat the retained dataset as clinically enriched controlled supervision, not fully clinician-ranked RLHF data. Detailed audit tables are reported in Appendices~\ref*{ms-app:negative_protocol}, \ref*{ms-app:preference_label_details}, and~\ref*{ms-app:same_answer_clinician_details}.}

\added{The added audits are mostly retrospective and partly automatic; they do not establish broad hallucination detection or deployment safety, and the retained negatives do not cover the full distribution of naturally generated medical-LLM errors.}
\end{displayquote}

\newpage
\begin{commentbox}{R1.13 Statistical Significance}
The reported 1.2--1.6 percentage-point gains are a range, not confidence intervals.
Section 6.2.1 does not report subset size. Improvements across six backbones should be
tested with paired t-tests or Wilcoxon tests.
\end{commentbox}

\responseheading{Response and Revision Summary}
We agree that the statistical presentation needed clearer paired uncertainty. We
revised the manuscript to state that the 1.2--1.6 point accuracy range is an observed
cross-backbone range, not a confidence interval.

We also added paired cluster-bootstrap intervals for same-answer process metrics and
benchmark-stratified paired-bootstrap intervals for downstream macro-average
accuracy. The Wilcoxon test over six backbone deltas is reported only as a
supplementary direction summary because the six evaluations are not fully independent.


\responseheading{Revision Details}
\textbf{Sec.~6.1: add paired statistical-analysis rules.}

\begin{displayquote}
\added{All comparisons between Standard DPO-Retained-34K and Fa-DPO-Retained-34K are paired because both systems are evaluated on the same questions under matched decoding, answer extraction, and verifier-side processing. For same-answer process metrics, each retained case contributes one response-level metric per system. We report case-weighted means and paired deltas, with 95\% confidence intervals from 10,000 paired cluster-bootstrap replicates over cases. Each bootstrap sample keeps both system responses and all nested ARUs from the sampled cases together; ARU counts are reported as measurement denominators, not as independent resampling units.}

\added{For downstream accuracy, macro-average uncertainty is estimated with benchmark-stratified paired bootstrap confidence intervals. Each replicate resamples paired test questions within each benchmark, recomputes the four benchmark accuracies, and then recomputes the simple four-benchmark macro-average. Benchmark-level paired tests, Holm-adjusted p-value screening, and the six-backbone Wilcoxon direction summary are reported as supplementary diagnostics; the main text emphasizes paired deltas and confidence intervals, and the 1.2--1.6 point range denotes the observed cross-backbone range rather than an uncertainty interval.}

\added{For representative ablation analyses, we use the same paired question-level bootstrap protocol for macro-average drops and apply Holm adjustment within the leave-one-out family. Detailed drop-difference checks for closely spaced components are reported with the appendix ablation analyses.}

\added{Together, these statistical choices keep the main-text uncertainty reporting focused on paired effect estimates while leaving detailed p-value screening to the supplementary analyses.}
\end{displayquote}

\textbf{Sec.~6.2.1 and Table~2: add subset sizes, ARU denominators, and paired CIs.}

\begin{displayquote}
\deleted{The main text reports paired deltas only, while detailed subset-construction rules and the representative detailed table are deferred to Appendix~D.}\added{The same-answer table reports paired deltas with paired uncertainty; full-test and correctness-stratified checks are reported below and in Appendix~\ref*{ms-app:process_metrics}.}

\deleted{Here, $n$ denotes the retained subset size for each pair.}\added{Here, $n$ denotes retained same-answer cases, responses denotes $2n$, and ARUs $N$ is the total ARU measurement count, not the resampling unit. Supplementary multiplicity-adjusted p-value screening is reported in Appendix~\ref*{ms-app:process_metrics}.}

\added{Table~\ref*{ms-tab:paired_process_summary} shows consistent process-metric improvements across all six matched pairs. On same-answer subsets of 68--100 cases per pair, Support Rate increases by an average of 5.1 points, while Unsupported Medical Reasoning Rate, Contradiction Rate, and Reasoning-Answer Inconsistency Rate decrease by 7.4, 3.0, and 5.6 points, respectively; all paired cluster-bootstrap intervals exclude zero.}
\end{displayquote}

\textbf{Sec.~6.2.2: distinguish observed accuracy range from uncertainty.}

\begin{displayquote}
Table~3 shows that the \deleted{process-faithfulness}\added{process-metric} gains \deleted{translate into}\added{are accompanied by higher answer accuracy}. \deleted{Fa-DPO improves macro-average accuracy over Standard DPO}\added{Fa-DPO-Retained-34K improves macro-average accuracy for all six base models, with observed gains of 1.2--1.6 points; the benchmark-stratified paired bootstrap intervals reported in the last column exclude zero for every backbone.}

\added{Benchmark-level evidence is heterogeneous rather than uniform: 13 of 24 McNemar tests have raw $p<0.05$, and 5 remain Holm-significant. Paired bootstrap intervals exclude zero in 16 of 24 benchmark cells, with 6 retained after Holm screening, so the positive benchmark cells should be read as a directional pattern rather than as uniform per-benchmark significance.}
\end{displayquote}

\noindent For reviewer verification, we also summarize the final paired statistical outputs in the following response-letter table:

\begin{table}[htbp]
\centering
\scriptsize
\caption{R1.13 paired statistical-significance summary. Same-answer process metrics use paired cluster-bootstrap confidence intervals over cases. Downstream macro-average accuracy uses benchmark-stratified paired bootstrap confidence intervals over test questions. The six-backbone Wilcoxon result is reported only as a supplementary direction summary.}
\label{tab:response_r113_statistical_significance}
\textit{Panel A. Same-answer automatic process-metric deltas with paired cluster-bootstrap uncertainty.}\\[0.2em]
\begin{adjustbox}{max width=\textwidth}
\begin{tabular}{lrrrllll}
\toprule
\textbf{Backbone} & \textbf{$n$} & \textbf{Resp.} & \textbf{ARUs $N$} & \textbf{Support $\Delta$ [95\% CI]} & \textbf{Contrad. $\Delta$ [95\% CI]} & \textbf{Unsupported $\Delta$ [95\% CI]} & \textbf{Claim Inc. $\Delta$ [95\% CI]} \\
\midrule
Llama-3.1-8B-Instruct & 68 & 136 & 1984 & +3.8 [+0.9, +6.8] & -2.1 [-4.5, -0.2] & -5.3 [-9.7, -1.8] & -4.2 [-8.1, -0.8] \\
Qwen3-8B & 74 & 148 & 2196 & +4.1 [+1.2, +7.3] & -2.4 [-4.9, -0.4] & -5.9 [-10.5, -2.0] & -4.6 [-8.6, -1.1] \\
Llama-3.1-8B-Instruct (w m23k SFT) & 82 & 164 & 2528 & +4.9 [+1.8, +8.0] & -2.8 [-5.2, -0.6] & -6.8 [-11.5, -2.6] & -5.3 [-9.6, -1.6] \\
Qwen3-8B (w m23k SFT) & 88 & 176 & 2744 & +5.2 [+2.0, +8.6] & -3.0 [-5.6, -0.8] & -7.3 [-12.3, -2.9] & -5.8 [-10.2, -2.0] \\
UltraMedical-3.1-8B & 92 & 184 & 2968 & +5.6 [+2.4, +9.4] & -3.4 [-6.1, -1.0] & -8.1 [-13.6, -3.2] & -6.1 [-11.0, -2.1] \\
II-Medical-8B & 100 & 200 & 3342 & +6.7 [+3.3, +11.0] & -4.2 [-7.3, -1.7] & -10.7 [-17.4, -4.8] & -7.6 [-13.7, -1.2] \\
\bottomrule
\end{tabular}
\end{adjustbox}

\vspace{0.45em}
\textit{Panel B. Downstream macro-average accuracy differences with benchmark-stratified paired uncertainty.}\\[0.2em]
\begin{adjustbox}{max width=\textwidth}
\begin{tabular}{lrrrrr}
\toprule
\textbf{Backbone} & \textbf{Std Macro Avg.} & \textbf{Fa-DPO Macro Avg.} & \textbf{$\Delta$ pp} & \textbf{95\% CI} & \textbf{Bootstrap $p$} \\
\midrule
Llama-3.1-8B-Instruct & 44.8 & 46.3 & +1.5 & [+0.6, +2.4] & 0.004 \\
Qwen3-8B & 46.9 & 48.4 & +1.5 & [+0.6, +2.5] & 0.005 \\
Llama-3.1-8B-Instruct (w m23k SFT) & 51.0 & 52.6 & +1.6 & [+0.7, +2.7] & 0.003 \\
Qwen3-8B (w m23k SFT) & 52.0 & 53.5 & +1.5 & [+0.6, +2.6] & 0.005 \\
UltraMedical-3.1-8B & 52.7 & 53.9 & +1.2 & [+0.4, +2.2] & 0.012 \\
II-Medical-8B & 62.0 & 63.2 & +1.2 & [+0.3, +2.2] & 0.018 \\
\midrule
\multicolumn{6}{l}{Supplementary exact Wilcoxon direction summary over six macro-average deltas: $V=21$, two-sided $p=0.031$.} \\
\bottomrule
\end{tabular}
\end{adjustbox}
\end{table}

In Panel B, the bootstrap p-values are two-sided empirical tail probabilities
from the paired bootstrap distribution around zero; the 95\% confidence intervals
remain the primary uncertainty summary. Supplementary appendix screening reports
the largest Holm-adjusted process-metric p-value for each backbone
(0.041, 0.034, 0.026, 0.021, 0.018, and 0.012). At the individual benchmark
level, 13 of 24 paired McNemar tests have raw $p<0.05$, and 5 remain significant
after Holm correction. Paired bootstrap intervals exclude zero in 16 of 24
benchmark cells, with 6 retained after Holm screening. The benchmark-level
pattern is directionally consistent but not uniformly significant after
multiplicity correction; MedXpertQA deltas are consistently positive but have
wider intervals on this harder benchmark.

\textbf{Appendix process-metric Holm-adjusted screening table.}
\begin{displayquote}
\begin{center}
\footnotesize
\captionof{table}{\added{Supplementary Holm-adjusted screening for the same-answer process-metric table. The main text reports 95\% confidence intervals; this appendix table reports the largest Holm-adjusted p-value among the four process-metric paired tests for each backbone after correction over the 24 backbone-metric tests.}}
\begingroup\color{blue}
\begin{tabular}{lc}
\toprule
\textbf{Backbone} & \textbf{Max Holm $p$} \\
\midrule
Llama-3.1-8B-Instruct & 0.041 \\
Qwen3-8B & 0.034 \\
Llama-3.1-8B-Instruct w m23k SFT & 0.026 \\
Qwen3-8B w m23k SFT & 0.021 \\
UltraMedical-3.1-8B & 0.018 \\
II-Medical-8B & 0.012 \\
\bottomrule
\end{tabular}
\endgroup
\end{center}
\end{displayquote}

\newpage
\begin{commentbox}{R1.14 Ablation Interaction Effects}
The ablation study uses one-by-one removal but does not consider inter-module coupling.
It lacks progressive addition experiments such as only negative filtering, only token
weighting, and only margin, and lacks statistical significance tests.
\end{commentbox}

\responseheading{Response and Revision Summary}
We agree. We expanded Sec.~6.3 beyond one-by-one ablations by adding leave-one-out,
progressive objective, and granularity--routing interaction analyses, with
ablation-specific statistical checks.

We interpret these as representative attribution analyses on II-Medical-8B rather
than six-backbone significance claims. The hyperparameter-sensitivity check is now
folded into the ablation section, with full details left to the appendix.


\responseheading{Revision Details}
\textbf{1. Statistical Analysis: add ablation-specific testing rule.}
\begin{displayquote}
\added{For representative ablation analyses, we use the same paired question-level bootstrap protocol for macro-average drops and apply Holm adjustment within the leave-one-out family. Detailed drop-difference checks for closely spaced components are reported with the appendix ablation analyses.}
\end{displayquote}

\textbf{2. Sec.~6.3 setup: state the ablation structure.}

\begin{displayquote}
\deleted{we remove one design choice at a time while keeping the rest of the training pipeline fixed}

\deleted{We next test which components are responsible for the paired gains in Table~\ref*{ms-tab:matched_main_results}. On the II-Medical-8B~\cite{intelligentinternet2025iimedical8b} backbone, we remove one design choice at a time while keeping the rest of the training pipeline fixed, so the comparison isolates the contribution of each module. All variants reuse the same verifier-side retrieval pipeline, NLI verifier, train/dev split, and optimization budget as the full model. We group the variants by optimization objective, data construction, and verification mechanism:}

\added{After the main-result block establishes that Fa-DPO improves matched performance, this subsection asks which design choices carry that effect. On the representative II-Medical-8B~\cite{intelligentinternet2025iimedical8b} backbone, the main text reports the compact leave-one-out ablation table, while Appendix~\ref*{ms-app:ablation_details} reports the progressive negative-filtering/objective matrix and the reasoning-unit granularity/routing matrix. These ablations are single-backbone attribution analyses and should not be read as evidence that every drop has the same significance pattern across all six models. All variants reuse the same verifier-side retrieval pipeline, NLI verifier, train/dev split, training budget, and evaluation settings unless explicitly stated.}
\end{displayquote}

\textbf{3. Leave-one-out ablation table: update the table with
uncertainty and Holm-adjusted significance.}

\begin{table}[htbp]
\centering
\scriptsize
\caption{\deleted{Ablation analysis on component effectiveness within Fa-DPO on the II-Medical-8B~\cite{intelligentinternet2025iimedical8b} backbone. We report Accuracy (\%) across the same four evaluation settings used in Table~\ref*{ms-tab:matched_main_results}.}\added{Leave-one-out ablation analysis on the II-Medical-8B~\cite{intelligentinternet2025iimedical8b} backbone. We report Accuracy (\%) on MedBullets-op4, MedBullets-op5, MedXpertQA, and MedQA, their macro-average, the paired macro-average delta relative to the full model, paired-bootstrap 95\% confidence intervals, and Holm-adjusted $p$ values over the six leave-one-out comparisons. The full model is the reference row; confidence intervals and Holm-adjusted $p$ values are reported only for leave-one-out comparisons against it.}}
\label{tab:response_r114_leave_one_out}
\label{tab:ablation}
\begin{adjustbox}{max width=\textwidth}
\begin{tabular}{lcccccccc}
\toprule
\textbf{Method} & \textbf{op4} & \textbf{op5} & \textbf{MedXpertQA} & \textbf{MedQA} & \textbf{Avg.} & \added{\textbf{$\Delta$}} & \added{\textbf{95\% CI}} & \added{\textbf{Holm $p$}} \\
\midrule
\textbf{Fa-DPO (Full)} & 75.0 & 70.2 & 24.1 & 83.4 & 63.2 & \added{0.0} & \added{Reference} & \added{Reference} \\
\midrule
\multicolumn{9}{l}{\textit{Optimization Objectives}} \\
\quad \textit{w/o Sample-Level Margin} & 73.6 & 69.4 & 23.1 & 82.8 & 62.2 & \added{-1.0} & \added{[-1.8, -0.2]} & \added{0.041} \\
\quad \textit{w/o Token-Level Dynamic Reweighting} & 73.2 & 68.7 & 22.5 & 82.2 & 61.7 & \added{-1.5} & \added{[-2.5, -0.5]} & \added{0.018} \\
\midrule
\multicolumn{9}{l}{\textit{Data Construction}} \\
\quad \textit{w/o Clinical Negative Responses} & 71.1 & 67.0 & 21.4 & 82.3 & 60.5 & \added{-2.7} & \added{[-4.1, -1.4]} & \added{0.004} \\
\midrule
\multicolumn{9}{l}{\textit{Verification Mechanism}} \\
\quad \textit{w/o Function-Specific Routing} & 72.8 & 68.3 & 22.8 & 82.4 & 61.6 & \added{-1.6} & \added{[-2.7, -0.5]} & \added{0.017} \\
\quad \textit{w/o ARU-Level Granularity} & 72.1 & 67.9 & 22.1 & 82.2 & 61.1 & \added{-2.1} & \added{[-3.5, -0.8]} & \added{0.009} \\
\quad \textit{w/o Retrieval-Based Verification} & 71.4 & 66.8 & 21.9 & 82.6 & 60.7 & \added{-2.5} & \added{[-3.9, -1.1]} & \added{0.006} \\
\bottomrule
\end{tabular}
\end{adjustbox}
\end{table}

\begin{displayquote}
\added{\textbf{w/o Clinical Negative Responses:} Replaces the clinical error-driven
generator (Module 1) with \textit{generic negative sampling}. The source prompts,
preferred responses, automatic QC, retention/scoring pipeline, training budget,
and evaluation protocol are kept fixed, but the rejected candidates are generic
incorrect-answer samples rather than F1--F4 clinical-error rewrites.}
\end{displayquote}

\begin{displayquote}
\added{Table~\ref*{ms-tab:ablation} shows that every tested removal lowers
macro-average accuracy on the representative II-Medical-8B backbone, and all six
drops remain significant after Holm correction within this ablation family. The
largest drops come from removing clinical negatives (-2.7), retrieval-based
verification (-2.5), and ARU-level granularity (-2.1). Because several
drop-difference intervals overlap, we do not strictly rank these close components.
Overall, clinical negatives and retrieval-grounded ARU verification provide the
strongest structural signal, while token-level weighting and the sample-level
margin add smaller but consistent optimization gains.}
\end{displayquote}

\textbf{4. Appendix objective-matrix table: add the requested
progressive filtering and objective-module matrix outside the main text.}

\begin{table}[htbp]
\centering
\small
\caption{\added{Progressive ablation separating response-level high-risk negative retention from the two Fa-DPO objective terms on the II-Medical-8B~\cite{intelligentinternet2025iimedical8b} backbone. Entries are macro-average Accuracy (\%) over the four main evaluation settings.}}
\label{tab:response_r114_objective_matrix}
\label{tab:ablation_objective_matrix}
\begin{adjustbox}{max width=\textwidth}
\begingroup\color{blue}
\begin{tabular}{lcccc}
\toprule
\textbf{Negative pool} & \textbf{Vanilla DPO loss} & \textbf{+ token weighting only} & \textbf{+ margin only} & \textbf{+ token weighting + margin} \\
\midrule
All QC candidates (47,393) & 61.6 & 62.1 & 61.9 & 62.4 \\
High-risk retained (34,204) & 62.0 & 62.2 & 61.7 & 63.2 \\
\bottomrule
\end{tabular}
\endgroup
\end{adjustbox}
\end{table}

\begin{displayquote}
\added{Table~\ref*{ms-tab:ablation_objective_matrix} separates response-level high-risk
negative retention from the two Fa-DPO objective terms in the representative
II-Medical-8B setting. Filtering alone is small and its interval includes zero
(+0.4, 95\% CI [-0.4, +1.2]). Within the high-risk retained pool, token weighting
alone gives a small gain over retained-pool vanilla DPO, whereas the margin-only
variant does not improve by itself in this representative setting. Their
combination gives the best result, suggesting that the response-level margin is
most useful when paired with localized token weighting rather than used as a
standalone objective change.}
\end{displayquote}

\begin{displayquote}
\added{These retained-pool values should therefore not be read as two independent
additive gains. Instead, they indicate coupling: the margin-only setting is
slightly below retained-pool vanilla DPO, but the same margin works best when
combined with localized token weighting.}
\end{displayquote}

\textbf{5. Appendix granularity-routing table: add the
granularity $\times$ routing interaction analysis outside the main text.}

\begin{table}[htbp]
\centering
\small
\caption{\added{Interaction between reasoning-unit granularity and routing strategy on the representative II-Medical-8B~\cite{intelligentinternet2025iimedical8b} setting. Accuracy is the macro-average over the four main evaluation settings; Support, Unsupported Reasoning, and Reasoning-Answer Inconsistency are automatic ARU/NLI-based process metrics (\%).}}
\label{tab:response_r114_granularity_routing}
\label{tab:ablation_granularity_routing}
\begin{adjustbox}{max width=\textwidth}
\begingroup\color{blue}
\begin{tabular}{llcccc}
\toprule
\textbf{Reasoning unit} & \textbf{Routing strategy} & \textbf{Avg. Acc.} & \textbf{Support} & \textbf{Unsupported} & \textbf{Claim Inc.} \\
\midrule
Whole-response block & Uniform routing & 60.8 & 80.1 & 26.4 & 22.8 \\
Sentence-level units & Uniform routing & 61.4 & 81.2 & 24.8 & 21.5 \\
Sentence-level units & Predicted role routing & 61.8 & 82.0 & 23.6 & 20.6 \\
ARU-level units & Uniform routing & 61.9 & 82.4 & 22.9 & 20.1 \\
ARU-level units & Role-specific routing & 63.2 & 84.0 & 20.9 & 18.5 \\
\bottomrule
\end{tabular}
\endgroup
\end{adjustbox}
\end{table}

\begin{displayquote}
\added{Table~\ref*{ms-tab:ablation_granularity_routing} shows that granularity and routing
work together. All rows use the same source prompts, retained pair count,
train/dev split, DPO training budget, final-answer extraction, NLI model, and
evaluation protocol; the row variables are the verifier preprocessing policy and
routing policy. Thus, whole-response, sentence-level, and ARU-level variants are
data- and training-matched, but they are not parser-cost-matched because parser use
and routing are deliberately varied. Relative to whole-response uniform
verification, the full ARU-level role-specific design improves average accuracy by
+2.4 points, Support Rate by +3.9, Unsupported Medical Reasoning Rate by -5.5, and
Reasoning-Answer Inconsistency Rate by -4.3. Role routing gives a larger gain at
ARU level (+1.3 accuracy) than at sentence level (+0.4), supporting local
role-aware verification.}
\end{displayquote}

\textbf{6. Sec.~6.3 and Appendix~D.6: merge the hyperparameter-sensitivity note
into the ablation section.}

The separate main-text sensitivity subsection is removed. The figure environment
itself remains in Appendix~D.6 without changing the curve caption.

\begin{displayquote}
\deleted{We next test whether Fa-DPO relies on a narrow hyperparameter setting. On
the representative II-Medical-8B~\cite{intelligentinternet2025iimedical8b}
paired setting, we vary one representative coefficient from each
faithfulness-aware component while keeping all other training and evaluation
settings fixed to Appendix~\ref*{ms-app:training_details}. Specifically, we sweep
the high-risk negative-response retention threshold $\gamma$ applied to the
response-level risk score in Eq.~\ref*{ms-eq:faithfulness_risk_score}, the ARU
risk threshold $\tau$ in Eq.~\ref*{ms-eq:severity_components} and
Eq.~\ref*{ms-eq:topology_mask}, and the sample-level margin scale $\eta$ in
Eq.~\ref*{ms-eq:severity_margin}.}

\added{To address possible coupling between modules, Appendix~D reports two
interaction-oriented ablations: one separating high-risk negative retention from
token weighting and margin, and another crossing reasoning-unit granularity with
routing strategy. These analyses show coupling between token weighting and the
response-level margin in the retained pool, and the strongest joint effect from
ARU-level granularity with role-specific routing.
Appendix~\ref*{ms-app:sensitivity} also sweeps the retention threshold $\gamma$,
the ARU-level verification-score threshold $\tau$, and the margin scale $\eta$
around their defaults; macro-average accuracy changes by at most 0.9 points and
Support Rate by at most 1.0 point, with the defaults near the best joint point.}
\end{displayquote}

\begin{displayquote}
\deleted{This appendix reports the representative results corresponding to the
sensitivity curves in Sec.~6.4. We keep the same
II-Medical-8B~\cite{intelligentinternet2025iimedical8b} paired setting and vary
only one hyperparameter at a time while fixing all others to the defaults above.
The default row in each block remains the same as the reported Fa-DPO
configuration in the main text.}

\added{This appendix reports the sensitivity curves and representative table
supporting the concise robustness statement in Sec.~\ref*{ms-sec:ablation}. We
keep the same II-Medical-8B~\cite{intelligentinternet2025iimedical8b} paired
setting and vary only one hyperparameter at a time while fixing all others to the
defaults above. The default row in each block remains the same as the reported
Fa-DPO configuration in the main text.}

\added{Figure~\ref*{ms-fig:sensitivity} shows that the default settings are near
the best joint point across the three sweeps. Accuracy and Support Rate peak
around $\gamma=0.5$, $\tau=0.1$, and $\eta=0.2$; moving away from these values
either weakens the risk signal or makes the retained/penalty signal too narrow or
too strong. Across the tested ranges, macro-average accuracy varies within 0.9
points and Support Rate within 1.0 point. The sensitivity result therefore
supports local robustness rather than introducing a separate source of gain.}
\end{displayquote}

\newpage
\begin{commentbox}{R1.15 LLM Judge Disclosure}
The blind LLM judge model, prompt template, and criteria are not disclosed. The
relationship between LLM-judge logical coherence and automatic claim inconsistency is
not discussed.
\end{commentbox}

\responseheading{Response and Revision Summary}
We thank the reviewer for noting that the blind LLM-judge evaluation lacked the
disclosure and interpretive detail needed for an evidence claim; in response, we
removed the blind LLM-judge subsection, GPT-5.5 preference results, win-rate figure,
and prompt-listing appendix from the manuscript. We chose deletion rather than
expanded disclosure because the judge was only an auxiliary proxy and the revised
manuscript now has direct clinician evidence for same-answer reasoning quality and
safety-relevant process outcomes. The evidence hierarchy now rests on matched
final-answer accuracy, automatic ARU/NLI process metrics, held-out parser/verifier
validation, the same-answer clinician audit, and the retrospective clinician
harm-severity audit.

\responseheading{Revision Details}
\textbf{1. Sec.~6.5.2: we removed the blind LLM-judge result from the main text.}

\begin{displayquote}
\deleted{This subsection asks where the process gains appear and whether they are visible beyond automatic scores. We first measure ARU-level unfaithfulness by reasoning function to locate the changes inside the reasoning chain. We then use blind pairwise reasoning evaluation on the same-answer setting to test whether the same direction is visible in holistic comparison.}

\added{To localize the process-metric gains, we further examine ARU-level unfaithfulness rates by ARU type, including Observation, Warrant, Differentiation, and Claim ARUs.}

\deleted{We next test whether the representative same-answer gains remain visible under holistic pairwise judgment. To do so, we use a blind LLM judge as an auxiliary comparative reader at fixed answer correctness, since final accuracy alone does not distinguish faithful from unfaithful reasoning.}

\deleted{A blind LLM judge performs a pairwise comparison across three dimensions:}

\deleted{The blind-judge win-rate figure shows a consistent Fa-DPO preference across all three judged dimensions. The strongest margin is on \textit{Hallucination Freedom} (68/20/12 win/tie/loss), followed by \textit{Overall Preference} (63/22/15); \textit{Logical Coherence} remains positive but has more ties (57/25/18).}

Table~\ref*{ms-tab:iimedical_aru_shift} \added{localizes the process gain. Overall ARU-level unfaithfulness drops from 22.1\% to 15.4\%, and every ARU type improves with confidence intervals excluding zero. The largest reductions appear in differentiation (-10.9 pp) and warrant ARUs (-9.6 pp), while claim ARUs also decrease by 8.4 pp. The gain is therefore strongest in medically inferential reasoning units, but it remains visible across the full chain. We next test whether these process-level improvements coincide with better final-answer accuracy.}
\end{displayquote}

\textbf{2. Discussion: we removed blind-pairwise evidence from the mechanism interpretation.}

\begin{displayquote}
\deleted{The function-wise ARU analysis and blind pairwise evaluation further show that the gains are visible both in verifier-derived local process indicators and in holistic judgments of reasoning quality.}

\added{The ARU-type analysis shows improvements across Observation, Warrant, Differentiation, and Claim units, with the strongest reductions in inferential reasoning units.}
\end{displayquote}

\textbf{3. Appendix: we removed the blind-judge protocol, table, and prompt listing.}

\begin{displayquote}
\deleted{\textbackslash section\{LLM-Based Blind Pairwise Evaluation Protocol\}}

\deleted{This appendix records the protocol used for blind pairwise judging. The evaluation starts from the shared same-answer hard subset and applies no additional filtering. For each example, the two anonymized responses are randomly assigned to slot A or B with seed 42.}

\deleted{The judge call uses a two-message chat prompt. In the reported experiments, we request \texttt{gpt-5.5} as the blinded LLM judge through a chat API at temperature 0.0 with a maximum completion length of 900 tokens.}

\deleted{Prompt-sensitivity and separate LLM-judge checks for blind pairwise preferences on the II-Medical-8B same-answer hard subset ($n=100$). The 95\% CIs are for non-tie Fa-DPO win rates.}

\deleted{Listing 1: Blind pairwise judge prompt template used in the released evaluation script.}
\end{displayquote}

\newpage
\begin{commentbox}{R1.16 NLI Held-Out Distribution Shift}
The NLI verifier reports 0.838 accuracy, but it is unclear whether the held-out set
matches the training set. If the held-out set contains constructed negatives while
real inference faces natural model errors, distribution shift may occur.
\end{commentbox}

\responseheading{Response and Revision Summary}
We thank the reviewer for raising this point. We now distinguish the held-out
in-pipeline verifier split from natural-output distribution shift, so the originally
reported verifier accuracy is not read as evidence of equal reliability on natural
model errors.

We added a natural-output stress test and revised the interpretation accordingly. The
verifier remains a diagnostic scoring component, but the manuscript now explicitly
limits automatic process-metric claims under natural-output degradation.

\responseheading{Revision Details}
\textbf{1. Sec.~6.6: we separated natural-output verifier shift from the original
in-pipeline verifier validation in the main summary.}

\begin{displayquote}
\deleted{We next assess whether the evidence used in the preceding analyses is reliable.}
\added{This section validates the process evidence used in the preceding analyses. It asks three linked questions: whether automatic same-answer improvements align with clinician judgments and data-audit evidence; whether the ARU parser, routed NLI verifier, and token masks are reliable enough for proxy process measurement; and whether the same signal remains directionally useful on natural outputs, answer-changing diagnostics, and retrospective safety checks. These experiments support bounded use of the process metrics, not prospective clinical validation.}

\deleted{Second, we validate the routed NLI verifier directly on held-out ARU-premise
pairs.}\added{Pipeline audits validate the automatic scoring components. The held-out parser meets the main span and role targets at the primary IoU threshold, but role macro-F1 and final-answer claim-span F1 limit exact role-wise attribution. The routed NLI verifier reaches 0.838 accuracy and 0.827 macro-F1 on the held-out in-pipeline medical split, while natural-output stress testing shows moderate degradation, including a high-risk false-negative increase from 0.118 to 0.172. Clinician false-positive and training-mask audits further show that high-risk flags are mostly meaningful but can include clinically acceptable evidence-unsupported cases. Table~\ref*{ms-tab:pipeline_validation_summary} summarizes these component-level checks. These results justify using the automatic metrics as diagnostic indicators, not as clinical ground truth; detailed protocols and tables are in Appendices~\ref*{ms-app:aru_parser_validation} and~\ref*{ms-app:nli_validation}.}
\end{displayquote}

\textbf{2. Sec.~6.6.2: we clarified that the 0.838 accuracy comes from a held-out
in-pipeline split, not from natural model outputs.}

\begin{displayquote}
\added{The in-pipeline protocol samples $n=200$ medical ARU-premise pairs with stratification over route type, ARU reasoning function, predicted label, and difficulty band. We reserve 40 pairs as a small stratified development split and exclude them from all reported held-out metrics; the final verifier results are reported on the remaining 160 held-out in-pipeline medical evaluation pairs.}

\deleted{held-out evaluation}\added{held-out in-pipeline medical evaluation}
\end{displayquote}

\begin{displayquote}
\textit{Manuscript table-caption excerpt:} \added{Verifier-level NLI validation on the held-out in-pipeline medical evaluation split ($n=160$). A separate stratified development split ($n=40$) is reserved and excluded from the reported held-out metrics.}
\end{displayquote}

\textbf{3. Sec.~6.6.2 and appendix: we added the natural-output distribution-shift
stress test and reported the constructed-to-natural performance gap.}

\begin{displayquote}
\captionof{table}{\added{Natural-output distribution-shift analysis for the routed NLI verifier. High-risk FPR is the fraction of gold-entailment ARUs incorrectly flagged as high risk; high-risk FNR is the fraction of gold non-entailment ARUs missed by the high-risk rule.}}
\begin{adjustbox}{max width=\linewidth}
\begingroup\color{blue}
\begin{tabular}{llccccc}
    \toprule
    Split & Source & $n$ & Acc. & Macro-F1 & High-risk FPR & High-risk FNR \\
    \midrule
    Constructed held-out & In-pipeline controlled construction & 160 & 0.838 & 0.827 & 0.063 & 0.118 \\
    Natural in-domain & MedCaseReasoning model outputs & 180 & 0.794 & 0.775 & 0.093 & 0.172 \\
    \bottomrule
\end{tabular}
\endgroup
\end{adjustbox}
\end{displayquote}

\begin{displayquote}
\added{Together, the held-out split, natural-output stress test, false-positive audit, and training-mask audit define the verifier's usable scope. They support the verifier as a proxy scoring component for process analysis, while limiting claims about clinical ground truth, exhaustive medical correctness, or complete robustness on natural outputs.}
\end{displayquote}

\textbf{4. Conclusion and Future Work: we retained the natural-output boundary in the
compressed limitation paragraph.}

\begin{displayquote}
\added{The added audits are mostly retrospective and partly automatic; they do not establish broad hallucination detection or deployment safety, and the retained negatives do not cover the full distribution of naturally generated medical-LLM errors.}
\end{displayquote}

\newpage
\begin{commentbox}{R1.17 Controlled Taxonomy Coverage}
Only 40.8\% of external errors can be mapped to F1--F4, meaning 59.2\% are uncovered.
This challenges the completeness of the taxonomy and the claim that the method
simulates realistic clinical reasoning failures.
\end{commentbox}

\responseheading{Response and Revision Summary}
We thank the reviewer for raising this concern. After re-examining the analysis, we
agree that the Med-HALT pilot was not well aligned with the central evaluation construct
of our paper. Fa-DPO is designed for evidence-grounded MQA, where the key question is
whether model reasoning faithfully uses provided clinical evidence. In contrast,
Med-HALT is a hallucination-oriented benchmark that includes medical exam reasoning,
fake-question detection, none-of-the-above tests, and memory-based PubMed retrieval
tasks. These tasks do not consistently provide the patient-specific evidence units
required for ARU-level evidence-use analysis.

Therefore, mapping Med-HALT errors into the F1--F4 evidence-use taxonomy introduces a
construct mismatch: low coverage may reflect benchmark mismatch rather than a limitation
of the proposed taxonomy in the intended evidence-grounded MQA setting. We removed the
Med-HALT pilot, the corresponding coverage table, and the 40.8\%/59.2\% coverage
discussion from the revised manuscript. We removed this construct-mismatched analysis
to avoid drawing an unsupported inference about the intended evidence-grounded MQA
setting.

We also narrowed the external-validity claim and replaced the pilot with a
task-aligned held-out natural-error analysis on evidence-grounded MQA prompts. The
new analysis is presented as bounded transfer evidence, not as proof that F1--F4 is a
complete taxonomy of natural medical errors. We also expanded the annotation-protocol
description and kept the limitation that these labels are automatic and multi-label.


\responseheading{Revision Details}
\textbf{Experimental settings: removed Med-HALT and narrowed the external-evaluation scope.}

\begin{displayquote}
\deleted{We further use a sampled Med-HALT reasoning pilot set as a supplementary external check of natural-error coverage.}

\added{The evidence also shows bounded task-family generalization. The matched benchmark results cover multiple evidence-grounded MQA datasets and general-instruction, medical-adapted, and reasoning-enhanced backbones.}
\end{displayquote}

\textbf{Reasoning-analysis set and appendix protocol: removed the Med-HALT pilot material.}

\begin{displayquote}
\deleted{Med-HALT reasoning is an external benchmark designed to probe medical hallucinations and reasoning errors. We use a 120-response pilot sample from it as a supplementary natural-error coverage check, asking whether sampled external outputs can be mapped to the controlled F1--F4 taxonomy. The appendix reports the pilot sampling and response-level labeling protocol.}

\deleted{Pilot Natural-Error Coverage Protocol on Med-HALT.}

\deleted{This appendix records the pilot protocol used for the Med-HALT study. The goal is not to compare Fa-DPO against Standard DPO under matched conditions, but to ask a narrower external-validity question: whether sampled errors from an external medical reasoning benchmark can be described using the controlled F1--F4 taxonomy introduced in Sec. 3.}
\end{displayquote}

\textbf{Held-out natural-error analysis: task-aligned replacement for the Med-HALT pilot.}
\begin{displayquote}
\added{To test transfer beyond controlled counterfactual rewrites, we evaluate the representative II-Medical-8B matched pair on 3,338 model-sampled cases from MedXpertQA-test, MedBullets-op4, MedBullets-op5, and MedQA-test. This analysis compares Standard DPO-Retained-34K and Fa-DPO-Retained-34K responses on the same held-out prompts without local rewriting. The labels are automatic rather than double-clinician-adjudicated and may co-occur, so this analysis supports a fixed-protocol consistency check rather than definitive clinical ground truth.}

\added{The held-out natural-error analysis is conducted as a frozen post hoc evaluation over already generated model responses. Official answers are used only after generation to determine answer correctness and assign response-level error labels for analysis; they are not included in prompts, training data, candidate retention, hyperparameter tuning, or model selection. Each row contains one generated response with the prompt, answer options, gold option, predicted option, and visible reasoning. The label definitions follow the evidence-grounded F1--F4 schema but are applied to natural outputs. F1 marks omission or substantive weakening of clinically important question evidence, F2 marks incorrect medical interpretation of evidence that the response uses, F3 marks overstatement beyond what the available facts support, and F4 marks a selected answer that is unsupported by or inconsistent with the preceding reasoning. An unmapped error is a substantive medical reasoning or knowledge error not cleanly captured by F1--F4; unsupported medical claim and harmful recommendation are recorded as separate safety-relevant flags. The labels are automatic rather than double-clinician-adjudicated, and they are multi-label: F1--F4 may co-occur with one another and with unmapped, unsupported-claim, or harmful-recommendation flags.}

\added{Together, Tables~\ref*{ms-tab:heldout_natural_error_analysis} and~\ref*{ms-tab:heldout_natural_standard_strata} show bounded transfer to natural model generations. Fa-DPO increases micro accuracy by 1.75 points, reduces any mapped F1--F4 error by 7.60 points, and lowers unsupported medical claims from 5.20\% to 3.92\%. The largest mapped reduction is F4 reasoning-answer mismatch. In the Standard-DPO mapped-error stratum, Fa-DPO is correct without a mapped error in 27.4\% of cases, but still has a mapped error in 66.8\%. Thus, the effect is meaningful but far from complete error removal, and the small unmapped stratum remains boundary evidence rather than a stable estimate of broad natural-error coverage.}
\end{displayquote}

\textbf{Discussion and Conclusion: task-aligned generalization boundary.}
\begin{displayquote}
\added{The evidence also shows bounded task-family generalization. The matched benchmark results cover multiple evidence-grounded MQA datasets and general-instruction, medical-adapted, and reasoning-enhanced backbones. The held-out natural-error analysis on 3,338 model-sampled MQA cases further shows fewer mapped F1--F4 errors and fewer unsupported medical claims for the representative II-Medical-8B pair. These findings indicate transfer beyond controlled rewrites within the same task family, while keeping the scope limited: F1--F4 is an evidence-use taxonomy for this setting, not a complete account of all medical hallucination or clinical reasoning errors.}

\added{The evaluation remains bounded: experiments use 8B models, evidence-grounded multiple-choice MQA, and controlled negative construction with ARU parsing, retrieval, and NLI verification.}

\added{The added audits are mostly retrospective and partly automatic; they do not establish broad hallucination detection or deployment safety, and the retained negatives do not cover the full distribution of naturally generated medical-LLM errors.}
\end{displayquote}

\newpage
\tocsection{Reviewer \#2}

\begin{commentbox}{R2 General Assessment}
Reviewer \#2 finds the manuscript technically novel and well executed, but identifies
four major weaknesses: circularity in process-faithfulness evaluation, limited
generalizability beyond the controlled taxonomy, insufficient validation of the ARU
parser, and missing comparisons to non-DPO faithfulness baselines.
\end{commentbox}

\responseheading{Response and Revision Summary}
We thank the reviewer for the positive assessment of the manuscript's technical
motivation and empirical execution. We agree that the four weaknesses identified by
the reviewer required clearer validation boundaries and additional analyses. In the
revision, we therefore (i) reframe ARU/NLI process metrics as reasoning-process
diagnostic indicators rather than independent clinical ground truth and add a
    completed same-answer clinician validation, (ii) remove the construct-mismatched
    Med-HALT pilot and replace it with task-aligned MQA evaluations across six
    backbones plus a held-out natural-error analysis on the representative
    II-Medical-8B matched pair, (iii) complete a held-out clinician-adjudicated ARU parser validation
with explicit span and role metrics, and (iv) add a faithfulness-oriented baseline
protocol covering SFT-only, full-candidate DPO, retained-pair DPO, Fa-DPO, and
inference-time verifier reranking.

\newpage
\begin{commentbox}{R2.1 Circularity in Process-Faithfulness Evaluation}
The main process-faithfulness metrics use the same NLI verifier and ARU decomposition
pipeline that provides Fa-DPO's training signal. The authors should validate these
metrics against clinician judgments or explicitly reframe them as diagnostic-indicator
metrics.
\end{commentbox}

\responseheading{Response and Revision Summary}
We agree with the circularity concern: the automatic process metrics are computed
by the same ARU parser and NLI verifier used to construct Fa-DPO's risk signal.
The revised manuscript therefore states that these metrics are diagnostic indicators,
not independent clinical ground truth or standalone proof of faithfulness.

We also added a blinded same-answer clinician audit to test whether the automatic
direction is clinically plausible when final-answer correctness is fixed. Clinicians
preferred Fa-DPO's reasoning across all five criteria, matching the direction of the
automatic metrics. We use this audit as directional support only; the detailed
validation protocol and tables are reported in R2.5.


\responseheading{Revision Details}
\textbf{Evaluation framing and limitations: directly acknowledge the same-source
circularity.}
\begin{displayquote}
We evaluate Fa-DPO with three groups of metrics: final-answer accuracy,
\deleted{process faithfulness}\added{automatic process metrics}, and ARU
unfaithfulness shift by reasoning function.

\deleted{These automatic process metrics provide mechanism-oriented proxy evidence rather than direct measures of absolute clinical correctness.}

\added{Because these process metrics are computed by the same ARU parser and NLI verifier used to construct the Fa-DPO risk signal, they are interpreted as diagnostic indicators of evidence-use and consistency behavior rather than independent clinical ground truth or direct measures of absolute clinical correctness. The same-answer clinician audit in Appendix~\ref*{ms-app:same_answer_clinician_details} provides directional validation of these indicators, not a replacement for clinician-adjudicated evaluation.}

\added{Some process metrics are computed using the same ARU/NLI verification framework that provides training signals. Therefore, these metrics should not be interpreted as fully independent proof of faithfulness. Reverse entailment is likewise treated only as a possible coverage or specificity diagnostic, not as a reported result or training signal. We address these limitations by reporting answer accuracy and clinician audits as complementary evidence.}
\end{displayquote}

\textbf{Same-answer clinician audit: support the direction, not clinician labels.}
\begin{displayquote}
\added{Human audits validate the data and same-answer mechanism evidence. In the same-answer audit, clinicians prefer Fa-DPO's reasoning when final-answer correctness is fixed, matching the direction of the automatic process metrics.}

\added{To validate the same-answer process metrics, we test whether their direction matches clinician judgments when answer correctness is fixed.}

\added{Clinicians prefer Fa-DPO on all five same-answer criteria, with non-tie win rates of 77.3--80.6\%. Response-level pass rates move in the same direction: the overall clinically coherent reasoning rate increases from 70.0\% for Standard DPO-Retained-34K to 84.0\% for Fa-DPO-Retained-34K.}
\end{displayquote}

\newpage
\begin{commentbox}{R2.2 Generalizability Beyond Controlled Taxonomy}
The F1--F4 taxonomy covers only 40.8\% of sampled external errors in the earlier pilot.
The authors should report Fa-DPO vs Standard DPO performance on mapped and unmapped
subsets separately and discuss the limitation.
\end{commentbox}

\responseheading{Response and Revision Summary}
We agree that the original Med-HALT pilot was not an appropriate basis for assessing generalization beyond the controlled F1--F4 taxonomy. Med-HALT mixes hallucination-oriented, exam-style, and memory-based settings, whereas Fa-DPO targets evidence-grounded MQA. We therefore removed the Med-HALT pilot, its coverage table, and the 40.8\%/59.2\% interpretation.

To address the reviewer's concern in a task-aligned way, we added a held-out natural-generation analysis on 3,338 evidence-grounded MQA cases using the representative II-Medical-8B pair. The analysis compares Standard DPO-Retained-34K and Fa-DPO-Retained-34K responses on the same prompts, without local rewriting, and records mapped F1--F4 errors, unmapped errors, unsupported medical claims, and harmful recommendations.

To avoid any possible concern about answer leakage, we clarified that official answers are used only after model outputs are generated, for answer-correctness determination and response-level error annotation. They are not used in prompts, training, candidate retention, hyperparameter tuning, or model selection.

Fa-DPO increases micro accuracy from 47.20\% to 48.95\%, reduces any mapped F1--F4 error from 70.00\% to 62.40\%, and lowers unsupported medical claims from 5.20\% to 3.92\%. Unmapped errors remain rare in this analysis and change from 1.10\% to 0.96\%; we use this only as a boundary check, not as evidence of broad improvement on taxonomy-outside errors.

We also clarify two levels of generalization. Results across multiple evidence-grounded MQA datasets and backbone types show that Fa-DPO is not limited to one benchmark or model, and the held-out natural-generation analysis shows transfer beyond controlled rewrites to mapped natural errors. We do not claim that F1--F4 covers all medical errors or that Fa-DPO improves every error outside this taxonomy.


\responseheading{Revision Details}
\textbf{Experimental settings: remove construct-mismatched Med-HALT pilot.}
\begin{displayquote}
\deleted{Med-HALT reasoning is an external benchmark designed to probe medical hallucinations and reasoning errors. We use a 120-response pilot sample from it as a supplementary natural-error coverage check, asking whether sampled external outputs can be mapped to the controlled F1--F4 taxonomy. The appendix reports the pilot sampling and response-level labeling protocol.}

\added{The evidence also shows bounded task-family generalization. The matched benchmark results cover multiple evidence-grounded MQA datasets and general-instruction, medical-adapted, and reasoning-enhanced backbones.}
\end{displayquote}

\textbf{Appendix protocol: delete the Med-HALT coverage protocol.}
\begin{displayquote}
\deleted{Pilot Natural-Error Coverage Protocol on Med-HALT.}

\deleted{This appendix records the pilot protocol used for the Med-HALT study. The goal is not to compare Fa-DPO against Standard DPO under matched conditions, but to ask a narrower external-validity question: whether sampled errors from an external medical reasoning benchmark can be described using the controlled F1--F4 taxonomy introduced in Sec. 3.}
\end{displayquote}

\textbf{Held-out natural-error analysis added to the validation section.}
\begin{displayquote}
\added{To test transfer beyond controlled counterfactual rewrites, we evaluate the representative II-Medical-8B matched pair on 3,338 model-sampled cases from MedXpertQA-test, MedBullets-op4, MedBullets-op5, and MedQA-test. This analysis compares Standard DPO-Retained-34K and Fa-DPO-Retained-34K responses on the same held-out prompts without local rewriting. The labels are automatic rather than double-clinician-adjudicated and may co-occur, so this analysis supports a fixed-protocol consistency check rather than definitive clinical ground truth.}

\added{The held-out natural-error analysis is conducted as a frozen post hoc evaluation over already generated model responses. Official answers are used only after generation to determine answer correctness and assign response-level error labels for analysis; they are not included in prompts, training data, candidate retention, hyperparameter tuning, or model selection. Each row contains one generated response with the prompt, answer options, gold option, predicted option, and visible reasoning. The label definitions follow the evidence-grounded F1--F4 schema but are applied to natural outputs. F1 marks omission or substantive weakening of clinically important question evidence, F2 marks incorrect medical interpretation of evidence that the response uses, F3 marks overstatement beyond what the available facts support, and F4 marks a selected answer that is unsupported by or inconsistent with the preceding reasoning. An unmapped error is a substantive medical reasoning or knowledge error not cleanly captured by F1--F4; unsupported medical claim and harmful recommendation are recorded as separate safety-relevant flags. The labels are automatic rather than double-clinician-adjudicated, and they are multi-label: F1--F4 may co-occur with one another and with unmapped, unsupported-claim, or harmful-recommendation flags.}

\added{Together, Tables~\ref*{ms-tab:heldout_natural_error_analysis} and~\ref*{ms-tab:heldout_natural_standard_strata} show bounded transfer to natural model generations. Fa-DPO increases micro accuracy by 1.75 points, reduces any mapped F1--F4 error by 7.60 points, and lowers unsupported medical claims from 5.20\% to 3.92\%. The largest mapped reduction is F4 reasoning-answer mismatch. In the Standard-DPO mapped-error stratum, Fa-DPO is correct without a mapped error in 27.4\% of cases, but still has a mapped error in 66.8\%. Thus, the effect is meaningful but far from complete error removal, and the small unmapped stratum remains boundary evidence rather than a stable estimate of broad natural-error coverage.}
\end{displayquote}

\textbf{Discussion: bounded task-aligned transfer interpretation.}
\begin{displayquote}
\added{The evidence also shows bounded task-family generalization. The matched benchmark results cover multiple evidence-grounded MQA datasets and general-instruction, medical-adapted, and reasoning-enhanced backbones. The held-out natural-error analysis on 3,338 model-sampled MQA cases further shows fewer mapped F1--F4 errors and fewer unsupported medical claims for the representative II-Medical-8B pair. These findings indicate transfer beyond controlled rewrites within the same task family, while keeping the scope limited: F1--F4 is an evidence-use taxonomy for this setting, not a complete account of all medical hallucination or clinical reasoning errors.}
\end{displayquote}

\newpage
\begin{commentbox}{R2.3 ARU Parser Validation}
The ARU decomposition step is foundational, but Appendix E provides only heuristic
quality controls. The authors should report parser performance on clinician-annotated
ARUs, including span boundary F1 and role classification accuracy.
\end{commentbox}

\responseheading{Response and Revision Summary}
We agree that the ARU parser requires direct validation because it supports evidence routing, NLI scoring, risk estimation, and token-level weighting. We therefore added a held-out clinician-adjudicated parser validation on 100 responses, disjoint from parser prompt-development examples.

Two clinicians annotated ARU span boundaries and role labels, with adjudication before scoring. The adjudicated gold set contains 1,124 ARUs. Parser outputs were aligned to gold ARUs by maximum bipartite matching over token-level span IoU.

At the primary IoU $\geq 0.5$ threshold, the parser reaches span boundary F1 0.863 and role classification accuracy 0.802, meeting the requested targets of span F1 $\geq 0.85$ and role accuracy $\geq 0.80$. We also report a stricter IoU $\geq 0.7$ sensitivity check. These results support the parser as reliable enough for directional ARU-based process analysis.

This revision also addresses the specific threshold request in R2.7. To avoid overclaiming, we clarify that parser-derived process metrics are diagnostic indicators, while benchmark final-answer accuracy is measured independently of the parser.

We added the clinician-adjudicated parser-validation protocol, the matching and scoring procedure, the validation table, and a short interpretation of the remaining limits of parser-derived process metrics.


\responseheading{Revision Details}
\textbf{Sec.~6.6 and Appendix~E: add completed parser-validation protocol and results.}
\begin{displayquote}
\added{To validate the ARU decomposition parser used by the process metrics, this protocol evaluates parser output against clinician-adjudicated ARU boundaries and role labels. The held-out validation set contains 100 responses that were not used to design parser prompts or tune quality-control heuristics. The split includes constructed clinical reasoning negatives covering the controlled failure types and the four ARU roles in Sec.~\ref*{ms-sec:decomposition}.}

\added{Two clinicians independently annotate each response. Given the question, case or evidence, model response, reasoning trace, and final answer, annotators mark the minimal clinical reasoning units in the response. For every unit, they record the span boundary, one role label in \{Observation, Warrant, Differentiation, Claim\}, whether the span is the final-answer claim, and whether the span is eligible to enter the token-level risk mask if it is later judged unsupported, contradictory, or reasoning-answer inconsistent. Clinicians do not assign NLI labels in this parser-validation task. Disagreements are resolved by discussion or by a third adjudicator before parser performance is computed.}

\added{Predicted and gold ARUs are aligned by maximum bipartite matching over token-level span overlap. Character offsets are first converted to token spans with the tokenizer used for Eq.~\ref*{ms-eq:weighted_rejected_logratio}. A predicted ARU is considered matched to a gold ARU when their token-level IoU is at least 0.5; IoU $\geq 0.7$ is additionally reported as a stricter sensitivity threshold. Span precision, recall, and F1 are computed from the matched predicted and gold ARU counts. Role accuracy and role macro-F1 are computed only on matched spans. The validation also reports final-answer claim-span F1 and token-mask F1 between the parser-derived and clinician-derived mask-eligible token spans.}

\added{Inter-annotator agreement is reported before adjudication using span-level F1 for boundary agreement and Cohen's $\kappa$ for role labels on annotator-matched spans. In the completed validation, pre-adjudication span F1 is 0.915 at IoU $\geq 0.5$ and role-label $\kappa$ is 0.818 at the same threshold. After adjudication, the parser reaches span F1 0.863 and role accuracy 0.802 at IoU $\geq 0.5$, meeting the reviewer-specified targets of span F1 $\geq 0.85$ and role accuracy $\geq 0.80$. At the stricter IoU $\geq 0.7$ sensitivity threshold, the parser reaches span F1 0.764 and role accuracy 0.822, indicating that the remaining errors are mainly strict boundary differences rather than broad role failures.}

\added{Overall, the parser is accurate enough for directional process analysis, but two limits remain. Role macro-F1 is lower than overall role accuracy because role labels are imbalanced and the metric is computed only on matched spans; minority-role errors therefore weaken exact role-wise attribution, especially for less frequent roles. Final-answer claim F1 is also low under strict span overlap because the final answer is extracted and normalized before being appended as a claim ARU, while clinicians may mark a longer or shorter answer phrase. Thus, the parser results do not suggest that Fa-DPO's accuracy gains are parser artifacts, but they do limit how strongly we can read parser-derived role-wise and claim-level process metrics.}
\end{displayquote}

\textbf{Appendix parser-validation table.}
\begin{displayquote}
\begin{center}
\footnotesize
\captionof{table}{\added{Held-out ARU parser validation against clinician-adjudicated annotations on 100 responses. The primary validation targets are span boundary F1 $\geq 0.85$ and role classification accuracy $\geq 0.80$ at IoU $\geq 0.5$; IoU $\geq 0.7$ is reported as a stricter sensitivity threshold.}}
\begin{adjustbox}{max width=\linewidth}
\begingroup\color{blue}
\begin{tabular}{lcc}
\toprule
\textbf{Metric} & \textbf{IoU $\geq 0.5$} & \textbf{IoU $\geq 0.7$} \\
\midrule
Responses & 100 & 100 \\
Gold ARUs & 1,124 & 1,124 \\
Parser ARUs & 1,134 & 1,134 \\
Matched ARUs & 974 & 863 \\
Span precision & 0.859 & 0.761 \\
Span recall & 0.867 & 0.768 \\
Span F1 & \textbf{0.863} & 0.764 \\
Role accuracy on matched spans & \textbf{0.802} & 0.822 \\
Role macro-F1 & 0.627 & 0.668 \\
Final-answer claim F1 & 0.471 & 0.471 \\
Token-mask F1 & 0.784 & 0.722 \\
\bottomrule
\end{tabular}
\endgroup
\end{adjustbox}
\end{center}
\end{displayquote}

\newpage
\begin{commentbox}{R2.4 Missing Non-DPO Faithfulness Baselines}
The paper compares to Standard DPO but not to simpler alternatives such as supervised
fine-tuning on faithful reasoning examples, standard DPO with filtered negatives but
without risk-adaptive weighting, or inference-time verification plus rejection
sampling.
\end{commentbox}

\responseheading{Response and Revision Summary}
We agree that the original evaluation did not fully separate Fa-DPO's objective from
simpler uses of faithfulness information. We therefore added a representative
faithfulness-oriented baseline suite on II-Medical-8B, including positive-only SFT,
full-candidate DPO, random-size-matched DPO, retained-pair DPO, and verifier
reranking.

The added comparison shows that verifier reranking is a strong inference-time
alternative, which supports the usefulness of the ARU/NLI faithfulness signal itself.
Fa-DPO achieves the best overall single-generation result and a slightly higher
macro-average accuracy than K=8 reranking, while avoiding the multi-generation and
online ARU/NLI verification cost required by reranking. We therefore present this
result as evidence that Fa-DPO can internalize much of the verifier-guided preference
signal during training, rather than as a large-margin win over reranking.


\responseheading{Revision Details}
\textbf{Baselines subsection: add faithfulness-oriented control layers.}
\begin{displayquote}
\deleted{We organize the comparison systems into three categories and use them in two ways: controlled matched Fa-DPO vs.\ Standard DPO comparisons and a broader released-baseline comparison.}

\added{We use three baseline layers. Released-model baselines contextualize competitiveness across general instruction, medical-adapted, and reasoning-oriented medical models; the complete released-model list is reported in Table~\ref*{ms-tab:competitiveness_reference}. Matched retained-pair DPO controls provide the primary causal comparison, and representative faithfulness-oriented controls test simpler explanations for the observed gain.}
\end{displayquote}

\textbf{Main Results: baseline table and deployment-cost fields.}
\begin{displayquote}
\added{On the representative II-Medical-8B backbone, we additionally compare against SFT-faithful, Standard DPO-QC-All, Standard DPO-Random-34K, Standard DPO-Retained-34K, and verifier reranking with $K=4$ or $K=8$. These controls separate positive-only imitation, full-candidate DPO, sample-count effects, retained-pool effects, and inference-time verification. Appendix~\ref*{ms-app:faithfulness_baseline_protocol} gives the control protocol.}
\end{displayquote}

\begin{table}[htbp]
\centering
\scriptsize
\setlength{\tabcolsep}{3.0pt}
\caption{\added{Faithfulness-oriented baseline controls on the representative II-Medical-8B backbone. Accuracy is the macro-average over MedBullets-op4, MedBullets-op5, MedXpertQA, and MedQA. Process metrics are automatic ARU/NLI-based percentages on the MedCaseReasoning same-answer hard subset. The table also records inference generations and whether deployment-time ARU/NLI verification is required; it does not report measured reranking latency for $K=4$ or $K=8$.}}
\label{tab:response_r24_faithfulness_baseline_results}
\begin{adjustbox}{max width=\textwidth}
\begin{tabular}{lccccccc}
\toprule
\textbf{Method} & \textbf{Avg. Acc.} & \shortstack[c]{\textbf{Support}\\\textbf{Rate $\uparrow$}} & \shortstack[c]{\textbf{Contrad.}\\\textbf{Rate $\downarrow$}} & \shortstack[c]{\textbf{Unsup.}\\\textbf{Reasoning $\downarrow$}} & \shortstack[c]{\textbf{Reason.-Ans.}\\\textbf{Incons. $\downarrow$}} & \shortstack[c]{\textbf{Infer.}\\\textbf{gen.}} & \shortstack[c]{\textbf{Deploy}\\\textbf{ARU/NLI?}} \\
\midrule
\added{Base model} & \added{61.5} & \added{75.8} & \added{10.7} & \added{34.2} & \added{28.4} & \added{1} & \added{no} \\
\added{SFT-faithful} & \added{62.2} & \added{79.1} & \added{8.3} & \added{28.7} & \added{23.9} & \added{1} & \added{no} \\
\added{Standard DPO-QC-All} & \added{61.6} & \added{76.6} & \added{9.8} & \added{32.7} & \added{27.1} & \added{1} & \added{no} \\
\added{Standard DPO-Random-34K} & \added{61.7} & \added{76.9} & \added{9.6} & \added{32.0} & \added{26.6} & \added{1} & \added{no} \\
\added{Standard DPO-Retained-34K} & \added{62.0} & \added{77.3} & \added{9.4} & \added{31.6} & \added{26.1} & \added{1} & \added{no} \\
\added{Standard DPO-Retained-34K + rerank, $K=4$} & \added{62.8} & \added{82.2} & \added{5.8} & \added{23.4} & \added{20.5} & \added{4} & \added{yes} \\
\added{Standard DPO-Retained-34K + rerank, $K=8$} & \added{63.0} & \added{83.1} & \added{5.4} & \added{22.2} & \added{19.8} & \added{8} & \added{yes} \\
\added{\textbf{Fa-DPO-Retained-34K}} & \added{\textbf{63.2}} & \added{\textbf{84.0}} & \added{\textbf{5.2}} & \added{\textbf{20.9}} & \added{\textbf{18.5}} & \added{1} & \added{no} \\
\bottomrule
\end{tabular}
\end{adjustbox}
\end{table}

\begin{displayquote}
\added{Verifier reranking narrows the gap, especially at $K=8$, but requires multiple generations and online ARU/NLI verification; Fa-DPO keeps single-generation inference.}
\end{displayquote}

\textbf{Appendix protocol: fair setup and verifier reranking.}
\begin{displayquote}
\added{The representative baseline controls in Table~\ref*{ms-tab:faithfulness_baseline_results} use the II-Medical-8B~\cite{intelligentinternet2025iimedical8b} backbone and the same benchmark splits, prompts, answer-extraction rule, and process-metric pipeline as the matched II-Medical-8B comparison. SFT-faithful trains on the exact preferred-response side ($y_w$) used in the retained preference pairs, deduplicated by prompt and target response; it uses the same tokenizer, optimizer family, LoRA setup, sequence limits, and one-epoch supervised schedule, but it is not an update-matched DPO run because it is a positive-only supervised control.}

\added{Verifier reranking does not add a new training run. It uses the Standard DPO-Retained-34K generator, samples $K$ responses per question under the same decoding profile, decomposes each candidate into ARUs, scores candidates with the same routed-premise NLI verifier, and selects the response with the lowest response-level risk score.}
\end{displayquote}

\textbf{Discussion: added controls support Fa-DPO.}
\begin{displayquote}
\added{Baseline controls further show that the gains are not explained by positive-only imitation, vanilla DPO, sample-count effects, or post-hoc verifier reranking alone. Overall, Fa-DPO works best when localized verifier risk is built into offline preference optimization.}
\end{displayquote}

\newpage
\begin{commentbox}{R2.5 Specific Comment: Clinician Validation of Metrics}
Add a validation subsection where clinicians evaluate 100--200 responses from the
MedCaseReasoning same-answer subsets, or explicitly reframe the metrics as
diagnostic-indicator metrics and temper claims.
\end{commentbox}

\responseheading{Response and Revision Summary}
This comment is addressed together with R2.1. We added a new same-answer clinician
validation audit and explicitly reframed the automatic metrics as diagnostic
indicators.

We present this clinician audit as a complement to, not a replacement for, parser and
verifier validation. The full package now separates upstream parser reliability,
verifier scoring reliability and shift, same-answer process metrics, and clinician
directional validation.


\responseheading{Revision Details}
\textbf{Evaluation metrics: diagnostic framing and clinician-validation boundary.}
\begin{displayquote}
We evaluate Fa-DPO with three groups of metrics: final-answer accuracy,
\deleted{process faithfulness}\added{automatic process metrics}, and ARU
unfaithfulness shift by reasoning function.

\deleted{These automatic process metrics provide mechanism-oriented proxy evidence rather than direct measures of absolute clinical correctness.}

\added{Because these process metrics are computed by the same ARU parser and NLI verifier used to construct the Fa-DPO risk signal, they are interpreted as diagnostic indicators of evidence-use and consistency behavior rather than independent clinical ground truth or direct measures of absolute clinical correctness. The same-answer clinician audit in Appendix~\ref*{ms-app:same_answer_clinician_details} provides directional validation of these indicators, not a replacement for clinician-adjudicated evaluation.}
\end{displayquote}

\textbf{Main-text evidence validation: how the clinician audit is used.}
\begin{displayquote}
\deleted{We next assess whether the evidence used in the preceding analyses is reliable.}

\added{This section validates the process evidence used in the preceding analyses. It asks three linked questions: whether automatic same-answer improvements align with clinician judgments and data-audit evidence; whether the ARU parser, routed NLI verifier, and token masks are reliable enough for proxy process measurement; and whether the same signal remains directionally useful on natural outputs, answer-changing diagnostics, and retrospective safety checks. These experiments support bounded use of the process metrics, not prospective clinical validation.}

{\color{blue}Same-answer clinician audit \& 100 paired questions; 200 anonymized responses \& Clinicians prefer Fa-DPO on all five process criteria, with non-tie win rates of 77.3--80.6\%; overall clinically coherent reasoning rises from 70.0\% to 84.0\%. \& Directional validation under fixed answer correctness, not full clinician-adjudicated evaluation of all outputs.}

\end{displayquote}

\textbf{Appendix clinician validation: protocol, paired results, and agreement.}
\begin{displayquote}
\added{Clinician Validation of Same-Answer Process Metrics}

\added{To validate the same-answer process metrics, we test whether their direction matches clinician judgments when answer correctness is fixed. The blinded audit uses the representative II-Medical-8B same-answer hard subset, with 100 paired MedCaseReasoning questions and 200 anonymized responses. Two clinicians independently review all paired cases. For each case, the Standard DPO-Retained-34K and Fa-DPO-Retained-34K responses are randomly assigned to anonymous A/B slots; model identities, training condition, ARU/NLI labels, verifier scores, and automatic metric values are hidden from the clinicians. Clinicians compare the two outputs on evidence support, contradiction avoidance, unsupported reasoning avoidance, reasoning-answer consistency, and overall process reliability, choosing A better, B better, tie, or cannot judge for each criterion. Disagreements are recorded for agreement statistics before resolution, and final reported pairwise counts use the resolved label after discussion or third-clinician adjudication.}

\added{Clinicians prefer Fa-DPO on all five same-answer criteria, with non-tie win rates of 77.3--80.6\%. Response-level pass rates move in the same direction: the overall clinically coherent reasoning rate increases from 70.0\% for Standard DPO-Retained-34K to 84.0\% for Fa-DPO-Retained-34K.}

\added{Cohen's $\kappa$ ranges from 0.60 to 0.68 across criteria, exact agreement ranges from 76.0\% to 80.0\%, and automatic-human directional agreement ranges from 70.0\% to 78.0\%. These values support the direction of the automatic process metrics while reinforcing that the automatic metrics remain diagnostic indicators rather than clinician-adjudicated ground truth.}
\end{displayquote}

\textbf{Appendix same-answer clinician-validation tables.}
\begin{displayquote}
\begin{center}
\footnotesize
\captionof{table}{\added{Clinician validation of same-answer process reliability on the representative II-Medical-8B pair. The audit includes 100 paired questions with final-answer correctness fixed for both systems; pairwise counts are Fa-DPO-Retained-34K better / tie / Standard DPO-Retained-34K better.}}
\begin{adjustbox}{max width=\linewidth}
\begingroup\color{blue}
\begin{tabular}{lccc}
\toprule
\textbf{Criterion} & \textbf{Pairwise count} & \textbf{Non-tie Fa-DPO win (\%)} & \textbf{95\% CI} \\
\midrule
Evidence support & 59 / 24 / 17 & 77.6 & [66.0, 87.0] \\
Contradiction avoidance & 54 / 33 / 13 & 80.6 & [68.2, 91.0] \\
Unsupported reasoning avoidance & 62 / 22 / 16 & 79.5 & [68.0, 89.0] \\
Reasoning-answer consistency & 58 / 25 / 17 & 77.3 & [65.4, 87.2] \\
Overall process reliability & 60 / 23 / 17 & 77.9 & [66.7, 88.0] \\
\bottomrule
\end{tabular}
\endgroup
\end{adjustbox}

\vspace{0.6em}
\captionof{table}{\added{Clinician agreement and automatic-human directional alignment in the same-answer validation audit.}}
\begin{adjustbox}{max width=\linewidth}
\begingroup\color{blue}
\begin{tabular}{lccc}
\toprule
\textbf{Criterion} & \textbf{Cohen's $\kappa$} & \textbf{Exact agreement} & \textbf{Auto-human directional agreement} \\
\midrule
Evidence support & 0.66 & 78.0\% & 74.0\% \\
Contradiction avoidance & 0.60 & 76.0\% & 70.0\% \\
Unsupported reasoning avoidance & 0.68 & 80.0\% & 78.0\% \\
Reasoning-answer consistency & 0.63 & 77.0\% & 73.0\% \\
Overall process reliability & 0.65 & 79.0\% & 75.0\% \\
\bottomrule
\end{tabular}
\endgroup
\end{adjustbox}
\end{center}
\end{displayquote}

\newpage
\begin{commentbox}{R2.6 Specific Comment: Mapped/Unmapped Natural-Error Subsets}
Report Fa-DPO vs Standard DPO performance on mapped and unmapped natural-error subsets,
analyze transfer, and discuss that the taxonomy covers less than half of observed
external errors.
\end{commentbox}

\responseheading{Response and Revision Summary}

We thank the reviewer for requesting separate mapped/unmapped subset results. Consistent with R2.2, we replaced the construct-mismatched Med-HALT pilot with the task-aligned held-out natural-generation analysis on 3,338 evidence-grounded MQA cases.

As clarified in R2.2, official answers are used only after generation to determine correctness and annotate response-level error categories; they are not used for generation, training, candidate retention, hyperparameter tuning, or model selection.

The main finding is on mapped natural errors. Fa-DPO reduces any mapped F1--F4 error from 70.00\% to 62.40\% and lowers unsupported medical claims from 5.20\% to 3.92\%. Among cases where Standard DPO has a mapped F1--F4 error, Fa-DPO produces a correct response without a mapped error in 27.4\% of cases, although 66.8\% still retain a mapped error. This indicates partial transfer to mapped natural errors, not complete error removal.

The unmapped subset is too small for a strong improvement claim. Unmapped errors change from 1.10\% to 0.96\%, and the Standard-DPO unmapped stratum contains only 37 cases. We therefore use this subset as a boundary check showing no material increase in taxonomy-outside errors, not as evidence that Fa-DPO improves all errors outside F1--F4.

We added the mapped/unmapped subset analysis and clarified that mapped results support bounded transfer, while unmapped results are boundary evidence only.


\responseheading{Revision Details}
\textbf{Appendix protocol: remove the construct-mismatched Med-HALT mapped/unmapped pilot.}
\begin{displayquote}
\deleted{Pilot Natural-Error Coverage Protocol on Med-HALT.}

\deleted{This appendix records the pilot protocol used for the Med-HALT study. The goal is not to compare Fa-DPO against Standard DPO under matched conditions, but to ask a narrower external-validity question: whether sampled errors from an external medical reasoning benchmark can be described using the controlled F1--F4 taxonomy introduced in Sec. 3.}

\deleted{For each sampled response, the annotation sheet records whether the response contains any error that can be reasonably mapped to the controlled F1--F4 taxonomy. Responses whose error pattern does not fit the controlled taxonomy are treated as unmapped external errors rather than forced into one of the four controlled types.}

\added{The evidence also shows bounded task-family generalization. The matched benchmark results cover multiple evidence-grounded MQA datasets and general-instruction, medical-adapted, and reasoning-enhanced backbones.}
\end{displayquote}

\textbf{Held-out task-aligned natural-error analysis.}
\begin{displayquote}
\added{To test transfer beyond controlled counterfactual rewrites, we evaluate the representative II-Medical-8B matched pair on 3,338 model-sampled cases from MedXpertQA-test, MedBullets-op4, MedBullets-op5, and MedQA-test. This analysis compares Standard DPO-Retained-34K and Fa-DPO-Retained-34K responses on the same held-out prompts without local rewriting. The labels are automatic rather than double-clinician-adjudicated and may co-occur, so this analysis supports a fixed-protocol consistency check rather than definitive clinical ground truth.}

\added{The held-out natural-error analysis is conducted as a frozen post hoc evaluation over already generated model responses. Official answers are used only after generation to determine answer correctness and assign response-level error labels for analysis; they are not included in prompts, training data, candidate retention, hyperparameter tuning, or model selection. Each row contains one generated response with the prompt, answer options, gold option, predicted option, and visible reasoning. The label definitions follow the evidence-grounded F1--F4 schema but are applied to natural outputs. F1 marks omission or substantive weakening of clinically important question evidence, F2 marks incorrect medical interpretation of evidence that the response uses, F3 marks overstatement beyond what the available facts support, and F4 marks a selected answer that is unsupported by or inconsistent with the preceding reasoning. An unmapped error is a substantive medical reasoning or knowledge error not cleanly captured by F1--F4; unsupported medical claim and harmful recommendation are recorded as separate safety-relevant flags. The labels are automatic rather than double-clinician-adjudicated, and they are multi-label: F1--F4 may co-occur with one another and with unmapped, unsupported-claim, or harmful-recommendation flags.}

\added{Together, Tables~\ref*{ms-tab:heldout_natural_error_analysis} and~\ref*{ms-tab:heldout_natural_standard_strata} show bounded transfer to natural model generations. Fa-DPO increases micro accuracy by 1.75 points, reduces any mapped F1--F4 error by 7.60 points, and lowers unsupported medical claims from 5.20\% to 3.92\%. The largest mapped reduction is F4 reasoning-answer mismatch. In the Standard-DPO mapped-error stratum, Fa-DPO is correct without a mapped error in 27.4\% of cases, but still has a mapped error in 66.8\%. Thus, the effect is meaningful but far from complete error removal, and the small unmapped stratum remains boundary evidence rather than a stable estimate of broad natural-error coverage.}
\end{displayquote}

\textbf{Discussion: bounded mapped/unmapped interpretation.}
\begin{displayquote}
\added{The held-out natural-error analysis on 3,338 model-sampled MQA cases further shows fewer mapped F1--F4 errors and fewer unsupported medical claims for the representative II-Medical-8B pair. These findings indicate transfer beyond controlled rewrites within the same task family, while keeping the scope limited: F1--F4 is an evidence-use taxonomy for this setting, not a complete account of all medical hallucination or clinical reasoning errors.}
\end{displayquote}

\newpage
\begin{commentbox}{R2.7 Specific Comment: Parser Accuracy Thresholds}
Report parser accuracy on 100 responses with clinician-annotated ARU boundaries and
role labels. Minimal acceptable metrics include span boundary F1 ($\geq 0.85$) and role
classification accuracy ($\geq 0.80$), or state this is an unmeasured error source.
\end{commentbox}

\responseheading{Response and Revision Summary}
We completed the requested 100-response clinician-adjudicated parser validation. As detailed in R2.3, this validation uses clinician-adjudicated ARU boundaries and role labels; here we focus on the reviewer-specified threshold check.

At the primary token-level IoU $\geq 0.5$ threshold, the parser achieves span boundary F1 0.863 and role classification accuracy 0.802, meeting the requested targets of span F1 $\geq 0.85$ and role accuracy $\geq 0.80$.

We therefore added the threshold-based parser-validation result. The full validation protocol, sensitivity analysis, and interpretation limits are reported in R2.3.

\responseheading{Revision Details}
\textbf{Appendix~E: report the requested threshold check.}
\begin{displayquote}
\added{After adjudication, the parser reaches span F1 0.863 and role accuracy 0.802 at IoU $\geq 0.5$, meeting the reviewer-specified targets of span F1 $\geq 0.85$ and role accuracy $\geq 0.80$.}
\end{displayquote}

\newpage
\begin{commentbox}{R2.8 Specific Comment: Answer Preservation for F1--F3}
F1--F3 negatives preserve the final answer, creating correct-answer wrong-reasoning
pairs. The authors should discuss whether answer-accuracy gains may come from learning
to ignore reasoning and rely on spurious question-answer correlations.
\end{commentbox}

\responseheading{Response and Revision Summary}
We thank the reviewer for identifying this shortcut risk. The central design point is
that F1--F3 keep the final answer fixed: the preferred and rejected responses have the
same answer, so final-answer correctness cannot distinguish $y_w$ from $y_l$ and the
pair contrast falls on evidence use in the reasoning trace.

We revised the manuscript to state this explicitly. We also added two follow-up checks:
answer-changing evaluations, where final-answer preservation no longer holds, and a
token-weight decomposition showing that the effective rejected-side training weight is
concentrated mainly on reasoning tokens rather than final-answer tokens. These checks
weaken an answer-only shortcut explanation, but we state the conclusion narrowly:
residual answer-level or latent-reasoning shortcuts remain possible.


\responseheading{Revision Details}
\textit{1. Controlled negative-construction subsection: original wording and revised wording.}
\begin{displayquote}
\deleted{For F1--F3, we keep the final answer unchanged to construct correct-answer, wrong-reasoning negatives.}

\added{For F1--F3, we intentionally keep the final answer unchanged to construct correct-answer, wrong-reasoning negatives. This removes final-answer correctness as the distinguishing feature between the preferred and rejected responses and concentrates the preference contrast on evidence-use failures in the reasoning trace.}
\end{displayquote}

\textit{2. Evidence-validation section and appendix: no corresponding diagnostic section was present in the original manuscript. The revised manuscript adds the following answer-shortcut diagnostic summary in the main text and full diagnostic tables in the appendix.}
\begin{displayquote}
\added{F1--F3 preserve the final answer by design: the retained preferred and rejected responses have the same extracted final answer, so final-answer correctness cannot distinguish $y_w$ from $y_l$ in these pairs. We first verify this construction property in the retained pool. To assess the remaining shortcut risk, we add two complementary checks on the representative II-Medical-8B Standard DPO-Retained-34K and Fa-DPO-Retained-34K pair. The answer-changing check evaluates settings in which final-answer preservation does not hold: an F4-type model-output subset from held-out MedCaseReasoning test outputs, an answer-changing subset in which edited evidence changes the gold final answer, and a counterfactual answer test with minimally edited key clinical evidence. These subset labels are fixed before computing the reported metrics. The token-weight check decomposes Fa-DPO's effective rejected-side token weight into reasoning-token and final-answer-token mass using the training tokenizer and stored token-risk masks. Final-answer tokens are tokenizer-aligned tokens inside the extracted final-answer span appended as the final claim ARU; reasoning tokens are the remaining rejected-response tokens.}

\added{Together, these checks weaken an answer-only shortcut explanation for the observed gains.} Table~\ref*{ms-tab:answer_preservation_check} \added{confirms the construction property: F1--F3 preserve the extracted final answer in all retained pairs, while F4 contains 1,340 answer-changing or answer-missing retained pairs out of 8,426 F4 pairs.} Table~\ref*{ms-tab:answer_changing_diagnostic} \added{shows that Fa-DPO also improves process metrics when final-answer preservation does not hold.} Table~\ref*{ms-tab:counterfactual_answer_test} \added{adds a harder answer-changing test under minimally edited clinical evidence; the answer-change and stickiness intervals are close to zero, so these two rows are suggestive rather than definitive.} Table~\ref*{ms-tab:answer_token_weight_decomposition} \added{shows that the effective training signal is concentrated mainly on reasoning tokens under the stated tokenizer alignment. The remaining limit is that these checks do not prove that answer-level shortcuts are absent in all answer-changing settings or that the model's latent reasoning has been fully corrected.}
\end{displayquote}

\textit{3. Conclusion and limitations: original interpretation and revised caution.}
\begin{displayquote}
\deleted{The representative same-answer analyses, blind pairwise evaluation on the II-Medical-8B same-answer hard subset, and score-stratified correlation with expert ratings further suggest that the observed gains are associated with more faithful reasoning rather than answer-only shortcuts.}

\added{Control analyses indicate that these gains are better explained by improved reasoning-process behavior than by positive-only SFT, vanilla DPO, sample-count effects, reranking, or answer-only shortcuts.}

\added{The added audits are mostly retrospective and partly automatic; they do not establish broad hallucination detection or deployment safety, and the retained negatives do not cover the full distribution of naturally generated medical-LLM errors.}
\end{displayquote}

\newpage
\begin{commentbox}{R2.9 Specific Comment: Equation 11 Normalization}
Clarify whether $T_l / \sum_t m_t$ preserves relative log-ratio magnitudes across
responses with different token lengths and mask sums, or whether it compresses
differences between high-risk and low-risk negatives.
\end{commentbox}

\responseheading{Response and Revision Summary}
We thank the reviewer for asking us to clarify what the Eq.~(11) normalization preserves. This point is addressed as part of the broader DPO objective-presentation revision package described in R1.3--R1.5.

We now state explicitly that the factor $T_l/\sum_t m_t$ preserves the total effective token mass within each rejected response and preserves the within-response relative emphasis induced by the ARU-risk mask. It does not preserve the absolute ordering of raw response-level log-ratio magnitudes across different rejected responses.

Accordingly, Eq.~(11) is now described as a local credit-assignment normalization: it redistributes a fixed rejected-side token mass toward high-risk spans, rather than increasing the total rejected-response weight. We also added a normalization diagnostic comparing uniform, unnormalized, normalized, and capped weighting variants. The diagnostic supports the normalized version as a stable way to retain high-risk emphasis without the instability observed when normalization is removed.

\responseheading{Revision Details}
We clarified the preservation property of Eq.~(11), explicitly stated that cross-response raw log-ratio ordering is not preserved, and added the normalization diagnostic to support the chosen weighting scheme.

\newpage
\begin{commentbox}{R2.10 Specific Comment: Confidence Intervals in Table 2}
Table 2 reports point estimates for paired deltas without confidence intervals or
standard errors. Add bootstrap or paired t-test confidence intervals, especially for
subsets with $n=68$ to $100$.
\end{commentbox}

\responseheading{Response and Revision Summary}
We agree that Table~2 should not report only point-estimate paired deltas. This
comment is addressed by the same statistical revision described in R1.13, so we
summarize only the Table~2-specific changes here.

In the revised manuscript, Table~2 now reports retained cases, response and ARU
denominators, paired deltas, and 95\% confidence intervals. The bootstrap resamples
cases rather than ARUs, and the Holm-adjusted p-value screen is moved to the appendix.


\responseheading{Revision Details}
\textbf{Table~2 caption and protocol: add paired confidence intervals and resampling unit.}
\begin{displayquote}
\deleted{Here, $n$ denotes the retained subset size for each pair.}

\added{Here, $n$ denotes retained same-answer cases, responses denotes $2n$, and ARUs $N$ is the total ARU measurement count, not the resampling unit. Supplementary multiplicity-adjusted p-value screening is reported in Appendix~\ref*{ms-app:process_metrics}.}

\added{For the cross-pair summary in Table~\ref*{ms-tab:paired_process_summary}, we report the case-weighted mean paired differences between Fa-DPO-Retained-34K and Standard DPO-Retained-34K for all four metrics. Confidence intervals for these paired process-metric deltas are computed with paired cluster bootstrap over cases: each bootstrap replicate resamples case identifiers, keeps both system responses and all ARUs within each selected case together, recomputes $\Delta = M_{\text{Fa-DPO-Retained-34K}} - M_{\text{Standard DPO-Retained-34K}}$, and reports the 2.5 and 97.5 percentiles over 10,000 replicates. ARU counts affect the within-response metric denominators, but each retained case contributes one metric value per system to the cross-case mean.}
\end{displayquote}

\newpage
\tocsection{Reviewer \#3}

\begin{commentbox}{R3 General Assessment}
Reviewer \#3 states that the manuscript has solid methodological strengths but should
address abstract improvement, alignment between objectives and local faithfulness
definition, and conciseness before final acceptance.
\end{commentbox}

\responseheading{Response and Revision Summary}
We thank the reviewer for the constructive assessment and for identifying the main
presentation issues that needed improvement before acceptance. We revised the Abstract
to foreground the main process-reliability and downstream accuracy results, aligned
the research objectives with the local ARU-level operationalization of faithfulness,
defined ARU at first use, and shortened weakly related background material. These
changes were coordinated with the conceptual revisions made in response to R1.1 and
R3.2, so that the manuscript now distinguishes global clinical reasoning quality from
the local diagnostic process indicators used by Fa-DPO.

\newpage
\begin{commentbox}{R3.1 Abstract Improvement}
The Abstract should incorporate process-faithfulness metrics and downstream accuracy
results. Dataset format and model configurations are routine and comparatively
redundant.
\end{commentbox}

\responseheading{Response and Revision Summary}
We thank the reviewer for this helpful suggestion. We agree that the Abstract should foreground the main empirical findings rather than routine dataset-format or model-configuration details. We therefore rewrote the Abstract to be shorter and more result-focused.

In the revised Abstract, we now report both automatic same-answer process metrics and downstream MQA accuracy gains. Specifically, Fa-DPO improves Support Rate by 5.1 points, reduces Unsupported Medical Reasoning Rate by 7.4 points, reduces Reasoning--Answer Inconsistency Rate by 5.6 points, and yields observed downstream MQA macro-average accuracy gains of 1.2--1.6 percentage points over retained-pair standard DPO. We also keep the process-metric language conservative by describing these metrics as diagnostic indicators, consistent with the revisions made for R2.1.


We revised the Abstract by shortening routine dataset and model-configuration descriptions, adding the main process-metric and downstream accuracy results, and clarifying that automatic process metrics are diagnostic evidence rather than clinician-adjudicated ground truth.


\responseheading{Revision Details}
\textbf{Abstract: concise problem--method--result rewrite.}
\begin{displayquote}
\deleted{However, DPO's sequence-level objective will mask localized faithfulness errors in long, mostly correct clinical responses.}\added{However, DPO's sequence-level objective can dilute the suppressive signal from localized evidence-use errors in long, mostly correct clinical responses.}

\deleted{In this paper, we propose \textbf{Fa-DPO}, a \textbf{F}aithfulness-\textbf{A}ware \textbf{D}irect \textbf{P}reference \textbf{O}ptimization framework that converts fine-grained faithfulness risk into preference signals for evidence-grounded MQA.}\added{In this paper, we propose \textbf{Fa-DPO}, a \textbf{F}aithfulness-\textbf{A}ware \textbf{D}irect \textbf{P}reference \textbf{O}ptimization framework for evidence-grounded MQA.}

\added{Fa-DPO identifies localized evidence-use and reasoning--answer consistency failures in long clinical reasoning chains and uses them to strengthen learning signals for risky reasoning spans and high-risk responses.}

\deleted{We then decompose each response into atomic reasoning units, verify each unit through role-specific evidence routing and natural language inference (NLI)-based checking, and use the resulting risk scores to reweight high-risk spans and set risk-adaptive margins.}\added{Responses are decomposed into reasoning-function-labeled Atomic Reasoning Units (ARUs), verified via natural language inference (NLI), to guide risk-aware preference optimization.}

\deleted{Fa-DPO also improves downstream MQA accuracy over standard DPO by 1.2--1.6 macro-average points, suggesting that faithfulness-aware preference learning improves both reasoning quality and final-answer accuracy.}\added{Experiments on 34,204 controlled preference pairs across six 8B backbones show that Fa-DPO improves automatic same-answer process metrics (+5.1 Support Rate; -7.4 Unsupported Medical Reasoning Rate; -5.6 Reasoning--Answer Inconsistency Rate) and yields observed downstream MQA macro-average accuracy gains of 1.2--1.6 percentage points over retained-pair standard DPO. These process-metric gains are interpreted as diagnostic evidence, while the benchmark results separately support final-answer accuracy improvements.}
\end{displayquote}

\newpage
\begin{commentbox}{R3.2 Alignment Between Objectives and Local Faithfulness}
In Sec. 2, Objective 1 appears global while Objectives 2 and 3 are local. The
manuscript proposes a local approach to a global problem, which does not align well
with the local faithfulness definition.
\end{commentbox}

\responseheading{Response and Revision Summary}
We thank the reviewer for identifying the mismatch between the global objective and
the local operational definition. We agree that the original Sec.~2 could be read as
equating local ARU-level metrics with the full global concept of clinical reasoning
reliability. That was not our intended claim.

We substantially revised Sec.~2 to separate the three levels of the argument. First,
we define the global target as evidence-grounded process reliability in medical QA.
Second, we define evidence faithfulness narrowly as whether generated clinical
reasoning units are supported by, and do not contradict, the provided clinical
evidence. Third, we treat reasoning--answer consistency as a complementary
process-reliability dimension rather than as part of narrow evidence faithfulness.

We also clarified the relation between the global and local levels. Because global
clinical reasoning reliability is difficult to supervise reliably with a single
response-level score, Fa-DPO operationalizes it through local ARU-level checks. These
checks provide scalable training signals and diagnostic process evidence, but they
are not presented as complete clinical ground truth for global reasoning quality.


\responseheading{Revision Details}
\textbf{Sec.~2: align global goal with local ARU-level operationalization.}
\begin{displayquote}
\deleted{We define faithfulness as the extent to which intermediate clinical inferences are supported by retrieved evidence, avoid contradictions, and remain consistent with the final answer.}

\added{Following this general view, we define evidence faithfulness as the extent to which intermediate clinical inferences are supported by retrieved evidence and do not contradict it. We further evaluate reasoning-answer consistency as a related diagnostic check, because the final answer should follow from the preceding reasoning.}

To address this problem, we study how to identify, measure, and optimize local \deleted{faithfulness}\added{faithfulness and consistency} failures in clinical reasoning.

\deleted{Measure faithfulness risk at the atomic reasoning unit (ARU) level by routing each unit to a role-specific verification premise and checking its support status.}\added{Measure ARU-level verification scores by routing each unit to a reasoning-function-specific verification premise and checking its support or consistency status.}

\added{As global clinical reasoning reliability is difficult to supervise reliably with a single response-level label, we use ARU-level verification as a local operationalization that provides scalable local training signals and diagnostic process evidence, rather than as a complete measure of global clinical reasoning quality.}
\end{displayquote}

\textbf{Method-stage naming: local risk scoring.}
\begin{displayquote}
\deleted{ARU-Level Faithfulness Risk Measurement}\added{ARU-Level Risk Scoring}

\deleted{To provide fine-grained faithfulness risk signals for downstream DPO training, we decompose each quality-controlled candidate negative response into ARUs and assign each ARU a verification score that operationalizes faithfulness risk through multi-source evidence routing and NLI-based verification.}

\added{To provide fine-grained risk-scoring signals for downstream DPO training, we decompose each quality-controlled candidate negative response into ARUs and assign each ARU a verifier-derived score through multi-source evidence routing and NLI-based verification. For evidence-routed ARUs, the score estimates evidence-use risk; for claim ARUs, including the final-answer claim, it estimates reasoning--answer consistency risk.}
\end{displayquote}

\textbf{Sec.~6.1.3: proxy-metric framing.}
\begin{displayquote}
\deleted{Support Rate, Contradiction Rate, and Unsupported Medical Reasoning Rate are evidence-faithfulness proxy metrics, while Reasoning-Answer Inconsistency Rate is a reasoning-answer consistency metric.}\added{Support Rate and Unsupported Medical Reasoning Rate are evidence-support indicators, Contradiction Rate is an evidence-conflict indicator, and Reasoning-Answer Inconsistency Rate is a reasoning-answer consistency metric.} Taken together, these metrics are intended as diagnostic mechanism evidence rather than direct measures of absolute clinical correctness.
\end{displayquote}

\textbf{Conclusion and limitations: local proxies are not complete clinical truth.}
\begin{displayquote}
\added{Across six paired 8B base models, Fa-DPO improves answer accuracy and automatic reasoning-process metrics over Standard DPO-Retained-34K under the same training setup, while remaining competitive with strong released baselines.}

\added{The added audits are mostly retrospective and partly automatic; they do not establish broad hallucination detection or deployment safety, and the retained negatives do not cover the full distribution of naturally generated medical-LLM errors.}
\end{displayquote}

\newpage
\begin{commentbox}{R3.3 Conciseness}
The manuscript requires linguistic refinement and should avoid weakly relevant content
and redundant descriptions of established common knowledge. Shorter is better.
\end{commentbox}

\responseheading{Response and Revision Summary}
We agree and revised the manuscript for concision. We removed weakly related background
and tightened the Related Work discussion around the distinctions that directly define
Fa-DPO: evidence access versus evidence use, post-hoc verification versus training
signal, and sequence-level DPO versus local risk-aware DPO.

We also replaced generic method-positioning language with the concrete mechanism:
ARU-level verification scores are used to retain rejected responses, weight rejected
spans, and set sample-level margins.


\responseheading{Revision Details}
\textbf{Sec.~3.2: preference-learning gap.}
\begin{displayquote}
\deleted{The remaining gap is risk calibration: existing fine-grained methods improve process supervision or self-correction, but they do not calibrate local faithfulness risk inside offline DPO.}
\added{Fine-grained, context-faithful, and hallucination-aware variants move preference signals toward tokens, steps, context conflicts, or response-level hallucination labels~\cite{gu2025maskdpo, zhou2025treg, yang2025sepo, xu2025fullstepdpo, bi2025contextdpo, zhao2024hadpo}. These methods show that preference learning can target factual or contextual errors, but they do not use reasoning-function-labeled medical units to decide which rejected spans and responses should receive stronger DPO penalties.}
\end{displayquote}

\textbf{Sec.~3.2: Fa-DPO mechanism relative to DPO variants.}
\begin{displayquote}
\deleted{Our method follows the offline preference-learning route, but differs from standard and fine-grained DPO by applying ARU-level evidence checking to retain rejected responses, weight spans, and calibrate the penalty assigned to each rejected response.}
\added{Fa-DPO fills this gap by using ARU-level verification scores to retain high-risk negatives, weight rejected spans, and set risk-adaptive margins in offline DPO.}
\end{displayquote}

\textbf{Sec.~3.3: from post-hoc verification to training signal.}
\begin{displayquote}
\added{Self-checking, verifier-RAG, atomic factuality evaluation, and biomedical RAG checkers decompose outputs into claims and verify them against generated, retrieved, or external evidence~\cite{manakul2023selfcheckgpt, phasook2025vereafine, min-etal-2023-factscore, wei2024longform, song-etal-2024-veriscore, ji2026medragchecker}.}

\added{Fa-DPO adopts the decompose-then-verify principle but changes its role: ARUs encode clinical reasoning function, evidence route, local verification score, and token span, allowing verification results to become DPO weights and margins rather than only diagnostic labels.}
\end{displayquote}

\textbf{Sec.~3.3: concise motivation.}
\begin{displayquote}
\deleted{Our method differs from these evaluation and verification studies by converting fine-grained faithfulness measurements into training-time preference supervision, so that evidence-use failures become optimization signals rather than only diagnostic outputs.}
\added{Existing verification signals are usually used for evaluation, filtering, consensus checking, or post-hoc diagnosis.}

\deleted{Taken together, these gaps motivate our method, which targets the unresolved bottleneck of faithful clinical reasoning in evidence-grounded MQA by converting fine-grained faithfulness risk into stable preference learning signals.}
\added{Fa-DPO addresses this bottleneck by moving from evidence access to evidence use, from post-hoc faithfulness checking to training-time optimization, and from sequence-level preference learning to risk-adaptive span-aware preference learning.}
\end{displayquote}

\newpage
\begin{commentbox}{R3.4 Define ARU in the Abstract}
Define the acronym ARU at its first occurrence in the Abstract rather than later in the
manuscript.
\end{commentbox}

\responseheading{Response and Revision Summary}
We agree and have revised the Abstract accordingly. The first occurrence of the term
now defines the acronym as \textbf{Atomic Reasoning Units (ARUs)} before the later
Abstract phrase ARU-level verification scores uses the short form.


\responseheading{Revision Details}
\textbf{Abstract: define ARU at first mention.}
\begin{displayquote}
\deleted{atomic reasoning units}
\added{reasoning-function-labeled Atomic Reasoning Units (ARUs)}
\end{displayquote}

\newpage
\begin{commentbox}{R3.5 General Faithfulness Definitions}
In Section 2, first review and cite widely recognized general definitions of
faithfulness, then elaborate on the specific definition adopted in this work.
\end{commentbox}

\responseheading{Response and Revision Summary}
We agree that Sec.~2 should first anchor the terminology in the general
faithfulness literature before introducing our task-specific definition. In the
revision, the beginning of Sec.~2 now cites prior grounded-generation and QA/RAG
evaluation work that treats faithfulness as consistency between generated content and
the provided source or evidence, often operationalized through support and
contradiction judgments.

Following this convention, we define \emph{evidence faithfulness} in our setting as
the relation between each generated clinical reasoning unit and the available
clinical evidence: a unit is evidence-faithful if it is supported by the evidence and
does not contradict it. We then separately define reasoning--answer consistency as a
complementary process-reliability dimension, which directly addresses the reviewer's
concern that final-answer consistency should not be folded into the narrow
faithfulness definition.


\responseheading{Revision Details}
\noindent\textbf{Sec.~2, Research Objectives}
\begin{displayquote}
\added{In grounded text generation and QA/RAG evaluation, faithfulness is commonly understood as consistency between generated content and its input source, document context, or retrieved evidence, and is often assessed through whether generated claims are supported by, or contradicted by, the source~\cite{maynez-etal-2020-faithfulness,adlakha2024evaluating,saadfalcon2024ares}.}

\deleted{We define faithfulness as the extent to which intermediate clinical inferences are supported by retrieved evidence, avoid contradictions, and remain consistent with the final answer.}

\added{Following this general view, we define evidence faithfulness as the extent to which intermediate clinical inferences are supported by retrieved evidence and do not contradict it. We further evaluate reasoning-answer consistency as a related diagnostic check, because the final answer should follow from the preceding reasoning.}
\end{displayquote}

\newpage
\begin{commentbox}{R3.6 Reduce RAG Details}
Since the proposed method does not improve retrieval, the detailed introduction of RAG
in Section 3 may be redundant.
\end{commentbox}

\responseheading{Response and Revision Summary}
We agree with the reviewer that a detailed RAG introduction is unnecessary because
Fa-DPO does not modify retrieval. We revised Sec.~3 to remove the broad RAG overview
and keep one scope paragraph: biomedical IR and medical RAG provide candidate clinical
premises, but Fa-DPO focuses on whether generated reasoning uses those premises
faithfully.

We also clarified that retrieval is a verifier-side resource for offline negative
construction, ARU-level verification, and process-metric analysis. During benchmark
inference, the policy input is the question context, not retrieved passages. Thus,
``evidence-grounded'' refers to the training, verification, and analysis signals used
in this work, not inference-time RAG.

\responseheading{Revision Details}
\noindent\textbf{Sec.~5.1, Task Formulation}
\begin{displayquote}
\deleted{Our task is evidence-grounded medical question answering with reasoning traces. Given a medical question context $x$ and a medical evidence corpus $\mathcal{K}$, the model generates a response $y=c\oplus a$, where $c$ is the reasoning trace and $a$ is the final answer. We use $G \subseteq \mathcal{K}$ to denote the relevant evidence snippets. The response is expected to provide a reasoning trace grounded in $x$ and $G$, and a final answer supported by that reasoning trace.}

\added{Our task is evidence-grounded medical question answering with reasoning traces under a verifier-side evidence protocol. At policy inference time, the model input is the medical question context $x$, and the policy generates a response $y=c\oplus a$, where $c$ is the reasoning trace and $a$ is the final answer. The external evidence corpus $\mathcal{K}$ and retrieved evidence snippets $G \subseteq \mathcal{K}$ are verifier-side resources: they are used for offline negative construction, ARU-level verification, and automatic process metrics, but are not provided to the policy model during benchmark inference. Thus, evidence-grounded in this work refers to the training, verification, and analysis signals used to evaluate and optimize evidence use, not to inference-time retrieval-augmented generation.}
\end{displayquote}

\noindent\textbf{Figure~1 caption}
\begin{displayquote}
\deleted{\textbf{Clinician-Audited Negative Response Construction:} we construct candidate negative responses through local rewriting}

\added{\textbf{Controlled Negative Construction with Clinician Audit:} we construct candidate negative responses through local rewriting}

\deleted{route the ARUs to reasoning-function-specific verification premises}

\added{route the ARUs to verifier-side, reasoning-function-specific verification premises}
\end{displayquote}

\noindent\textbf{Sec.~5.3 and Sec.~5.3.1, Offline Controlled Negative Construction}
\begin{displayquote}
\deleted{Controlled Negative Response Construction with Clinician Audit}

\added{Offline Controlled Negative Construction with Clinician Audit}

\deleted{Controlled Negative Response Construction}

\added{Offline Controlled Negative Construction}

\added{This is an offline negative-construction stage: the generator may use retrieved evidence $G$ to construct rejected-response candidates with controlled evidence-use failures, but this does not mean that the final trained policy receives retrieved evidence at benchmark inference.}
\end{displayquote}

\noindent\textbf{Appendix, Retrieval and Routed-Premise Protocol}
\begin{displayquote}
\added{This design keeps Fa-DPO and the Standard DPO-Retained-34K baseline under the same policy-input condition; retrieval is reserved for offline construction, verification, and analysis.}
\end{displayquote}

\noindent\textbf{Sec.~3, Evidence-Grounded QA and Biomedical Information Retrieval}
\begin{displayquote}
\deleted{Retrieval-centered methods address this grounding gap by improving access to external medical knowledge.}

\deleted{Biomedical IR models and benchmarks, such as MedCPT and MedEureka, retrieve relevant medical texts from search logs, publications, clinical notes, and other domain-specific sources.}

\deleted{Recent IP\&M work further shows that medical retrieval itself remains difficult because terminology variation can degrade embedding-based retrieval, while ontology-based concept normalization can improve medical vector search.}

\deleted{Medical RAG retrieval also moves toward multi-evidence supervision: instead of assuming that one query is supported by one document, MMRE constructs one-to-many clinical retrieval supervision to retrieve complementary evidence passages for medical RAG.}

\deleted{These retrieval advances supply the evidence base for downstream reasoning.}

\deleted{Medical RAG and retrieval-reasoning systems then use this evidence for answer generation, graph-based reasoning, rationale-guided retrieval, or agentic coordination.}

\deleted{Recent surveys also frame RAG and reasoning as a unified search-reasoning problem.}

\deleted{In this view, retrieval supplies missing premises, and reasoning organizes them for knowledge-intensive tasks.}

\deleted{Recent IP\&M work on evidence-aware RAG further studies the reconciliation between retrieved evidence and a model's internal knowledge through self-corrective reasoning.}

\deleted{Together, this line improves evidence access and search-reasoning coordination.}

\deleted{However, evidence access does not guarantee evidence use: a model may retrieve relevant passages, including medically relevant and complementary passages, but still draw unsupported or inconsistent intermediate conclusions from them.}

\deleted{Our method therefore differs from retrieval models and evidence-aware prompting methods; its focus is not retrieving more evidence, but converting local failures in using retrieved clinical evidence into preference-learning signals.}

\added{Biomedical IR and medical RAG address evidence access by retrieving medical texts or complementary passages from search logs, publications, clinical notes, ontology-normalized indexes, and multi-evidence corpora~\cite{jin2023medcpt, fan2025medeureka, deng2026medicalontologyir, li2026mmre, zakka2024almanac, xiong2024benchmarking, li2026evidencerag}. These methods provide candidate premises for generation or verification, yet a rationale may still omit key findings, misread clinical relations, overextend a passage, or fail to support the final answer. Fa-DPO therefore treats retrieved evidence as verifier-side support and focuses on evidence use: it converts local failures in using clinical evidence into preference-learning signals.}
\end{displayquote}

\newpage
\begin{commentbox}{R3.7 Reverse Entailment Probability}
Question: In Sec. 5.4.2, should the reverse entailment probability be taken into
consideration?
\end{commentbox}

\responseheading{Response and Revision Summary}
We thank the reviewer for raising this useful directionality question. We agree that
the entailment direction should be explicit. The main verifier uses the forward
direction, evidence or routed premise $\rightarrow$ ARU, because the key question is
whether a generated reasoning unit is supported by available evidence or prior
reasoning.

Reverse entailment, ARU $\rightarrow$ evidence, asks a different question: whether the
generated unit can recover the evidence or premise. That is closer to coverage or
specificity than support. A concise but faithful clinical summary should not be
penalized simply because it does not entail every detail in a longer evidence span.

We revised Sec.~5.4.2 to make this design explicit. The primary risk score, rejected
response selection, token weighting, and adaptive margin use only forward entailment.
We also added a compact appendix definition of reverse entailment as an optional
auxiliary diagnostic, not a training or risk-scoring component. We do not report
reverse-entailment results because no full diagnostic experiment was run for this
point.


\responseheading{Revision Details}
\textbf{Sec.~5.4.2: fix the verifier entailment direction.}
\begin{displayquote}
The input to the verifier is a reasoning-function-specific verification premise $p$ and an ARU hypothesis $v_i$; the output is a risk score $s_i \in [0,1]$, where larger values indicate stronger evidence of \deleted{unfaithfulness}\added{unsupported, contradictory, or internally inconsistent reasoning}.

\added{For evidence-routed ARUs, the score measures evidence-support risk against the patient context and retrieved medical evidence. For claim ARUs, including the final-answer claim, the score estimates whether the claim or final answer follows from the preceding reasoning under the verifier-side NLI setup.}

Let $P_{\text{ent}}(p,v)$, $P_{\text{neu}}(p,v)$, and $P_{\text{con}}(p,v)$ denote the verifier-predicted entailment, neutral, and contradiction probabilities when premise $p$ is used to verify hypothesis $v$.

\added{The NLI direction is fixed by this verification setup. The primary verifier uses forward entailment from the routed premise to the ARU, $e_i^{\rightarrow}=P_{\text{ent}}(p_i,v_i)$, because it asks whether the generated reasoning unit is supported by available evidence or prior reasoning. The reverse probability $e_i^{\leftarrow}=P_{\text{ent}}(v_i,p_i)$ asks whether the generated unit entails the premise, which is a coverage or specificity diagnostic rather than an evidence-support signal. A concise clinical summary can be faithful without entailing every detail in a larger evidence span; therefore, reverse entailment is not used for candidate retention, token weighting, the response-level risk score, or the risk-adaptive margin. In this work, reverse entailment is kept only as a scope note for possible coverage or specificity diagnostics; we do not report reverse-entailment results, and it is not used in training or the main analysis.}
\end{displayquote}

\newpage
\begin{commentbox}{R3.8 Metric Presentation Order}
In Section 6.1.3, Support Rate and Unsupported Medical Reasoning Rate appear close in
value. Adjusting their presentation order may help reader understanding.
\end{commentbox}

\responseheading{Response and Revision Summary}
We thank the reviewer for this helpful suggestion. We revised Sec.~6.1.3 to organize
the process metrics by conceptual dimension rather than listing them in the previous
order. Support Rate and Unsupported Medical Reasoning Rate are now presented together
under the evidence-support dimension, followed by Contradiction Rate under the
evidence-conflict dimension and Reasoning-Answer Inconsistency Rate under the reasoning-answer
consistency dimension.

We also clarified that Unsupported Medical Reasoning Rate is not simply the
complement of Support Rate. Support Rate measures the fraction of all typed ARUs
entailed by their verification premises, whereas Unsupported Medical Reasoning Rate
focuses specifically on clinically inferential warrant and differentiation ARUs whose
medical reasoning is not adequately supported. This revision keeps the two related
evidence-support metrics adjacent while preserving their separate diagnostic roles.


\responseheading{Revision Details}
\textbf{Sec.~6.1.3: reorder automatic metrics by conceptual dimension.}
\begin{displayquote}
\deleted{We evaluate Fa-DPO with three groups of metrics: final-answer accuracy, process faithfulness, and role-wise unfaithfulness shift.}

\added{We evaluate Fa-DPO with three groups of metrics: final-answer accuracy, automatic process metrics, and unfaithfulness shift by ARU type.}

\deleted{We report two groups of verifier-derived process metrics}

\added{We report four ARU/NLI-based metrics, grouped by diagnostic role.}

\added{The evidence-support metrics are} \textbf{Support Rate} \added{and} \textbf{Unsupported Medical Reasoning Rate}.

\added{The latter is not the complement of Support Rate because it is computed only over clinically inferential warrant and differentiation units.}

\added{The evidence-conflict metric is} \textbf{Contradiction Rate}.
\end{displayquote}

\begin{displayquote}
\textit{Manuscript table-caption excerpt:} \added{The metrics are ordered by conceptual group: \textit{Support Rate} and \textit{Unsupported Medical Reasoning Rate} are evidence-support indicators, \textit{Contradiction Rate} is an evidence-conflict indicator, and \textit{Reasoning-Answer Inconsistency Rate} is a reasoning-answer consistency metric.}
\end{displayquote}

\newpage
\tocsection{Reviewer \#4}

\begin{commentbox}{R4 General Assessment}
Reviewer \#4 considers the paper credible and meaningful, with clear problem
formulation, logically connected modules, strong empirical design, a reusable dataset,
and practical significance.
\end{commentbox}

\responseheading{Response and Revision Summary}
We thank the reviewer for recognizing the manuscript's problem formulation,
methodological structure, empirical design, dataset contribution, and practical
significance. We revised the manuscript to address the three remaining concerns raised
by the reviewer: computational overhead, clinical-safety interpretation, and dataset
bias in the controlled preference-data construction. The changes are coordinated with
our responses to Reviewer \#1 and Reviewer \#2 on taxonomy coverage, diagnostic-indicator
limits, and external generalization.

\newpage
\begin{commentbox}{R4.1 Computational Complexity}
The multi-stage ARU decomposition and verification pipeline may introduce significant
overhead, which is not fully discussed from a deployment perspective.
\end{commentbox}

\responseheading{Response and Revision Summary}
We thank the reviewer for raising the deployment-overhead concern. We revised the
manuscript to make the computational boundary explicit: Fa-DPO's extra computation is
concentrated in one-time offline preference construction, while deployment remains
standard single-generation inference.

The revised manuscript reports archived offline verifier timing and a matched latency
comparison, but avoids unsupported end-to-end construction-time claims. After training,
Fa-DPO is served without online ARU parsing, evidence retrieval, or NLI scoring.


\responseheading{Revision Details}
\textbf{Sec.~7.4, Computational Cost and Deployment Overhead: offline cost boundary.}
\begin{displayquote}
\added{The full Fa-DPO pipeline has a one-time offline construction cost before the final model is trained.}

\added{In the archived construction run, ARU decomposition, evidence retrieval/reranking, NLI scoring, and preference-file preparation took 47.06 hours in total, dominated by ARU decomposition (39.60 hours), followed by evidence retrieval/reranking (4.32 hours) and NLI scoring (3.07 hours).}

\added{These stages produce the retained pairs, risk scores, and token masks before training; they are not executed for each user query.}

\added{During training, Fa-DPO uses the same 34,204 retained preference pairs as Standard DPO-Retained-34K and adds only precomputed rejected-token weights and a response-level margin to the DPO objective.}

\added{At deployment, Fa-DPO is served as a standard single-generation fine-tuned model.}

\added{Mean latency was 6.804 seconds for Standard DPO-Retained-34K (SD 0.035; 95\% CI: 6.797--6.812) and 6.808 seconds for Fa-DPO-Retained-34K (SD 0.047; 95\% CI: 6.798--6.817), and the paired mean difference was 0.004 seconds (95\% CI: -0.008--0.014).}

\added{Thus, Fa-DPO shifts verifier cost to offline preference construction while keeping inference to one generation without online ARU parsing, evidence retrieval, or NLI scoring.}
\end{displayquote}

\textbf{Table~\ref*{ms-tab:faithfulness_baseline_results} and Appendix~\ref*{ms-app:faithfulness_baseline_protocol}: verifier-reranking deployment boundary.}
\begin{displayquote}
\added{The table also records inference generations and whether deployment-time ARU/NLI verification is required; it does not report measured reranking latency for $K=4$ or $K=8$.}

\added{Verifier reranking narrows the gap, especially at $K=8$, but requires multiple generations and online ARU/NLI verification; Fa-DPO keeps single-generation inference.}
\end{displayquote}

\newpage
\begin{commentbox}{R4.2 Clinical Safety Implications}
The discussion of clinical safety implications is insufficient. The evaluation does not
show whether improved ARU-level faithfulness translates into safer recommendations or
reduced harm.
\end{commentbox}

\responseheading{Response and Revision Summary}
We thank the reviewer for pointing out that the original manuscript motivated Fa-DPO
through clinical safety but did not sufficiently distinguish automatic process-metric
improvement from direct evidence of real-world patient safety. We agree that reduced
ARU-level violations should not be equated with proven clinical safety or reduced
real-world clinical harm.

In the revision, we frame ARU-level metrics as safety-relevant process proxies and add
a retrospective clinician harm-severity audit. The audit provides preliminary support
that the process signal is safety-relevant, but we present it as retrospective
one-backbone evidence rather than prospective clinical safety validation or deployment
safety proof.


\responseheading{Revision Details}
\textbf{Clinical-safety discussion: add retrospective harm-severity evidence and limits.}
\begin{displayquote}
\deleted{Their reasoning often contradicts medical evidence, leading to unreliable diagnoses that compromise clinical safety.}\added{Their reasoning often contradicts medical evidence, creating safety-relevant risks for clinical decision support.}

\added{Clinical safety is a central motivation for evidence-grounded MQA, but reduced ARU-level violations should be interpreted as a safety-relevant process signal rather than direct evidence of reduced real-world harm or deployment safety.}

\added{The blinded retrospective clinician harm-severity audit uses 200 matched MedCaseReasoning question IDs from the representative II-Medical-8B pair, yielding 400 anonymized responses. In this audit, an unsupported high-stakes claim is a diagnosis-, triage-, treatment-, monitoring-, contraindication-, or reassurance-relevant assertion that is not supported by the case, retrieved evidence, or preceding reasoning and could plausibly affect clinical management if relied on without clinician oversight.}

\captionof{table}{\added{Retrospective clinician harm-severity audit on paired Standard DPO-Retained-34K and Fa-DPO-Retained-34K outputs for the representative II-Medical-8B pair. Deltas are Fa-DPO-Retained-34K minus Standard DPO-Retained-34K.}}
\begin{adjustbox}{max width=\linewidth}
\begin{tabular}{lcccc}
\toprule
\textbf{Outcome} & \textbf{Std. DPO-Retained-34K} & \textbf{Fa-DPO-Retained-34K} & \textbf{$\Delta$} & \textbf{95\% CI / test} \\
\midrule
Mean harm severity & 0.74 & 0.61 & -0.13 & [-0.24, -0.02] \\
Moderate/high harm (\%) & 22.5 & 17.0 & -5.5 & [-10.8, -0.8]; McNemar $p=0.042$ \\
Unsafe recommendation (\%) & 13.0 & 9.5 & -3.5 & [-7.8, +0.3]; McNemar $p=0.085$ \\
Unsupported high-stakes claim (\%) & 17.5 & 12.0 & -5.5 & [-10.1, -1.2]; McNemar $p=0.018$ \\
Contradictory recommendation (\%) & 10.5 & 7.0 & -3.5 & [-7.4, +0.4]; McNemar $p=0.074$ \\
Clinician oversight warning (\%) & 20.0 & 15.5 & -4.5 & [-9.2, +0.2]; McNemar $p=0.061$ \\
\bottomrule
\end{tabular}
\end{adjustbox}

\added{Table~\ref*{ms-tab:harm_severity_audit} shows modest safety-relevant shifts. Fa-DPO reduces mean harm severity by 0.13, unsupported high-stakes claims by 5.5 points, and moderate/high harm by 5.5 points; unsafe recommendations, contradictory recommendations, and oversight warnings move in the same direction but have intervals that include zero. In non-tie pairwise lower-harm judgments, clinicians prefer Fa-DPO in 62.2\% of cases, but 104 of 200 pairs are ties. The verifier-derived risk score also correlates with clinician-rated harm severity (Spearman $\rho=0.42$, $p<0.001$). These findings support the safety relevance of the process signal, but remain a retrospective audit of one backbone and subset rather than prospective clinical safety validation.}
\end{displayquote}

\newpage
\begin{commentbox}{R4.3 Dataset Bias}
The preference dataset derived from four controlled failure types is valuable, but the
data construction process needs further clarification regarding dataset bias.
\end{commentbox}

\responseheading{Response and Revision Summary}
We agree that the 34,204 preference pairs should not be read as an unbiased sample of
natural clinical reasoning errors. The dataset is a controlled counterfactual training
set: useful for known error types and fixed preferred/rejected contrasts, but not for
estimating natural error prevalence.

We added a retained-pool data card, source-reuse diagnostics, type-wise retention
rates, and a local-rewrite naturalness audit. The revised manuscript therefore presents
the retained pool as controlled, risk-enriched supervision, not as a complete or
prevalence-preserving natural-error sample.



\responseheading{Revision Details}
\textbf{Dataset-bias paragraph: define the retained pool as controlled and risk-enriched.}
\begin{displayquote}
\deleted{We use one constructed preference dataset for alignment training, four main-benchmark settings for final-answer accuracy, one paired reasoning-analysis dataset for controlled process-faithfulness evaluation, and one supplementary external pilot set for natural-error coverage.}

\added{The retained preference pool is controlled and risk-enriched rather than prevalence-preserving. All four F1--F4 failure types remain represented, but Appendix~\ref*{ms-app:dataset_bias_details} reports the retained-pool data card, retained/raw rates, local-rewrite naturalness audit, and bias-source table. We therefore use the retained pool as controlled supervision for evidence-use and reasoning-answer consistency failures, not as an unbiased natural-error sample.}

\added{In the retained-pool audit, most sampled constructed negatives remain coherent, clinically plausible, and model-like after risk-based retention, with a strict accept-for-training rate of 65.0\% (95\% Wilson CI: 60.2--69.5\%).}
\end{displayquote}

\textbf{Appendix bias-source table.}
\begin{displayquote}
\captionof{table}{\added{Main dataset-bias sources and how they are handled in the revised manuscript.}}
\begin{adjustbox}{max width=\linewidth}
\begin{tabular}{L{0.18\linewidth}L{0.39\linewidth}L{0.34\linewidth}}
\toprule
\textbf{Bias source} & \textbf{Observed boundary} & \textbf{Manuscript handling} \\
\midrule
Source dataset & All preference pairs come from MedCaseReasoning training positives, so source specialty and case-style coverage may not match broader clinical QA; \added{repeated perturbations of the same source cases may also reduce source diversity}. & We report the source size \added{and source-reuse distribution} and avoid population-prevalence claims. \\
Taxonomy & F1--F4 cover evidence-use and reasoning-answer consistency failures in evidence-grounded MQA, not all medical hallucination phenomena. & We restrict external evaluation to task-aligned MQA benchmarks and avoid treating hallucination-oriented benchmarks as evidence-use taxonomy tests. \\
Construction & Local rewrites provide target-type control, and the retained-pool audit shows 65.0\% strict accept-for-training, but constructed responses remain noisy controlled counterfactuals rather than natural model-sampled errors. & We report clinician/audit results, interpret constructed negatives as controlled counterfactuals, and avoid natural-prevalence claims. \\
\bottomrule
\end{tabular}
\end{adjustbox}
\end{displayquote}

\textbf{Appendix retained-pool and source-reuse diagnostics.}
\begin{displayquote}
\captionof{table}{\added{Type-wise retained-pool composition for the constructed preference dataset.}}
\begin{adjustbox}{max width=\linewidth}
\begin{tabular}{lrrrrr}
\toprule
\textbf{Type} & \textbf{Design-implied raw attempts} & \textbf{Retained pairs} & \textbf{Retained/raw} & \textbf{Share of retained pool} & \textbf{Mean ARUs} \\
\midrule
F1 & 13,092 & 6,715 & 51.3\% & 19.6\% & 20.46 \\
F2 & 13,092 & 9,607 & 73.4\% & 28.1\% & 20.98 \\
F3 & 13,092 & 9,456 & 72.2\% & 27.6\% & 20.70 \\
F4 & 13,092 & 8,426 & 64.4\% & 24.6\% & 20.87 \\
\midrule
Overall & 52,368 & 34,204 & 65.3\% & 100.0\% & 20.77 \\
\bottomrule
\end{tabular}
\end{adjustbox}

\added{Overall, the pipeline generated 52,368 raw rewrite attempts, 47,393 passed automatic QC, and 34,204 were retained after ARU-risk filtering. The overall QC pass rate was 90.5\%, the retained/QC rate was 72.2\%, and the retained/raw rate was 65.3\%.}

\added{The source-reuse distribution shows that the retained pool repeatedly perturbs many source positives: 62.4\% of source positives contribute three or four retained pairs, accounting for 77.3\% of retained pairs. This improves within-case counterfactual coverage but limits source diversity and may increase overfitting to source cases or case-specific reasoning patterns. The retained-pool composition table further shows that retention is not prevalence-preserving across controlled failure types: F2 and F3 have the highest retained/raw rates and retained-pool shares, F1 has the lowest retained/raw rate and retained-pool share, and F4 remains intermediate. Because type-wise QC/discarded counts are not part of the recorded type-wise construction outputs, we report QC flow only at the overall level and do not attribute type-wise retained/raw differences to a specific QC stage.}
\end{displayquote}

\textbf{Appendix retained-pool local-rewrite naturalness table.}
\begin{displayquote}
\captionof{table}{\added{Clinician validation of the risk-retained controlled negative pool. All rates are computed on 400 retained negatives, stratified as 100 items per F1--F4 type.}}
\begin{adjustbox}{max width=\linewidth}
\begin{tabular}{lcccccccccc}
\toprule
\textbf{Type} & \textbf{$n$} & \textbf{Type Match} & \textbf{Single Error} & \textbf{Policy} & \textbf{Coherence} & \textbf{Model-like} & \textbf{Plaus.} & \textbf{Faith. Degr.} & \textbf{Strict Accept} & \textbf{Strict 95\% CI} \\
\midrule
F1 & 100 & 74.0 & 68.0 & 98.0 & 84.0 & 76.0 & 82.0 & 72.0 & 56.0 & [46.2, 65.3] \\
F2 & 100 & 84.0 & 78.0 & 99.0 & 91.0 & 86.0 & 89.0 & 80.0 & 70.0 & [60.4, 78.1] \\
F3 & 100 & 86.0 & 80.0 & 98.0 & 90.0 & 85.0 & 88.0 & 84.0 & 72.0 & [62.5, 79.9] \\
F4 & 100 & 76.0 & 70.0 & 89.0 & 87.0 & 81.0 & 85.0 & 76.0 & 62.0 & [52.2, 70.9] \\
\midrule
Overall & 400 & 80.0 & 74.0 & 96.0 & 88.0 & 82.0 & 86.0 & 78.0 & 65.0 & [60.2, 69.5] \\
\bottomrule
\end{tabular}
\end{adjustbox}
\end{displayquote}

\textbf{Discussion and limitations.}
\begin{displayquote}
\added{The source-reuse distribution shows that the retained pool repeatedly perturbs many source positives: 62.4\% of source positives contribute three or four retained pairs, accounting for 77.3\% of retained pairs. This improves within-case counterfactual coverage but limits source diversity and may increase overfitting to source cases or case-specific reasoning patterns.}

\added{The retained-pool composition table further shows that retention is not prevalence-preserving across controlled failure types: F2 and F3 have the highest retained/raw rates and retained-pool shares, F1 has the lowest retained/raw rate and retained-pool share, and F4 remains intermediate.}

\added{The retained-pool clinician audit in Table~\ref*{ms-tab:negative_quality} shows that the constructed negatives are usually acceptable as controlled counterfactuals after retention. We therefore use the retained pool as controlled, risk-enriched supervision and avoid treating it as a natural-prevalence sample.}
\end{displayquote}

\bibliographystyle{cas-model2-names}
\bibliography{sample-base}

\end{document}
